\documentclass[a4paper,11pt]{article}

\usepackage[top=2cm,bottom=2cm,left=2cm,right=2cm]{geometry}
\usepackage[pdftex]{graphicx} %per poter inserire le figure
\usepackage{amssymb,amsmath,amsthm,amsfonts}
\usepackage{xspace}
\usepackage{tabularx}
\usepackage{indentfirst}
\usepackage{subfigure}
\usepackage[small]{caption}
\usepackage{eucal}
\usepackage{eso-pic}
\usepackage{url}
\usepackage{booktabs}
\usepackage{afterpage}
\usepackage{parskip}
\usepackage{listings}
\usepackage{fancyhdr}
\usepackage{textcomp}
\usepackage{cite}
\usepackage{multirow}
\usepackage{tikz}
\usepackage{tikz-cd}
\usepackage{cancel}
\usepackage{epsfig}
\usepackage{bbold}
\newcommand{\at}[2][]{#1|_{#2}}
\usepackage[utf8]{inputenc}   %per riuscire a scrivere gli accenti
\usepackage[italian]{babel}   %per riuscire a scrivere gli accenti

\pagestyle{fancy} 



\begin{document}

\newtheorem{Prop}{Proposizione}[section]
\newtheorem{Ass}{Assioma}
\newtheorem{Def}{Definizione}[section]
\newtheorem{Teo}{Teorema}[section]
\newtheorem{Lem}{Lemma}[section]
\newtheorem{Cor}{Corollario}[section]


\begin{titlepage}
\vspace{5mm}
\begin{figure}[hbtp]
\centering
\includegraphics[scale=0.13]{logo1.png}
\end{figure}
\vspace{5mm}
\begin{center}
{{\huge{\textsc{\bf UNIVERSIT\`A DEGLI STUDI DI PADOVA}}}\\}
\vspace{5mm}
{\Large{\bf Dipartimento di Fisica e Astronomia ``Galileo Galilei''}} \\
\vspace{5mm}
{\Large{\textsc{\bf Corso di Laurea Triennale in Fisica}}}\\
\vspace{20mm}
{\Large{\textsc{\bf Tesi di Laurea}}}\\
\vspace{30mm}

{\LARGE \textbf{Strutture di Poisson e sistemi Hamiltoniani sullo spazio dei loop formale}}\\

\vspace{8mm}
\end{center}

\vspace{20mm}

\begin{tabular}{ l  c  c c c  cc c c c c  l }
{\Large{\bf Relatore}} &&&&&&&&&&& {\Large{\bf Laureando}}\\
{\Large{\bf Prof./Dr. Paolo Rossi}} &&&&&&&&&&& {\Large{\bf Filippo Pesce}}\\
\end{tabular}

\vspace{15 mm}

\begin{center}
{\Large{\bf Anno Accademico 2021/2022}}
\end{center}
\end{titlepage}
\clearpage{\pagestyle{empty}\cleardoublepage}


\newpage


\tableofcontents


\newpage







\section{Introduzione}
In questa prima sezione diamo una prima introduzione alla geometria simplettica, la quale fornisce un linguaggio naturale per la descrizione Hamiltoniana della meccanica. Dai recenti sviluppi, questa branca di studio ha conquistato un enorme interesse e risulta tutt'ora un settore centrale della geometria differenziale.

$\textbf{IMPORTANTE}$:In tutto il testo adotteremo indici alti per esprimere le coordinate sulla varietà generica, e sfrutteremo la convenzione di Einstein sugli indici ripetuti quando non esplicitato.

\subsection{Spazi vettoriali simplettici}
Esponiamo alcuni risultati di algebra lineare preliminari, essenziali per i risultati che troveremo più avanti.

\begin{Def}
Sia V uno spazio vettoriale sul campo $\mathbb{R}$ di dimensione $dimV=m$, e sia $\Omega:V\times V\rightarrow\mathbb{R}$ una mappa bilineare. Quest'ultima si dice essere "antisimmetrica" se $\Omega(u,v)=-\Omega(v,u)$ $\forall u,v\in V$.
\end{Def}
\begin{Teo}
Se $\Omega$ è una forma bilineare santisimmetrica, allora esiste una base $\left\{u_{1},...,u_{k},e_{1},...,e_{n},f_{1},...,f_{n}\right\}$ di V che soddisfa
\[\Omega(u_{i},v)=0\hspace{0.4cm}\forall i=1,...,k\hspace{0.4cm}e\hspace{0.4cm}\forall  v  \in V;\]
\[\Omega(e_{i},e_{j})=\Omega(f_{i},f_{j})=0\hspace{0.4cm}\forall i,j=1,...,n;\]
\[\Omega(e_{i},f_{j})=\delta_{ij}\hspace{0.4cm}\forall i,j=1,...,n. \]
\end{Teo}
\begin{proof}[Dimostrazione]
Si consideri il sottospazio vettoriale $U:=\left\{u\in V: \Omega(u,v)=0\hspace{0.2cm} \forall v\in V \right\}$, e scegliamone una base $\left\{u_{1},...,u_{k}\right\}$.
Possiamo sempre scrivere 
\[V=U\oplus\ W\]
con W lo spazio complementare di U. Si prenda ora un elemento NON nullo $0\ne e_{1}\in W$, allora esisterà un elemento $f_{1}\in W$ tale per cui si abbia $\Omega(e_{1},f_{1})\ne 0$ (in caso contrario $e_{1}$ sarebbe per definizione un elemento di U;inoltre $f_{1}\in W$ perchè se appartenesse ad U avrei sempre $\Omega(e_{1},f_{1})=-\Omega(f_{1},e_{1})=0$ per definizione di U) che a costo di un riscalamento possiamo assumere essere $\Omega(e_{1},f_{1})=1$.Si definiscano ora i seuguenti sottospazi di W
\[ W_{1}:=\mathbb{R}\bigl\langle e_{1},f_{1}\bigr\rangle\]
\[W_{1}^{\Omega}:=\{w\in W : \Omega(w,v)=0\hspace{0.2cm}\forall v\in W_{1}\}\]
dove ricordiamo che $\mathbb{R}\bigl\langle e_{1},f_{1}\bigr\rangle$ coinciude con l'insieme delle combinazioni lineari formali di $\{e_{1},f_{1}\}$.Utilizziamo ora il seguente 
\begin{Lem}
Vale $W_{1}\cap W_{1}^{\Omega}=\{0\}$
\end{Lem}
\begin{proof}[Dimostrazione]
Infatti si noti che se $W_{1}\cap W_{1}^{\Omega}\ni v:=a e_{1}+b f_{1}$, allora
\begin{eqnarray}
0=\Omega(v,e_{1})=a\Omega(e_{1},e_{1})+b\Omega(f_{1},e_{1})=0-b\Omega(e_{1},f_{1})=-b \\
0=\Omega(v,f_{1})=a\Omega(e_{1},f_{1})+b\Omega(f_{1},f_{1})=a
\end{eqnarray}
essendo che $e_{1},f_{1}\in W_{1}$, ma allora $a=b=0\Rightarrow v=0$.
\end{proof}
Inoltre vale il seguente 
\begin{Lem}
Vale $W=W_{1}\oplus W_{1}^{\Omega}$
\end{Lem}
\begin{proof}[Dimostrazione]
Supponiamo che $v\in W$ sia tale che $\Omega(v,e_{1})=c$ e $\Omega(v,f_{1})=d$, allora posso sempre scrivere
\[v=(v\textcolor{blue}{+cf_{1}}\textcolor{red}{-de_{1}})\textcolor{blue}{-cf_{1}}\textcolor{red}{+de_{1}}\]
dove il primo termine tra parentesi sta in $W_{1}^{\Omega}$ ed il secondo in $W_{1}$. Infatti
\[\Omega(v+cf_{1}-de_{1},e_{1})=\Omega(v,e_{1})+c\Omega(f_{1},e_{1})+0=c-c=0\]
e lo stesso per $\Omega(v+cf_{1}-de_{1},f_{1})$.Infine per il secondo termine basta osservare che ogni elemento di $W_{1}$ è per definizione una combinazione lineare dei due $e_{1},f_{1}$.
\end{proof}
Procedimo con la dimostrazione: abbiamo visto sin'ora che $V=U\oplus W$ dove W si può scrivere a sua volta come somma diretta di due sottospazi $W=W_{1}\oplus{W_{1}^{\Omega}}$. Si prenda ora $W_{1}^{\Omega}\ni e_{2}\ne0$, dal momento che $\forall u\in U$ vale $\Omega(e_{2},u)=-\Omega(u,e_{2})=0$ e $\Omega(e_{2},w)=0\hspace{0.2cm} \forall w\in W_{1}$, esisterà un $0\ne f_{2}\in W_{1}^{\Omega}$ tale per cui $\Omega(e_{2},f_{2})\ne0$ (che assumeremo essere pari a 1 a costo di ruscalare il tutto) poichè altrimenti $e_{2}$ sarebbe un elemento di $U\oplus W_{1}$.Sia poi 
\[W_{2}:=\mathbb{R}\bigl\langle e_{2},f_{2}\bigr\rangle\]
\[W_{2}^{\Omega}:=...\]
e si itera il ragionamento.Dal momento che $dimV<+ \infty$, il processo terminerà con
\[V=U\oplus W_{1}\oplus W_{2}\oplus ... \oplus{W_{n}}\]
dove $\Omega(U,W_{i})=0\hspace{0.2cm}$,$\Omega(W_{i},W_{j})=0$ $\forall i,j=1,...,n$ e $\Omega(e_{i},f_{j})=\delta_{ij}$.
\end{proof}
$\textbf{NOTA}$:
\begin{itemize}
    \item La base trovata non è unica;
    \item In notazione matriciale si ha
    \[\Omega(u,v)=\left(\begin{array}{ccccc}
    u^{1}& . & . & . & u^{m} \\
    \end{array}\right) \left(\begin{array}{ccc}
    \mathbb{0}_{k} & \mathbb{0}_{k\times n}  & \mathbb{0}_{k\times n}\\
    \mathbb{0}_{n\times k} & \mathbb{0}_{n}  & \mathbb{1}_{n}\\
    \mathbb{0}_{n\times k} & \mathbb{-1}_{n} & \mathbb{0}_{n}\\
    \end{array}\right) \left(\begin{array}{c}
    v^{1}       \\
    .           \\
    .           \\
    .           \\
    v^{m}       \\
    \end{array}\right)\];
    \item Il sottospazio U non dipende invece dalla base scelta, perciò è un invariante di $(V,\Omega)$;
    \item Siccome $m=k+2n$, n è ancora un invariante numerico di $(V,\Omega)$.
    \end{itemize}

\begin{flushleft}    
$\underline{Esempio}$: Sia V uno spazio vettoriale 2n-dimensionale, e sia $\{A_{1},B_{1},...,A_{n},B_{n}\}$ la corrispondente base mentre si denoti con $\{\alpha_{1},\beta_{1},...,\alpha_{n},\beta_{n}\}$ la rispettiva base duale di $V^{*}$.Definito il tensore
\[\omega:=\sum_{i}\alpha_{i}\wedge\beta_{i}\in\Lambda^{2}(V)\]
si ha che 
\[\omega(A_{i},B_{i})=-\omega(B_{i},A_{i})=\delta_{ij}\]
\[\omega(A_{i},A_{j})=\omega(B_{i},B_{j})=0\]
Supponiamo poi che $X=\sum_{i}(a_{i}A_{i}+b_{i}B_{i})\in V$ soddisfi $\omega(X,Y)=0$ $\forall Y\in V$. Allora $0=\omega(X,B_{i})=a_{i}$ e $\omega(X,A_{i})=-b_{i}$, il che implica che $X=0$ e che perciò $\omega$ sia non degenere.
\end{flushleft}

\begin{Def}
Sia V uno spazio vettoriale sul campo $\mathbb{R}$ di dimensione $dimV=m$, e sia $\Omega:V\times V\rightarrow\mathbb{R}$ una mappa bilineare. Definiamo la mappa da essa indotta come 
\[ \widetilde{\Omega}:V\rightarrow V^{*}   \]
\[ \widetilde{\Omega}(v)(w):=\Omega(v,w) \hspace{0.5cm}\forall v,w\in V\]
\end{Def}

Si noti che in questo modo $\ker  {\widetilde{\Omega}}=\{{u\in V: \Omega(u,v)=0\hspace{0.2cm} \forall v\in V} \} =U$ dove U è lo spazio precedentemente introdotto.

\begin{Def}
Una mappa bilineare antisimmetrica $\Omega$ si dice "simplettica" se la mappa indotta $\widetilde{\Omega}$ è invertibile, ossia $\ker\widetilde{\Omega}=U={0}$. In tal caso la coppia $(V,\Omega)$ si dice essere uno "spazio vettoriale simplettico".
\end{Def}

$\textbf{NOTA}$:In tal caso \begin{itemize}
    \item La biiezione $\widetilde{\Omega}:V\rightarrow V^{*}$ permette di identificare V e $V^{*}$;
    \item Poichè $dimU=k=0\Rightarrow m=2n$ per il Teorema sopra enunciato;
    \item Vi è una "base canonica" $\{e_{1},...,e_{n},f_{1},...,f_{n}\}$ di V tale per cui
    \[\Omega(e_{i},e_{j})=\Omega(f_{i},f_{j})=0\hspace{0.4cm} \forall i,j=1,...,n\]
    \[\Omega(e_{i},f_{j})=\delta_{ij}\hspace{0.4cm} \forall i,j=1,...n\]
    \[\Omega(v,w)=\left(\begin{array}{ccccc}
    v^{1}& . & . & . & v^{2n} \\
    \end{array}\right) \left(\begin{array}{cc}
    \mathbb{0}_{n} & \mathbb{1}_{n}  \\
    \mathbb{-1}_{n} & \mathbb{0}_{n} \\
    \end{array}\right) \left(\begin{array}{c}
    w^{1}       \\
    .           \\
    .           \\
    .           \\
    w^{2n}       \\
    \end{array}\right)\]
\end{itemize}

\begin{Def}
Un sottospazio vettoriale $W\subseteq V$ di uno spazio vettoriale simplettico $(V,\Omega)$ si dice essere "simplettico" se $\Omega |_{W}$ è simplettica. Invece si dice "isotropo" se $\Omega |_{W}(u,v)=0\hspace{0.2cm}\forall u,v\in W$.
\end{Def}

\begin{flushleft}
$\underline{Esempio}$: Riprendendo la dimostrazione del Teorema precedente, abbiamo che $W_{1}=\mathbb{R}\bigl\langle e_{1},f_{1}\bigl\rangle $ è simplettico (infatti $\ker\widetilde{\Omega}|_{W_{1}}=\{ w\in W_{1} : \Omega(w,v)=0\hspace{0.2cm}\forall v\in W_{1}\}={0}$ come è ovvio per costruzione dei vettori $e_{1},f_{1}$, cioè $\Omega (e_{1},f_{1})=1$), mentre lo spazio $W:=\mathbb{R}\bigl\langle e_{1},e_{2}\bigl\rangle$ è invece isotropo (sempre per costruzione dei vettori $e_{1},e_{2}$).
\end{flushleft}

\begin{Def}
Un "simplettomorfismo" di due spazi vettoriali simplettici $(V,\Omega)$ e $(V',\Omega')$ è un isomorfismo lineare $\varphi :V\rightarrow V$ tale per cui $\varphi^{*}\Omega'=\Omega$, ricordando che per definizione $\varphi^{*}\Omega'(u,v)=\Omega'(\varphi(u),\varphi(v))$. In tal caso $(V,\Omega)$ e $(V',\Omega')$ si dicono essere "simplettomorfi".
\end{Def}

\begin{flushleft}
$\textbf{NOTA}$:Il prototipo di spazio vettoriale simplettico è $(\mathbb{R}^{2n},\Omega_{0})$, con $\Omega_{0}$ tale per cui la base 
\[e_{1}=(1,0,...,0) \hspace{0.3cm},..., \hspace{0.3cm} e_{n}=(0,...,0,1,0,...,0)\]
\[f_{1}=(0,...,0,1,0,...,0)\hspace{0.3cm},...,\hspace{0.3cm}f_{n}=(0,...,0,1)\]
sia canonica.
Per il Teorema appena visto, ogni spazio vettoriale simplettico $2n-$dimensionale $(V,\Omega)$ è simplettomorfo ad $(\mathbb{R}^{2n},\Omega_{0})$.
\end{flushleft}


\subsection{Varietà simplettiche}
Sia ora M una varietà differenziabile m-dimensionale,e sia $\omega\in \mathcal{A}^{2}(M)$ una $2$-forma definita su di essa.
\begin{Def}
La forma differenziale $\omega$ si dice "simplettica" se è chiusa (d$\omega=0$) e $\omega|_{p}$ è simplettica per ogni punto $p\in M$ della varietà.
\end{Def}
Si noti che se $\omega$ è simplettica, allora dimM deve essere pari.
\begin{Def}
Una "varietà simplettica" è una coppia $(M,\omega)$ dove M è una varietà liscia e $\omega$ una forma simplettica su M.
\end{Def}
$\underline{Esempio}$:\begin{itemize}
    \item Sia $M=\mathbb{R}^{2n}$ con coordinate lineari $\{x^{1},...,x^{n},y^{1},...,y^{n}\}$.La forma 
    \begin{equation}
    \omega_{0}=\sum_{i=1}^{n}dx^{i}\wedge dy^{i}
    \end{equation} è simplettica (d$\omega_{0}=\sum_{i}[d(dx^{i})\wedge dy^{i}+(-1)^{1}dx^{i}\wedge d(dy^{i})]=0$), e la base $\{(\frac{\partial}{\partial x^{1}}|_{p}),...,(\frac{\partial}{\partial x^{n}}|_{p})\}$ è una base canonica per $T_{p}M$. In particolare si osservi che corrisponde, in ogni punto, al tensore deifinito nell'esempio di pagina 3 e perciò risulta non degenere.
    \item Dal momento che quanto detto sin'ora si può estendere al caso di varietà complesse, si consideri $M=\mathbb{C}^{n}$ con coordinate lineari $\{z^{1},...,z^{n}\}$, la forma 
    \[\omega_{0}=\frac{i}{2}\sum_{k=1}^{n}dz^{k}\wedge d\bar{z}^{k}\]
    è simplettica, poichè $z^{k}=x^{k}+iy^{k}$ e dunque essa eguaglia la 2-forma precedente sotto l'isomorfismo $\mathbb{C}^{n}\simeq \mathbb{R}^{2n}$.Infatti 
    \[\omega_{0}=\frac{i}{2}\sum_{k=1}^{n}d(x^{k}+iy^{k})\wedge d(x^{k}-iy^{k})=\frac{1}{2}\sum_{k=1}^{n}[-i^{2}dx^{k}\wedge dy^{k}+i^{2}dy^{k}\wedge dx^{k}]=\frac{2}{2}\sum_{k=1}^{n}dx^{k}\wedge dy^{k}\]
    \item Sia $M=S^{2}$ vista come il sottoinsieme di vettori di $\mathbb{R}^{3}$ con norma unitaria.Si ha
    \[T_{p}S^{2}=\{v\in\mathbb{R}^{3} : v\cdot p=0\}=(\mathbb{R}\big\langle p \big\rangle)^{\bot}\]
    Se allora considero $\omega_{p}(u,v):=p\cdot (u\times v)$ con $u,v\in T_{p}S^{2}$, essa è chiusa (perchè $dimM=2$,vedasi Teorema 1.1) ed è non degenere poichè $p\cdot (u\times v)= 0$ se e solamente se $v$ oppure $u$ sono nulli (dunque $(u\times v)= 0\forall v\in\mathbb{R}^{3}\iff u=0$).
    \begin{figure}
        \centerline{\includegraphics[scale=.03]{TS2.jpg}}
        \caption{Raffigurazione grafica di $T_{p}S^{2}$}
        \label{fig:my_label}
    \end{figure}
\end{itemize}
\begin{Prop}
Data una varietà simplettica $(M,\omega)$ di dimensione $dimM=2n$,la forma differenziale $\omega^{n}:=\omega\wedge...\wedge\omega \in \Lambda^{2n}(M)$ è una forma di volume su M,ossia non si annulla in ogni punto di M.
\end{Prop}
\begin{proof}[Dimostrazione]
Per definizione $\omega|_{p}=\omega_{p}$ è simplettica $\forall p\in M$ dunque per il Teorema 1.1 esiste una base $\{e_{1},...,e_{n},f_{1},...,f_{n}\}$ di $T_{p}M$ tale per cui 
\[\omega^{n}_{p}(e_{1},f_{1},...,e_{n},f_{n})=\omega_{p}(e_{1},f_{1})\cdot ...\cdot\omega_{p}(e_{n},f_{n})=1^{n}=1\]
\end{proof}
\begin{Cor}
Ogni varietà simplettica è orientabile
\end{Cor}
\begin{Cor}
$S^{2n}$ NON ammette una struttura simplettica per qualsiasi $\mathbb{N}\ni n>1$.
\end{Cor}
\begin{proof}[Dimostrazione]
Si supponga che $\omega$ sia simplettica, ed assumiamo che $\omega=d\alpha$ per qualche $\alpha\in\Lambda (S^{2n})$, allora essendo che
\[d(\alpha\wedge\omega^{n-1})=d\alpha\wedge\omega^{n-1}+(-1)^{1}\alpha\wedge d(\omega^{n-1})=\omega^{n}-\alpha\wedge[d\omega\wedge\omega^{n-2}+(-1)^{2}\omega\wedge d\omega^{n-2}]=\omega^{n}-[0+...+0]=\omega^{n}\]
allora sfruttando Stokes ho
\[0\neq\int_{S^{2n}}\omega^{n}=\int_{ S^{2n}}d(\alpha\wedge\omega^{n-1})=\int_{\partial S^{2n}=\emptyset}\alpha\wedge\omega^{n-1}=0\]
che è una contraddizione!Ma allora la classe di coomologia $[\omega]\in H^{2}(S^{2n},\mathbb{R})$ dovrebbe essere non nulla, mentre si sa che $H_{dR}^{2}(S^{2n},\mathbb{R})={0}$ per n>1.
\end{proof}



\subsection{Struttura del fibrato cotangente}
Sia X una varietà liscia n-dimensionale ed $M=T^{*}X$ il suo fibrato cotangente. Se la varietà X è descritta localmente dalla carta $(U,\varphi)$ dove $\varphi:U\rightarrow\mathbb{R}^{n}$ è definita in coordinate locali come $\varphi(x)=(x^{1},...,x^{n})$, allora ad ogni punto $x\in X$ sappiamo che $T_{x}^{*}X$ è descritto dalla base $\{dx^{1}\at{x},...,dx^{n}\at{x}\}$.Se poi $\zeta\in T_{x}^{*}X$, posso sempre scrivere $\zeta=\sum_{i}\zeta_{i}dx^{i}\at{x}$ per opportuni coefficenti $\zeta_{i}$. Questo permette di definire una mappa
\[\Phi:T^{*}U\rightarrow \mathbb{R}^{2n}\]
\[(x,\zeta)\mapsto(x^{1},...,x^{n},\zeta^{1},...,\zeta^{n})\]
La coppia $(\pi^{-1}(U),\Phi)$ è così una carta di $T^{*}X$, dove $\pi:T^{*}X\rightarrow X$ è la proiezione naturale.
\begin{Def}
Sia X una varietà liscia, $(U,\varphi)$ una sua carta locale e $(\pi^{-1}(U),\Phi)$ la carta associata al fibrato cotangente. Definiamo la 2-forma $\mathcal{A}^{2}(T^{*}U)$ come
\begin{equation}
\omega=\sum_{i=1}^{n}dx^{i}\wedge d\zeta^{i}
\end{equation}
detta anche "forma canonica".
Definiamo invece la "1-forma di Liouville" la forma differeneziale
\begin{equation}
\mathcal{A}(T^{*}U)\ni\alpha=\sum_{i=1}^{n}\zeta^{i}dx^{i}.
\end{equation}
\end{Def}

Si noti che $\omega=-d\alpha$, e che 

\begin{Prop}
La forma di Liouville $\alpha\in\mathcal{A}(T^{*}U)$ è indipendente dalla scelta delle coordinate.
\end{Prop}

\begin{proof}[Dimostrazione]
Siano $(\pi^{-1}(U),x^{1},...,x^{n},\zeta^{1},...,\zeta^{n})$ e $(\pi^{-1}(\widetilde{U}),\widetilde{x}^{1},...,\widetilde{x}^{n},\widetilde{\zeta}^{1},...,\widetilde{\zeta}^{n})$ due carte del fibrato cotangente, in $\pi^{-1}(U\cap\widetilde{U})$ le coordinate sono collegate dalla relazione $\widetilde{\zeta}^{i}=\sum_{k}\zeta^{k}\frac{\partial x^{k}}{\partial \widetilde{x}^{i}}$. Dal momento che $d\widetilde{x}_{i}=\sum_{j}\frac{\partial \widetilde{x}^{i}}{\partial x^{j}}dx^{j}$ si trova
\[\alpha=\sum_{p}\widetilde{\zeta}^{p}d\widetilde{x}^{p}=\sum_{p,k,m}\zeta^{k}\frac{\partial x^{k}}{\partial \widetilde{x}^{p}}\frac{\partial\widetilde{x^{p}}}{\partial x^{m}}dx^{m}=\sum_{k}\zeta^{k}dx^{k}\]
\end{proof}

Si consideri nuovamemnte il fibrato $M=T^{*}X$ con la sua proiezione naturale $\pi:M\rightarrow X$
\[\begin{tikzcd}
T_{p}M \arrow[r,"\pi_{*}"]  & T_{x}X=T_{\pi(p)}X \\
T_{x}^{*}X=T_{\pi(p)}^{*}X \arrow[r,"\pi^{*}"] & T_{p}^{*}M 
\end{tikzcd}\]
Un punto di $T^{*}X$ sarà una coppia $(x,\zeta)=:p$ dove $\zeta\in T^{*}_{x}X$ per qualche $x\in X$.

La 1-forma di Liouville si può alternativamente definire come quella $\alpha\in\mathcal{A}(T^{*}X)$ tale per cui

\begin{equation}
\alpha_{p}=\alpha_{(x,\zeta)}:=\pi^{*}\zeta\in T_{p}^{*}M
\end{equation}

per qualche $\zeta\in T_{x}^{*}X$.Infatti abbiamo $\pi(x^{1},...,x^{n},\zeta^{1},...,\zeta^{n})=(x^{1},...,x^{n})$ dunque $\frac{\partial\pi^{\beta}(p)}{\partial\zeta^{\alpha}}=\frac{\partial x^{\beta}}{\partial \zeta^{\alpha}}=0\hspace{0.3cm} \forall \alpha,\beta=1,...,n$
\[\Rightarrow\pi^{*}(\sum_{p}\zeta^{p}dx^{p})=\sum_{p}\zeta^{p}d(x^{p}\circ\pi)=\sum_{p,k}\zeta^{p}[\frac{\partial x^{p}}{\partial x_{k}}dx^{k}+\cancel{\frac{\partial x^{p}}{\partial\zeta^{k}}}d\zeta^{k}]=\sum_{p}\zeta^{p}dx^{p}\]

In altre parole, se $X\in T_{p}(T^{*}X)$ è un vettore tangente, allora per definizione
\[\alpha_{p}(X)=\alpha_{(x,\zeta)}(X)=\pi^{*}\zeta(X)=\zeta(\pi_{*}X)\]

Mostriamo ora l'equivalenza delle due definizioni appena date per $\alpha$:

\begin{Prop}
Sia X una varietà liscia, la forma di Liouville $\alpha$ è liscia e $\omega=-d\alpha$ è una forma simplettica su $T^{*}X$.
\end{Prop}
\begin{proof}[Dimostrazione]
Siano $(x)^{i}$ le coordinate locali su X, e siano $(x^{i},\zeta^{i})$ le coordinate locali su $T^{*}X$, ricordiamo che le coordinate del generico $(x,\zeta)\in T^{*}X$ sono proprio le $(x^{i},\zeta^{i})$ dove le prime sono le coordinate rappresentative di x e $\zeta=\sum_{i}\zeta^{i}dx^{i}$. Siccome $\pi(x,\zeta)=x$, la rappresentazione in coordinate di $\alpha$ sarà
\[\alpha_{(x,\zeta)}=\pi^{*}(\sum_{i}\zeta^{i}dx^{i})=\sum_{i}\zeta^{i}dx^{i}\]
Segue subito che $\alpha$ è liscia dal momento che le sue funzioni che la compongono sono lineari.
Chiaramente poi $\omega=-d\alpha$ è simplettica visto che è chiusa (in particolare esatta), e mediante la naturale identificazione tra $T^{*}X$ ed un aperto di $\mathbb{R}^{2n}$ vedo che $\omega$ corrisponde alla forma simplettica standard (3), dunque è simplettica.
\end{proof}



\subsection{Sottovarietà Lagrangiane}

\begin{Def}
Sia $(M,\omega)$ una varietà simplettica 2n-dimensionale. Una sottovarietà Y di M si dice "Lagrangiana" se, per ogni $p\in Y$, $T_{p}Y$ è un sottospazio Lagrangiano di $T_{p}M$, ossia tale che $\omega_{p}|_{T_{p}Y}=0$ e $dimT_{p}Y=\frac{1}{2}dimT_{p}M$. In modo equivalente, se $i:Y\hookrightarrow M$ è la mappa di inclusione, allora Y è Lagrangiana se e solo se $i^{*}\omega=0$ e $dimY=\frac{1}{2}dimM$.
\end{Def}

\begin{Def}
Sia X una varietà differenziabile n-dimensionale, e $M=T^{*}M$ il suo fibrato cotangente.In riferimento alle notazioni precedenti, la "sezione nulla" di M è la sottovarietà 
\begin{equation}
X_{0}:=\{(x,\zeta) \in T^{*}X:\zeta=0\hspace{0.2cm}in \hspace{0.2cm}T_{x}^{*}X\}
\end{equation}
di $T^{*}X$ la cui intersezione con $T^{*}U\subseteq T^{*}X$ è data dalle equazioni $\zeta^{1}=...=\zeta^{n}=0$.
\end{Def}

\begin{flushleft}
Ovviamente la forma di Liouville $\alpha=\sum_{i}\zeta^{i}dx^{i}$ si annulla in $T^{*}U\cap X_{0}$, ed in particolare se $i_{0}:X_{0}\hookrightarrow T^{*}X$ è la mappa di inclusione, si ha $i_{0}^{*}d\alpha=i_{0}^{*}\omega=0$, dunque $X_{0}$ è Lagrangiana.

Ma quali sono le sottovarietà che "assomigliano" ad $X_{0}$?

Sia $X_{\mu}$ una sottovarietà n-dimensionale di $T^{*}X$ della forma
\begin{equation}
X_{\mu}:=\{(x,\mu):x\in X,\mu_{x}\in T_{x}^{*}X\}\subseteq T^{*}X
\end{equation}
dove il covettore $\mu_{x}$ dipende in modo differenziabile da x, e $\mu:X\rightarrow T^{*}X$ è la 1-forma di de Rham associata.Si ha allora la seguente 
\end{flushleft}

\begin{Prop}
Sia $X_{\mu}$ la sottovarietà appena definita, $s_{\mu}:X\rightarrow T^{*}X=M$ definita come $s_{\mu}(x):=(x,\mu_{x})$ ed $\alpha$ la forma di Liouville su $T^{*}X$. Allora
\[s_{\mu}^{*}\alpha=\mu .\]
\end{Prop}

\begin{proof}[Dimostrazione]
Dalla definizione di $\alpha$ ho $\alpha_{p}=\pi^{*}\zeta$ in $p=(x,\zeta)\in M$. Per $s_{\mu}(x)=(x,\mu_{x})=p$ ho allora $\alpha_{p}=(\pi^{*}\mu)_{p}$ e
\[(s_{\mu}^{*}\alpha)_{x}=s^{*}_{\mu}(\alpha_{s_{\mu}(x)})=s^{*}_{\mu}(\alpha_{(x,\mu_{x})})=s^{*}_{\mu}(\pi^{*}\mu_{x})\]
\[=\overbrace{(\pi\circ s_{\mu})}^{=id}^{*}\mu_{x}=\mu_{x}\]
\end{proof}

\begin{flushleft}
Si supponga che $X_{\mu}$ sia una sottovarietà n-dimensionale di $T^{*}X$ della forma definita in (8), con forma di de Rahm associata $\mu$. Allora $s_{\mu}:X\rightarrow T^{*}X$ è un embedding con immagine pari a $X_{\mu}$, e c'è un diffeomorfismo $\tau:X\rightarrow X_{\mu}$, dato da $\tau(x):=(x,\mu_{x})$, tale per cui il seguente diagramma 
\[\begin{tikzcd}
X \arrow[rr, "s_{\mu}"] \arrow[dr, "\tau", "\simeq"' ] &  & T^{*}X  \\
                                          & X_{\mu} \arrow[ru, "i"] &                               
\end{tikzcd}\]
commuti.

Osserviamo allora che $X_{\mu}$ è Lagrangiana
\[ \iff i^{*}\alpha=0\iff \tau^{*}i^{*}\alpha=0\iff (i\circ \tau)^{*}\alpha=0\iff s_{\mu}^{*}\alpha=0\]
\[\iff d(s_{\mu}^{*}\alpha)=0\iff d\mu=0 \]
ovverosia se e solo se $\mu$ è chiusa. Vi è dunque una corrispondenza biunivoca tra forme chiuse e sottovarietà Lagrangiane di $T^{*}X$ della forma (8).

Siano ora $(M_{1},\omega_{1})$ e $(M_{2},\omega_{2})$ due varietà simplettiche 2n-dimensionali;dato un diffeomorfismo $\varphi:M_{1}\rightarrow M_{2}$, ci si può chidedere quale sia la condizione affinchè esso risulti un simplettomorfismo. A tal scopo si considerino le due proiezioni
\[\begin{tikzcd}
  & M_{1} \\
  M_{1}\times M_{2} \arrow[ru, "\pi_{1}"] \arrow[rd, "\pi_{2}"] &\\
  & M_{2}
\end{tikzcd}\]
dove $\pi_{1}(p_{1},p_{2}):=p_{1}$ e $\pi_{2}(p_{1},p_{2}):=p_{2}$. La forma $\omega:=(\pi_{1})^{*}\omega_{1}+(\pi_{2})^{*}\omega_{2}$ è chiusa
\[d\omega=\pi_{1}^{*}d\omega_{1}+\pi_{2}^{*}d\omega_{2}=0\]
e simplettica.
\end{flushleft}

\begin{Def}
La "forma simplettica contorta" è la forma simplettica
\begin{equation}
\widetilde{\omega}:=\pi_{1}^*\omega_{1}-\pi_{2}^*\omega_{2}
\end{equation}
Invece il "grafico" $\Gamma_{\varphi}$ del diffeomorfismo $\varphi:M_{1}\rightarrow M_{2}$ è la sottovarietà 2n-dimensionale di $M_{1}\times M_{2}$ definita come
\begin{equation}
\Gamma_{\varphi}:=\{(p,\varphi(p)): p\in M_{1}\}.
\end{equation}
\end{Def}
Si osservi che il grafico $\Gamma_{\varphi}$ è l'immagine dell'embedding
\[\gamma:M_{1}\hookrightarrow M_{1}\times M_{2}\]
\[p\mapsto (p,\varphi(p))\]
\begin{Prop}
Il diffeomorfismo $\varphi$ è un simplettomorfismo se e solo se $\Gamma_{\varphi}$ è una sottovarietà lagrangiana di $(M_{1}\times M_{2}, \widetilde{\omega})$.
\end{Prop}
\begin{proof}[Dimostrazione]
Per definizione $\Gamma_{\varphi}$ è lagrangiana se e solo se $\gamma^{*}\widetilde{\omega}=0$.Ma
\[\gamma^{*}\widetilde{\omega}=\gamma^{*}(\pi_{1}^{*}\omega_{1})-\gamma^{*}(\pi_{2}^{*}\omega_{2})=(\pi_{1}\circ\gamma)^{*}\omega_{1}-(\pi_{2}\circ\gamma)^{*}\omega_{2}=id^{*}\omega_{1}-\varphi^{*}\omega_{2}\]
dunque 
\[\gamma^{*}\widetilde{\omega}=0\iff \omega_{1}=\varphi^{*}\omega_{2}\]
\[\begin{tikzcd}
M_{1} \arrow[rr, bend right, hook, "\gamma"]\arrow[r,"\varphi"] & M_{2}  & M_{1}\times M_{2} \arrow[ll, bend right,"\pi_{1}"] \arrow[l, bend right,"\pi_{2}"]\\ 
\end{tikzcd}\]
\end{proof}



\subsection{Teoremi di Moser}
Sia ora M una varietà liscia n-dimensionale, X una sottovarietà liscia k-dimensionale(k<n) e $i:X\hookrightarrow M$ la rispettiva mappa di inclusione. Ad ogni $x\in M$ lo spazio tangente di X è visto come un sottospazio di $T_{x}M$ mediante la mappa di inclusione $(i_{*})_{x}:T_{x}X\hookrightarrow T_{i(x)=x}M$.

\begin{Def}Il quoziente $N_{x}X:=T_{x}M/T_{x}X$ è uno spazio vettoriale (n-k)-dimensionale detto "Spazio normale di X in x".Il "fibrato normale" di X è
\begin{equation}
NX:=\{ (x,v):x\in X,v\in N_{x}X \} 
\end{equation}
\end{Def}

\begin{flushleft}
L'insieme NX ha una struttura di fibrato tangente di dimensione n-k sotto la proiezione naturale, dunque NX è una varietà n-dimensionale.La sezione nulla
\[i_{0}:X \hookrightarrow NX \]
\[x\mapsto (x,0)\]
costituisce un embedding tra X ed NX.
\end{flushleft}

\begin{Def}
Un intorno $U_{0}$ di NX si dice "convesso" se l'intersezione $U_{0}\cap N_{x}X$ con ogni fibra è convesso.
\end{Def}

\begin{Teo}[Degli intorni tubulari]
Esistono un intorno convesso $U_{o}$ di X in NX, un intorno U di X in M ed un diffeomorfismo $\varphi:U_{o}\rightarrow U$ tali per cui il seguente diagramma
\[\begin{tikzcd}
U_{0}\subseteq NX \arrow[rr, "\varphi","\simeq"']  &  & U\subseteq M  \\
                                          & X \arrow[ru, hook, "i"] \arrow[lu, hook, "i_{0}"] &                               
\end{tikzcd}\]
commuti.
\end{Teo}

\begin{flushleft}
$\textbf{NOTA}$: Restringendoci ad un sottoinsieme $U_{0}\in NX$, dal Teorema precedente, otteniamo una sommersione $\pi_{0}:U_{0}\rightarrow X$ con tutte le fibre $\pi^{-1}_{0}(x)$ convesse. Possiamo ottenere poi un fibrato da U definendo $\pi:=\pi_{0}\circ \varphi^{-1}$.
\[\begin{tikzcd}
U_{0}\subseteq NX \arrow[dr, "\pi_{0}"] &  & U\subseteq M \arrow[ll, "\varphi^{-1}"]\arrow[dl,"\pi"] \\
                                          & X &                               
\end{tikzcd}\]

Si consideri ora una varietà liscia M dotata di una mappa $\rho:M\times\mathbb{R}\rightarrow M$, e poniamo $\rho_{t}(p):=\rho(p,t)$.
\end{flushleft}

\begin{Def}
La mappa $\rho$ si dice "isotopia" se e solo se OGNI $\rho_{t}:M\rightarrow M$ è un diffeomorfismo e $\rho_{0}=id_{M}$.
\end{Def}

\begin{flushleft}
Data un isotopia $\rho$, si ottiene un campo vettoriale dipendente dal tempo: ossia una famiglia di campi vettoriali $V_{t}$ con $t\in\mathbb{R}$ tali per cui ad ogni $p\in M$ si abbia 
\[V_{t}(p)=\frac{d \rho_{s}(q)}{ds}\at[\big]{s=t}\hspace{0.5cm} q=\rho^{-1}_{t}(p)\]
Viceversa, dato un campo vettoriale dipendente dal tempo $V_{t}$, esiste (se M è compatta) un isotopia $\rho$ che soddisfa le relazioni precedenti.
Esponiamone ora un punto di vista un po' più euristico
\end{flushleft}

\begin{Def}
Un "campo vettoriale dipendente dal tempo" è una mappa 
\[V:\Omega\subset(\mathbb{R}\times M)\rightarrow TM\]
\[(t,p)\mapsto V_{t}(p)=V(t,p)\]
definita su di un aperto $\Omega$, tale per cui $\forall t\in\mathbb{R}$ l'insieme
\[\Omega_{t}:= \{p\in M : (t,p)\in\Omega \}\subset M \]
non sia vuoto.
Una "curva integrale" del campo è la mappa
\[\alpha:I\subset\mathbb{R}\rightarrow M\]
tale per cui, per ogni $t_{0}\in I$
\[(t_{0},\alpha(t_{0}))\in\Omega\hspace{0.5cm}e\hspace{0.5cm}\frac{d\alpha}{dt}\at[\big]{t=t_{0}}=V(t_{0},\alpha(t_{0}))\]
\end{Def}
Come per i campi vettoriali indipendenti dal tempo, possiamo definirne il flusso
\begin{Def}
Il "flusso" di un campo vettoriale dipendente dal tempo è l'unica mappa differenziabile
\[\theta:(J\subseteq\mathbb{R})\times\Omega\rightarrow M\]
tale per cui $\forall (t_{0},p)\in\Omega$ 
\[t\mapsto\theta(t,t_{0},p)\]
sia la curva integrale $\alpha$ di V che soddisfa $\alpha(t_{0})=p$.In altre parole $\theta(t_{0},t_{0},p)=p$ e
\[\frac{d}{dt}[\theta(t,t_{0},p)]\at[\big]{t=t_{0}}=V(t_{0},p)=V(t_{0},\theta(t_{0},t_{0},p))\]
Definiamo infine $\theta_{t,s}$ mediante la relazione $\theta_{t,s}(p):=\theta(t,s,p)$.
\end{Def}

\begin{flushleft}
$\textbf{NOTA}$: Valgono le seguenti due relazioni
\begin{equation}
\begin{aligned}
&\theta_{t_{2},t_{1}}\circ\theta_{t_{1},t_{0}}=\theta_{t_{2},t_{0}}\\
&\theta_{t,t}=id
\end{aligned}
\end{equation}
per ogni istante $t,t_{0},t_{1},t_{2}$ nel rispettivo dominio di definizione.Per capire da dove arriva la prima, ragioniamo in modo un po' sbrigativo ma piuttosto intuitivo:si ha che per unicità del flusso (si veda la Figura 2)
\[\theta_{t_{2},t_{1}}\circ\theta_{t_{1},t_{0}}(p)=\theta_{t_{2},t_{1}}(\theta(t_{1},t_{0},p))=\theta(t_{2},t_{1},\alpha(t_{1}))=\theta_{t_{2},t_{0}}(p)\]
Ma allora
\[\theta_{s,t}\circ\theta_{t,s}=\theta_{s,s}=id\Rightarrow\theta_{s,t}^{-1}=\theta_{t,s}\]
Segue che, definita l'isotopia $\theta_{t,0}=:\Theta_{t}$, posso scrivere 
\[\Theta_{t_{2}}\circ\Theta_{t_{1}}^{-1}=\theta_{t_{2},0}\circ\theta_{t_{1},0}^{-1}=\theta_{t_{2},0}\circ\theta_{0,t_{1}}=\theta_{t_{2},t_{1}}\]
\[\begin{tikzcd}
M \arrow[r,"\theta_{t_{1},0}^{-1}"] \arrow[rr,bend right,"\theta_{t_{2},t_{1}}"] & M \arrow[r,"\theta_{t_{2,0}}"] & M
\end{tikzcd}\]
Dunque 
\begin{equation}
V_{t}(p)=\frac{d}{d\epsilon}[\Theta_{t+\epsilon}\circ\Theta_{t}^{-1}(p)]\at[\bigg]{\epsilon=0}=\frac{d}{d\epsilon}[\theta_{t+\epsilon,t}(p)]\at[\bigg]{\epsilon=0}
\end{equation}
\end{flushleft}

\begin{Def}
Data una varietà differenziabile M,$V_{t}$ un campo vettoriale dipendente dal tempo su M, definiamo la "derivata di Lie" della forma $\alpha\in \mathcal{A}^{k}(M)$ 
\begin{figure}
    \centering
    \centerline{\includegraphics[scale=.03]{VF.jpg}}
    \caption{Rappresentazione grafica del campo $V_{t}$}
    \label{fig:my_label}
\end{figure}

\begin{equation}
\begin{split}
\mathcal{L}_{V_{t}}\alpha:=\frac{d}{d\epsilon}[\theta_{t+\epsilon,t}^{*}\alpha]\at[\bigg]{\epsilon=0}=\frac{d}{d\epsilon}[(\Theta_{t+\epsilon}\circ\Theta_{t}^{-1})^{*}\alpha]\at[\bigg]{\epsilon=0}=\frac{d}{d\epsilon}[(\Theta_{t}^{-1})^{*}(\Theta_{t+\epsilon})^{*}\alpha]\at[\bigg]{\epsilon=0}=\\
=\overbrace{(\Theta_{t}^{-1})^{*}}^{=(\Theta_{t}^{*})^{-1}}\frac{d}{d\epsilon}[(\Theta_{t+\epsilon}^{*})\alpha]\at[\bigg]{\epsilon=0}
\end{split}
\end{equation}
\end{Def}

Dimostriamo ora un risultato che ci servirà in seguito
\begin{Prop}
Per ogni campo vettoriale $X\in\mathcal{T}(M)$ ed ogni k-forma differenziale $\alpha\in\mathcal{A}^{k}(M)$ sulla varietà differenziabile M, vale la "formula magica di Cartan"
\begin{equation}
\mathcal{L}_{X}\alpha=X\lrcorner(d\alpha)+d(X\lrcorner\alpha)
\end{equation}
\end{Prop}
\begin{proof}
La dimostrazione procede per iduzione su k:
\begin{itemize}
    \item \underline{k=0}:iniziamo con una 0-forma $f$, per la quale si ha
    \[X\lrcorner(df)+d(X\lrcorner f)=X\lrcorner df=df(X)=X(f)=\mathcal{L}_{X}f\]
    dal momento che f non può essere contratta da X(ricordiamo che $\lrcorner$ abbassa di un grado le forme differenziali).Tale relazione coinicde con (15).
    \item \underline{k=1}:Siccome ogni 1-forma può essere scritta localmente come $udv$ per opportune funzioni liscie $u$ e $v$, possiamo limitarci ad $\alpha=udv$.In tal caso 
    \[\mathcal{L}_{X}(udv)=(\mathcal{L}_{X}u)dv+u(\mathcal{L}_{X}dv)=X(u)dv+ud(X(v))\]
    dove si è sfruttata la regola di Leibniz ed il seguente 
    \begin{Lem}
    Se $f\in C^{\infty}(M)$, allora $\mathcal{L}_{X}(df)=d(\mathcal{L}_{X}f)$.
    \end{Lem}
    \begin{proof}
    Partendo dalla nota formula (vedi appendice per una dimostrazione)
    \begin{equation}
    (\mathcal{L}_{X}\sigma)(Y_{1},...,Y_{k})=X(\sigma(Y_{1},...,Y_{k}))-\sigma([X,Y_{1}],Y_{2},...,Y_{k})-...-\sigma(Y_{1},...,Y_{k-1},[X,Y_{k}])
    \end{equation}
    dove $\sigma$ è un tensore covariante ed $Y_{i},X$ sono campi vettoriali , si ha immediatamente
    \[(\mathcal{L}_{X}df)(Y)=X(df(Y))-df([X,Y])=XY(f)-[X,Y](f)=\]
    \[=XY(f)-(XY-YX)(f)=YX(f)=d(Xf)(Y)=d(X(f))(Y)=\]
    \[=d(\mathcal{L}_{X}f)(Y)\]
    \end{proof}
    Tornando a noi, si ha anche che (sfruttando il fatto $X \lrcorner$ è un antiderivazione)
    \[X\lrcorner d(udv)+\overbrace{d(X\lrcorner(udv))}^{=d(uX\lrcorner dv)=d(uX(v))}=X\lrcorner(du\wedge dv)+\overbrace{d(uX(v))}^{d(fg)=fdg+gdf}=\]
    \[=(X\lrcorner du)\wedge dv+(-1)^{1}du\wedge(X\lrcorner dv)+ud(X(v))+(X(v))du=\]
    \[=(X(u))dv-(X(v))du+ud(X(v))+(X(v))du=(X(u))dv+ud(X(v))\]
    \[=\mathcal{L}_{X}(udv)\]
    \item\underline{k>1}:Sia ora k>1, e supponiamo che (15) sia soddisfatta per k-1.Sia ora $\alpha$ una k-forma arbitraria, espressa in coordinate locali come
    \[\alpha=\sum_{I}^{'}=\alpha_{I}dx^{i_{1}}\wedge...\wedge dx^{i_{k}}=\sum_{I}^{'}=\alpha_{I}dx^{I}\]
    dove la somma è vincolata agli indici $1\leq i_{1}< i_{2}< ...< i_{k}\leq n=dim(M)$.
    Ponendo $\eta:=\alpha dx^{i_{1}},\gamma:=dx^{i_{2}}\wedge...\wedge dx^{i_{k}}$ vedo che posso esprimere $\alpha$ come somma di termini del tipo $\eta\wedge\gamma$, dove $\eta$ è una 1-forma e $\gamma$ una (k-1)-forma sulla quale possiamo applicare l'ipotesi induttiva;quest'ultima in particolare implica che 
    \[\mathcal{L}_{X}(\eta\wedge\gamma)=(\mathcal{L}_{X}\eta)\wedge\gamma+\eta\wedge(\mathcal{L}_{X}\gamma)=\]
    \[=(X\lrcorner d\eta+d(X\lrcorner\eta))\wedge\gamma+\eta\wedge(X\lrcorner d\gamma+d(X\lrcorner\gamma))=\spadesuit\]
    D'altra parte, usando il fatto che sia d che $X\lrcorner$ sono antiderivazioni, possiamo scrivere
    \[X\lrcorner d(\eta\wedge\gamma)+d(X\lrcorner(\eta\wedge\gamma))=\]
    \[=X\lrcorner (d\eta\wedge\gamma+(-1)^{1}\eta\wedge d\gamma)+d((X\lrcorner\eta)\wedge\gamma+(-1)^{1}\eta\wedge(X\lrcorner\gamma))=\]
    \[=(X\lrcorner d\eta)\wedge\gamma+\textcolor{orange}{d\eta\wedge(X\lrcorner \gamma)}\textcolor{blue}{-(X\lrcorner\eta)\wedge d\gamma}+\eta\wedge(X\lrcorner d\gamma)+\]
    \[+d(X\lrcorner\eta)\wedge\gamma\textcolor{blue}{+(X\lrcorner\eta)\wedge d\gamma}\textcolor{orange}{-d\eta\wedge(X\lrcorner\gamma)}+\eta\wedge d(X\lrcorner\gamma)=\]
    \[=\spadesuit\]
    dove si ricordi che $\eta$ è una 1-forma dunque $d\eta$ è una 2-forma.
\end{itemize}
\end{proof}

\begin{Cor}
Se X è un campo vettoriale ed $\alpha$ è una forma differenziale, allora
\[\mathcal{L}_{X}(d\alpha)=d(\mathcal{L}_{X}\alpha)\]
\end{Cor}

\begin{proof}
L'enunciato segue dalla Proposizione precedente e dal fatto che $d^{2}=0$:da una parte
\[\mathcal{L}_{X}d\alpha=X\lrcorner d(d\alpha)+d(X\lrcorner d\alpha)=d(X\lrcorner d\alpha)\]
mentre 
\[d(\mathcal{L}_{X}\alpha)=d(X\lrcorner d\alpha+d(X\lrcorner\alpha))=d(X\lrcorner d\alpha)\]
\end{proof}

Vogliamo ora dimostrare il seguente
\begin{Teo}[Assoluto di Moser]
Sia $(M,\omega_{t}=(1-t)\omega_{0}+t\omega_{1})$ una varietà simplettica $\forall t\in[0,1]$ e tale che $[\omega_{0}]=[\omega_{1}]\in H_{dR}^{2}(M)$. Allora esiste una mappa $\rho:M\times\mathbb{R}\rightarrow M$ tale per cui
\[\rho_{t}:M\xrightarrow{\sim}M\hspace{0.5cm}\rho_{t}^{*}\omega_{t}=\omega_{0}\hspace{0.2cm}\forall t\in[0,1]\]
\end{Teo}

\begin{proof}
Partiamo anzitutto dal seguente 

\begin{Lem}
Esiste un unico campo vettoriale $V_{t}\in\mathcal{T}(M)$ dipendente in modo liscio da t che soddisfa l'equazione
\begin{equation}
\mathcal{L}_{V_{t}}\omega_{t}+\frac{d\omega_{t}}{dt}=0
\end{equation}
\end{Lem}

\begin{proof}
Si ha $\frac{d\omega_{t}}{dt}=\omega_{1}-\omega_{0}$ e dal momento che $[\omega_{1}]=[\omega_{0}]$ si ha $\omega_{1}-\omega_{0}=d\mu$ per qualche $\mu\in\mathcal{A}(M)$.
Dalla (15) so poi che
\[\mathcal{L}_{V_{t}}\omega_{t}=(V_{t}\lrcorner d+d (V_{t}\lrcorner )\omega_{t}=d(V_{t}\lrcorner\omega_{t})\]
perchè $\omega_{t}$ è combinazione lineare delle $\omega_{0,1}$ che sono chiuse per definizione. Dunque la (17) si riduce alla 
\[d(V_{t}\lrcorner\omega_{t}+\mu)=0\]così
\[d(V_{t}\lrcorner\omega_{t})+d\mu=0\iff\mathcal{L}_{V_{t}\omega_{t}}+\frac{d\omega_{t}}{dt}=0\]
E' dunque sufficiente risolvere $V_{t}\lrcorner\omega_{t}+\mu=0$, ma essnedo $\omega_{t}$ non degenere possiamo risolvere tale equazione punto per punto, ottenendo un unico $V_{t}$.
\end{proof}

In particolare, per la definizione data in (14)
\[\overbrace{\frac{d}{d\epsilon}(\rho_{t+\epsilon}^{*}\omega_{t})\at[\bigg]{\epsilon=0}}^{=\frac{d}{dt}(\rho_{t}^{*}\omega_{t})}=\rho_{t}^{*}\mathcal{L}_{V_{t}\omega_{t}}\]
A questo punto però ci avvaliamo del seguente

\begin{Lem}
Per ogni famiglia di d-forme differenziabili $\omega_{t},t\in\mathbb{R}$ si ha
\[\frac{d}{dt}(\rho_{t}^{*}\omega_{t})=\rho_{t}^{*}(\mathcal{L}_{V_{t}}\omega_{t}+\frac{d\omega_{t}}{dt})\]
\end{Lem}

\begin{proof}
Se $f(x,y)$ è una funzione reale in due variabili, per la chain rule
\[\overbrace{\frac{d}{dt}f(t,t)}^{=\frac{d}{dt}f(x(t),y(t))}=\nabla f(t,t)\cdot (1,1)=\frac{\partial}{\partial x}f(x,t)\at[\bigg]{x=t}+\frac{\partial}{\partial y}f(t,y)\at[\bigg]{y=t}\]
Dunque nel nostro caso
\[\frac{d}{dt}(\rho_{t}^{*}\omega_{t})=\underbrace{\frac{\partial}{\partial x}\rho_{x}^{*}\omega_{t}\at[\bigg]{x=t}}_{=\rho_{t}^{*}\mathcal{L}_{V_{t}}\omega_{t} per (14)}+\underbrace{\frac{\partial}{\partial y}\rho_{t}^{*}\omega_{y}\at[\bigg]{y=t}}_{=\rho_{t}^{*}\frac{d\omega_{y}}{dy}\at[\big]{y=t}}=\rho_{t}^{*}(\mathcal{L}_{V_{t}}\omega_{t}+\frac{d\omega_{t}}{dt})\]
\end{proof}

In questo modo
\[0=\rho_{t}^{*}(\mathcal{L}_{V_{t}}\omega_{t}+\frac{d\omega_{t}}{dt})=\frac{d}{dt}(\rho_{t}^{*}\omega_{t})\]
dove $\rho_{t}$ è l'isotopia associata al campo vettoriale $V_{t}$. Infine 
\[\frac{d}{dt}(\rho_{t}^{*}\omega_{t})=0\Rightarrow \rho_{t}^{*}\omega_{t}=\rho_{0}^{*}\omega_{0}=\omega_{0}\]
\end{proof}

\begin{flushleft}
Si osservi che quest'ultimo risultato può essere esteso ad una qualsiasi famiglia di forme simplettiche $\frac{d}{dt}[\omega_{t}]=0\in H_{dR}^{2}(M)$.
Per dimostrare l'ultimo risultato è necessaria la seguente 
\end{flushleft}

\begin{Prop}[Formula omotopica]
Sia U un intorno tubulare di una sottovarietà X in M.Se una d-forma $\omega$ chiusa su U soddisfa $i^{*}\omega=0$, allora $\omega$ è esatta ($\omega=d\mu$ per un opportuna $\mu\in\mathcal{A}^{d-1}(U)$).Inoltre, possiamo scegliere $\mu$ tale per cui $\mu\at{x}=0\hspace{0.3cm}\forall x\in X$.
\end{Prop}
\begin{proof}
Siccome $\varphi:U\xrightarrow{\sim} U_{0}$, lavorare su $U_{0}$ sarà equivalente.Definiamo ora, per ogni $t\in[0,1]$, la mappa
\[\rho_{t}:U_{0}\rightarrow U_{0}\]
\[(x,v)\mapsto(x,tv)\]
Essa è ben definita dal momento che $U_{0}$ è, per ipotesi, convesso.La mappa $\rho_{1}$ è l'identità, $\rho_{0}=i_{0}\circ\pi_{0}$, ed ogni $\rho_{t}$ fissa X ($\rho_{t}\circ i_{0}=i_{0}$). Diciamo allora che la famiglia $\{\rho_{t}:0\leq t\leq 1 \}$ è un "omotopia" tra $i_{0}\circ\pi_{0}$ ed $id_{X}$. La mappa $\pi_{0}:U_{0}\rightarrow X$ è chiamata una "ritrazione" poichè $i_{0}\circ\pi_{0}$ è l'identità (stiamo sfruttando i risultati del Teorema 1.2).

Un "operatore omotopico (di de Rham)" tra $\rho_{0}=i_{0}\circ\pi_{0}$ e $\rho_{1}=id$ è una mappa lineare
\[Q:\mathcal{A}^{d}(U_{0})\rightarrow\mathcal{A}^{d-1}(U_{0})\]
che soddisfa la "formula omotopica"
\begin{equation}
Id-(i_{0}\circ\pi_{0})^{*}=dQ+Qd
\end{equation}
Quando $d\omega=0$ e $i^{*}_{0}\omega=0$ (che nel nostro caso è vero per ipotesi), l'operatore Q fornisce $\omega=dQ(\omega)$, perciò possiamo prendere $\mu:=Q(\omega)$.

\begin{Lem}
L'operatore 
\begin{equation}
Q(\omega):=\int_{0}^{1}\rho_{t}^{*}(V_{t}\lrcorner\omega)dt
\end{equation}
dove $V_{t}(q)=V_{t}(\rho_{t}(p))$ è il vettore tangente alla curva $\rho_{s}(p)$ al tempo s=t, è un operatore omotopico.
\end{Lem}
\begin{proof}
Basta dimostrare che Q soddisfa la formula omotopica, ma in effetti si ha (la derivata esterna commuta con il pullback)
\[Q(d\omega)+dQ(\omega)=\int_{0}^{1}\rho_{t}(V_{t}\lrcorner d\omega)dt+d\int_{0}^{1}\rho_{t}^{*}(V_{t}\lrcorner\omega)dt=\]
\[=\int_{0}^{1}\rho_{t}^{*}\underbrace{(V_{t}\lrcorner d\omega+d(V_{t}\lrcorner\omega))}_{=\mathcal{L}_{V_{t}\omega}}dt\]
Siccome poi
\[\frac{d}{dt}(\rho_{t}^{*}\omega)=\rho_{t}\mathcal{L}_{V_{t}}\omega\]
dal Teorema fondamentale del calcolo integrale si avrà
\[Q(d\omega)+dQ(\omega)=\int_{0}^{1}\frac{d}{dt}(\rho_{t}^{*}\omega)dt=\rho_{1}\omega-\rho_{0}\omega=(Id-(i_{0}\circ\pi_{0})^{*})\omega\]
\end{proof}
Nel nostro caso allora, per $x\in X$, $\rho_{t}(x)=x$ è la curva costante e allora $V_{t}$ si annulla in ogni punto di X ed in ogni isante; dunque $\mu_{x}=Q(\omega)\at{x}=0$.
\end{proof}

Un ulteriore risultato che ci serve viene esposto nel seguente

\begin{Teo}[Relativo di Moser]
Siano $(M,\omega_{0})$ ed $(M,\omega_{1})$ due strutture simplettiche nella stessa varietà differenziabile, sia $i:X\hookrightarrow M$ l'embedding associato alla sottovarietà compatta X, e si supponga che $\omega_{0}\at{p}=\omega_{1}\at{p}$ $\forall p\in i(X)$.Allora esistono due intorni $U_{0},U_{1}$ di $i(X)$ ed un diffeomorfismo $\varphi:U_{0}\xrightarrow{\sim} U_{1}$ tale per cui il seguente diagramma
\[\begin{tikzcd}
U_{0} \arrow[rr, "\sim","\varphi"']  &  & U_{1}  \\
                                          & X \arrow[ru, hook, "i"] \arrow[lu, hook, "i"] &                               
\end{tikzcd}\]
commuti, e $\varphi^{*}\omega_{1}=\omega_{0}$.
\end{Teo}
\begin{proof}
Si scelga un intorno tubulare $U_{0}\subseteq M$ di X. La 2-forma $\omega_{1}-\omega_{0}$ è chiusa in $U_{0}$ e per ipotesi $(\omega_{1}-\omega_{0})_{p}=0$ $\forall p\in X$.Per la formula omotopica (Proposizione 1.7) esiste una 1-forma $\mu$ su $U_{0}$ tale che $\omega_{1}-\omega_{0}=d\mu$ e $\mu_{p}=0$ in ogni $p\in X$.

Si consideri la famiglia $\omega_{t}=(1-t)\omega_{0}+t\omega_{1}=\omega_{0}+td\mu$ di 2-forme chiuse in $U_{0}$. Restringendo quest'ultimo (se necessario), possiamo assumere che $\omega_{t}$ sia simplettica $\forall t\in[0,1]$.

A questo punto risolvendo "l'equazione di Moser"
\begin{equation}
    V_{t}\lrcorner \omega_{t}=-\mu
\end{equation}
si vede che deve essere $V_{t}=0$:questo perchè $\omega_{t}$ è non degenere
\[\omega_{t}(V_{t},X)=0\hspace{0.5cm}\forall X\in\mathcal{T}(M)\iff V_{t}=0\]
mentre $\mu\at{p}=0$ implica
\[\omega_{t}(V_{t},X)\at{p}=V_{t}\lrcorner\omega_{t}\at{p}(X)=-\mu\at{p}(X)=0\hspace{0.4cm}\forall X\Rightarrow V_{t}=0\]

Il Teorema assoluto di Moser garantisce poi l'esistenza di un isotopia $\rho_{t}:U_{0}\times[0,1]\rightarrow M$ tale per cui $\rho_{t}^{*}\omega_{t}=\omega_{0}$.Inoltre dal momento che $V_{t}=0$ si avrà $\rho_{t}\at{X}=id_{X}$ $\forall t\in[0,1]$. 
Possiamo allora scegliere $\varphi:=\rho_{1}$ e $U_{1}:=\rho_{1}(U_{0})$.
\end{proof}

Terminiamo dunque la sezione con un risultato fondamentale
\begin{Teo}[Di Darboux]
Sia $(M,\omega)$ una varietà simplettica, e $p\in M$ un punto qualsiasi della varietà. Allora possiamo trovare un sistema di coordinate $(U,x^{1},...,x^{n},y^{1},...,y^{n})$ centrate in p tali per cui, in U, si abbia
\[\omega=\sum_{i=1}^{n}dx^{i}\wedge dy^{i}\]
\end{Teo}
\begin{proof}
Basta applicare il Teorema 1.4 ad $X\equiv\{p\}$:si usi infatti una base simplettica di $T_{p}M$ per costruire delle coordinate $(x'^{i},...,x'^{n},y'^{1},...,y'^{n})$ centrate in p ed associate ad un opportuno intorno U'.

Ci sono allora due forme simplettiche su U': $\omega=:\omega_{0}$ ed $\omega_{1}:=\sum_{i}dx'^{i}\wedge dy'^{i}$. Dal Teorema 1.4 esistono allora due intorni $U_{0},U_{1}$ di p ed un diffeomorfismo $\varphi:U_{0}\rightarrow U_{1}$ tali per cui
\[\varphi(p)=p\hspace{0.5cm}e\hspace{0.5cm}\varphi^{*}(dx'^{i}\wedge dy'^{i})=\omega \]
Dal momento che $\varphi^{*}(\sum_{i}dx'^{i}\wedge dy'^{i})=\sum_{i}d(x'^{i}\circ\varphi)\wedge d(y'^{i}\circ\varphi)$, per concludere basterà usare le nuove coordinate $x^{i}:=x'^{i}\circ\varphi$ e $y^{i}:=y'^{i}\circ\varphi$. 
\end{proof}



\subsection{Struttura Hamiltoniana}
Utilizziamo quanto appena visto per definire i sistemi hamiltoniani su varietà simplettiche: anzitutto ricordiamo che, data una varietà simplettica $(M,\omega)$, è ben definita la mappa
\[\widetilde{\omega}:\mathcal{T}(M)\rightarrow\mathcal{A}(M)\]
\[X\mapsto\widetilde{\omega}(X)\]
dove $\widetilde{\omega}(X)(Y):=\omega(X,Y)$ per ogni $X,Y\in\mathcal{T}(M)$.
\begin{Def}
Sia $(M,\omega)$ una varietà simplettica. Per ogni funzione liscia $h\in C^{\infty}(M)$, definiamo il "campo vettoriale Hamiltoniano" quel campo vettoriale
\begin{equation}
X_{h}:=\widetilde{\omega}^{-1}(dh)    
\end{equation}
in modo che 
\begin{equation}
    \omega(X_{h},Y)=\overbrace{\widetilde{\omega}(X_{h})}^{=(\widetilde{\omega}\circ\widetilde{\omega}^{-1})(dh)}(Y)=dh(Y)=Y(h)
\end{equation}
per ogni scelta del campo vettoriale $Y\in\mathcal{T}(M)$.Un alternativa è definirlo come
\begin{equation}
X_{h}\lrcorner\omega=dh
\end{equation}
\end{Def}


\begin{flushleft}
$\underline{Esempio}$:In $\mathbb{R}^{2n}$ dotato della forma simplettica standard $\omega=\sum_{i=1}^{n}dx^{i}\wedge dy^{i}$, il campo $X_{h}$ può essere esplicitato come segue:scrivendo
\[X_{h}=\sum_{j=1}^{n}(A^{j}\frac{\partial}{\partial x^{j}}+B^{j}\frac{\partial}{\partial y^{j}})\]
per opportuni coefficenti $A^{j},B^{j}$ da determinare,si ha da una parte (si ricordi che $X\lrcorner $ è un antiderivazione per ogni campo X)
\[X_{h}\lrcorner\omega=[\sum_{j=1}^{n}(A^{j}\frac{\partial}{\partial x^{j}}+B^{j}\frac{\partial}{\partial y^{j}})]\lrcorner(\sum_{i=1}^{n}dx^{i}\wedge dy^{i})=\]
\[=\sum_{i,j=1}^{n}[(A^{j}\frac{\partial}{\partial x^{j}}+B^{j}\frac{\partial}{\partial y_{j}})\lrcorner dx^{i}]\wedge dy^{i}+(-1)^{1}dx^{i}\wedge[(A^{j}\frac{\partial}{\partial x^{j}}+B^{j}\frac{\partial}{\partial y^{j}})\lrcorner dy^{i}]=\]
\[=\sum_{i,j=1}^{n}(A^{j}\delta_{ij}dy^{i}-B^{j}\delta_{ij}dx^{i})=\sum_{i=1}^{n}(A^{i}dy^{i}-B^{i}dx^{i})\]
D'altra parte però
\[dh=\sum_{i=1}^{n}(\frac{\partial h}{\partial x^{i}}dx^{i}+\frac{\partial h}{\partial y^{i}}dy^{i})\]
e dunque, uguagliando le due, si ottiene che 
\[A^{j}=\frac{\partial h}{\partial y^{j}}\hspace{0.5cm}B^{j}=-\frac{\partial h}{\partial x^{j}}\]
cioè
\[X_{h}=\sum_{j=1}^{n}(\frac{\partial h}{\partial y^{j}}\frac{\partial}{\partial x^{j}}-\frac{\partial h}{\partial x^{j}}\frac{\partial}{\partial y^{j}})\]
\end{flushleft}

\begin{Prop}
Il gruppo ad un parametro di diffeomorfismi associato al campo vettoriale Hamiltoniano preserva la struttura simplettica.
\end{Prop}
\begin{proof}
Sia infatti $\theta_{t}:M\rightarrow M$ il flusso associato al campo $X_{h}$, dal momento che $\omega$ è per definizione chiusa 
\[\Rightarrow \frac{d}{dt}[(\theta_{t}^{*}\omega)]=\theta_{t}^{*}(\mathcal{L}_{X_{h}}\omega)=\theta_{t}^{*}(d(X_{h}\lrcorner\omega)+X_{h}\lrcorner \overbrace{d\omega}^{=0})=\theta_{t}^{*}d^{2}h=0\]
dove nel primo passaggio si è usato il Lemma 1.7 dimostrato poco più avanti.
\end{proof}

\begin{flushleft}
$\underline{Esempio}$: Sia $M=\mathbb{S}^{2}$ con $\omega_{p}(X,Y):=<p,X\times Y>=p\cdot(X\times Y)$ e descriviamone l'embedding in $\mathbb{R}^{3}$ come
\[\iota:\mathbb{S}^{2}\rightarrow\mathbb{R}^{3}\]
\[(\theta,h)\mapsto(\sqrt{1-h^{2}}\cos\theta, \sqrt{1-h^{2}}\sin\theta, h)=:(x(\theta,h),y(\theta,h),z(\theta,h))\]
In tal caso, ricrodando come traformano i campi vettoriali 
\[\iota_{*}(\frac{\partial}{\partial x^{i}}\at[\bigg]{p})=\frac{\partial\iota^{j}}{\partial x^{i}}(p)\frac{\partial}{\partial y^{j}}\at[\bigg]{\iota(p)}\]
si ottiene (si sono denominate con $x^{i}$ le variabili in $\mathbb{S}^{2}$ e con $y^{i}$ quelle in $\mathbb{R}^{3}$)
\[\iota_{*}(\frac{\partial}{\partial\theta})=\frac{\partial x}{\partial\theta}\frac{\partial}{\partial x}+\frac{\partial y}{\partial\theta}\frac{\partial}{\partial y}+\frac{\partial z}{\partial\theta}\frac{\partial}{\partial z}=-y\frac{\partial}{\partial x}+x\frac{\partial}{\partial y}\]
\[\iota_{*}(\frac{\partial}{\partial h})=\frac{\partial x}{\partial h}\frac{\partial}{\partial x}+\frac{\partial y}{\partial h}\frac{\partial}{\partial y}+\frac{\partial z}{\partial h}\frac{\partial}{\partial z}=-\frac{h}{\sqrt{1-h^{2}}}\cos\theta\frac{\partial}{\partial x}-\frac{h}{\sqrt{1-h^{2}}}\sin\theta\frac{\partial}{\partial y}+\]
\[+\frac{\partial}{\partial z}=\frac{-xz}{1-z^{2}}\frac{\partial}{\partial x}+\frac{-yz}{1-z^{2}}\frac{\partial}{\partial y}+\frac{\partial}{\partial z}\]
\[\Rightarrow\iota_{*}(\frac{\partial}{\partial\theta})\times\iota_{*}(\frac{\partial}{\partial h})=(-1)^{1+1}Det(A_{11})\frac{\partial}{\partial x}+(-1)^{2+1}Det(A_{21})\frac{\partial}{\partial y}+(-1)^{3+1}Det(A_{31})\frac{\partial}{\partial z}=\]
\[=x\frac{\partial}{\partial x}+y\frac{\partial}{\partial y}++\frac{(x^{2}+y^{2})z}{1-z^{2}}\frac{\partial}{\partial z}=x\frac{\partial}{\partial x}+y\frac{\partial}{\partial y}+z\frac{\partial}{\partial z}\]
e dunque la struttura simplettica standard su $\mathbb{S}^{2}$ è $\omega=d\theta\wedge dh$.
Si è usata la matrice 
\[A:=\begin{pmatrix}
e_{1} & -y & \frac{-xz}{1-z^{2}}\\
e_{2} & x & \frac{-yz}{1-z^{2}} \\
e_{3} & 0 & 1
\end{pmatrix}\]
\end{flushleft}

\begin{Prop}
Sia $(M,\omega)$ una varietà simplettica e sia $h\in C^{\infty}(M)$.Allora
\begin{itemize}
    \item [(i)] h è costante lungo il flusso di $X_{h}$, ovvero se $\theta$ è il suddetto flusso, si ha $h\circ\theta_{t}(p)=h(p)$ per ogni $(t,p)$ nel dominio di $\theta$.
    \item [(ii)] Ad ogni punto di h, il campo vettoriale $X_{h}$ è tangente all'insieme di livello di h.
\end{itemize}
\end{Prop}
\begin{proof}
Partiamo dal seguente Lemma, che da un analogo della relazione (14) ma per campi autonomi
\begin{Lem}
Sia M una varietà liscia, $X\in\mathcal{T}(M)$ e sia $\theta$ il suo flusso. Allora per ogni campo tensoriale $\tau$ ed ogni $(t_{0},p)$ nel dominio di $\theta$ vale
\begin{equation}
    \frac{d}{dt}\at[\bigg]{t=t_{0}}\theta_{t}^{*}(\tau_{\theta_{t_{0}}(p)})=\theta_{t_{0}}^{*}(\mathcal{L}_{X}\tau)_{\theta_{t_{0}}(p)}
\end{equation}
\end{Lem}
\begin{proof}
Si ha che il cambio di variabili $t=t_{0}+s$ implica
\[\frac{d}{dt}\at[\bigg]{t=t_{0}}\theta_{t}^{*}(\tau_{\theta_{t_{0}}(p)})=\frac{d}{ds}\at[\bigg]{s=0}\overbrace{(\theta_{t_{0}+s}^{*})}^{=(\theta_{s}\circ\theta_{t_{0}})^{*}}\tau_{\theta_{s+t_{0}}(p)}\]
\[=\frac{d}{ds}\at[\bigg]{s=0}(\theta_{t_{0}})^{*}(\theta_{s})^{*}\tau_{\theta_{s}(\theta_{t_{0}}(p))}\]
\[=(\theta_{t_{0}})^{*}\frac{d}{ds}\at[\bigg]{s=0}(\theta_{s})^{*}\tau_{\theta_{s}(\theta_{t_{0}}(p))}\]
\[=(\theta_{t_{0}})^{*}(\mathcal{L}_{X}\tau)_{\theta_{t_{0}}(p)}\]
\end{proof}
A questo punto possiamo dimostrare un ulteriore Lemma che ci servirà
\begin{Lem}
Sia M una varietà differenziabile e $X\in\mathcal{T}(M)$. Un campo tensoriale $\tau$ su M è invariante sotto il flusso di X se e solo se $\mathcal{L}_{X}\tau=0$.
\end{Lem}
\begin{proof}
Sia $\theta$ il flusso associato ad X. Se $\tau$ è invariante rispetto a $\theta$, allora $\theta_{t}^{*}\tau=\tau$ per ogni istante t. Inserendo quest'ultima relazione nella definizione della derivata di Lie, vedo che  $\mathcal{L}_{X}\tau=0$.

Viceversa si supponga che $\mathcal{L}_{X}\tau=0$. Per ogni punto p, si denoti con $\mathcal{D}_{p}$ il dominio di $\theta^{(p)}$. Si consideri poi la curva
\[T:\mathcal{D}_{p}\rightarrow T^{k}(T_{p}M)\]
\[t\mapsto T(t):=\theta_{t}^{*}(\tau_{\theta_{t}(p)})\]
Il Lemma precedente implica allora che $T'(t)=0\hspace{0.4cm}\forall t\in\mathcal{D}_{p}$. In altre parole $T(t)=T(0)=\tau_{p}$ per istante t. Ciò equivale a dire che 
\[\theta_{t}^{*}(\tau_{\theta_{t}(p)})=\tau_{p}\]
che è quanto voluto.
\end{proof}
Tornando ora alla Proposizione 1.9 iniziale,entrambe le affermazioni seguono dal Lemmma precedente assieme al fatto che
\[\mathcal{L}_{X_{h}}(h)=X_{h}(h)=dh(X_{h})=\omega(X_{h},X_{h})=0\]
essendo $\omega$ antisimmetrica.
\end{proof}

\begin{Def}
Data una varietà simplettica $(M,\omega)$, un campo vettoriale $X\in\mathcal{T}(M)$ si dice "simplettico" se e solo se soddisfa $\mathcal{L}_{X}\omega=0$. Si dice invece essere "Hamiltoniano" se esiste una funzione $f\in C^{\infty}(M)$ tale per cui $X=X_{f}$, mentre si dice "localmente Hamiltoniano" se ogni punto p ha un intorno in cui $X=X_{f}$.
\end{Def}

$\textbf{NOTA}$:Si noti che
\begin{itemize}
    \item Chiaramente ogni campo Hamiltoniano è anche localmente Hamiltoniano;
    \item X è simplettico $\iff d(X\lrcorner\omega)=0$;
    \item X è Hamiltoniano $\iff X\lrcorner\omega$ è esatta;
    \item $\Rightarrow H_{dR}^{1}(M)=0$ è la condizione affinchè ogni campo vettoriale simplettico sia Hamiltoniano (cioè ogni 1-forma $X\lrcorner\omega$ sia esatta).
\end{itemize}

\begin{Def}
Sia $(M,\omega)$ una varietà simplettica, $f,g\in C^{\infty}(M)$ due funzioni. Definiamo le loro "parentesi di Poisson" 
\[\{f,g\}:=X_{f}(g)=\omega(X_{g},X_{f})\]
Si dice poi che le due funzioni f,g "commutano" se $\{f,g\}=0$.
\end{Def}

\begin{Prop}
Sia $(M,\omega)$ una varietà simplettica, ed $f,g\in C^{\infty}(M)$. Allora valgono le seguenti affermazioni
\begin{itemize}
    \item [(i)]$\{f,g\}=-\{g,f\}$;
    \item [(ii)] g è costante lungo il flusso di $X_{f}$ se e solo se $\{f,g\}=0$;
    \item [(iii)] $X_{\{f,g\}}=[X_{f},X_{g}]$.
\end{itemize}
\end{Prop}
\begin{proof}
La (i) è ovvia dal momento che $\{f,g\}=\omega(X_{f},X_{g})$ e la (ii) dal fatto che $\{f,g\}=X_{f}(g)=\mathcal{L}_{X_{f}}g$. Per la (iii) si ha
\[d(\{f,g\})=d(X_{f}(g))=d(\mathcal{L}_{X_{f}}(g))\stackrel{Cor. 1.3}{=}\mathcal{L}_{X_{f}}dg=\mathcal{L}_{X_{f}}(X_{g}\lrcorner\omega)\]
\[\stackrel{(1)}{=}(\mathcal{L}_{X_{f}}X_{g})\lrcorner\omega+X_{g}\lrcorner(\mathcal{L}_{X_{f}}\omega)\stackrel{(2)}{=}[X_{f},X_{g}]\lrcorner\omega\]
dove si è sfruttata la relazione (1) dimostrata in appendice.Inoltre in (2) abbiamo sfruttato il fatto che
\[X_{g}\lrcorner(\mathcal{L}_{X_{f}}\omega)\stackrel{(15)}{=}X_{g}\lrcorner(X_{f}\lrcorner d\omega)+X_{g}\lrcorner d(X_{f}\lrcorner\omega)=X_{g}\lrcorner d^{2}f=0\]
dal momento che $\omega$ è chiusa per definizione.
\end{proof}

Possiamo dunque riassumere le proprietà delle parentesi di Poisson nel seguente

\begin{Teo}
Le parentesi di Poisson soddisfano le seguenti proprietà:
\begin{itemize}
    \item [(a)]ANTISIMMETRIA:$\{f,g\}=-\{g,f\}$;
    \item [(b)]BILINEARITA':\{$\lambda_{1}f+\lambda_{2}g$,h\}=$\lambda_{1}$\{f,h\}+$\lambda_{2}$\{g,h\} per ogni costante $\lambda_{1,2}\in\mathbb{R}$;
    \item [(c)]REGOLA DI LEIBNIZ:$\{f,gh\}=g\{f,h\}+h\{f,g\}$;
    \item [(d)]IDENTITA' DI JACOBI:
    \begin{equation}
    \{f,\{g,h\}\}+\{h,\{f,g\}\}+\{g,\{h,f\}\}=0
    \end{equation}
\end{itemize}
\end{Teo}
\begin{proof}
La (a) è già stata dimostrata mentre la (b) è ovvia. Per la (c) invece, dalla definizione
\[\{f,gh\}=X_{f}(gh)=gX_{f}(h)+hX_{f}(g)=g\{f,h\}+h\{f,g\}\]
essendo $X_{f}$ una derivazione (in ogni punto della varietà). Infine per la (d) bastà sfruttare le proprietà appena esposte assieme alla Prop.1.10(iii)
\[\{f,\{g,h\}\}+\textcolor{brown}{\{h,\{f,g\}\}}+\{g,\{h,f\}\}=X_{f}X_{g}(h)-\{g,\{f,h\}\}\textcolor{brown}{-X_{\{f,g\}}(h)}\]
\[\stackrel{Prop1.10}{=}[X_{f},X_{g}](h)-[X_{f},X_{g}](h)=0\]
\end{proof}

\begin{Def}
Un "algebra di Lie" è uno spazio vettoriale $\mathfrak{g}$ dotato di parentesi di Lie $[\cdot,\cdot]$, ovverosia una mappa bilineare $[\cdot,\cdot]:\mathfrak{g}\times\mathfrak{g}:\rightarrow\mathfrak{g}$ che soddisfa
\begin{itemize}
    \item [(1)]ANTISIMMETRIA:$[x,y]=-[y,x]\hspace{0.4cm}\forall x,y\in\mathfrak{g}$;
    \item [(2)]IDENTITA' DI JACOBI:$[x,[y,z]]]+[z,[x,y]]+[y,[z,x]]=0\hspace{0.4cm}\forall x,y,z\in\mathfrak{g}$;
\end{itemize}
\end{Def}

\begin{flushleft}
$\underline{Esempio}$:Lo spazio vettoriale $M(n,\mathbb{R})$ delle matrici reali $n\times n$ è un algebra di Lie $n^{2}$-dimensionale definita dal sempice commutatore
\[[A,B]:=AB-BA\]
rispetto al prodotto tra matrici. Quando penseremo ad $M(n,\mathbb{R})$ come algebra di Lie, lo denoteremo $\mathfrak{gl}(n,\mathbb{R})$.
\end{flushleft}

\begin{Def}
Un "algebra di Poisson" sul campo $\mathbb{K}$ è una terna $(\mathcal{P},\cdot,\{\})$ tale che:
\begin{itemize}
    \item [(1)]$\mathcal{P}$ sia una $\mathbb{K}$-algebra associativa rispetto al prodotto $\cdot$;
    \item [(2)]$\mathcal{P}$ sia una $\mathbb{K}$-algebra di Lie rispetto alle parentesi $\{\}$;
    \item [(3)]$\forall a,b,c\in\mathcal{P}$ sia soddisfatta la regola di Leibniz
    \[\{a,bc\}=b\{a,c\}+c\{a,b\}\]
\end{itemize}
\end{Def}

\begin{flushleft}
$\textbf{NOTA}$:Possiamo dunque concludere che, se $(M,\omega)$ è una varietà simplettica, allora $(C^{\infty}(M),\cdot,\{\})$ costituisce un algebra di Poisson.
\end{flushleft}



\subsection{Sistemi integrabili}

Vediamo ora come poter definire i sistemi integrabili, i quali ammettono un sistema di coordinate (le variabili di azione-angolo) in cui l'Hamiltoniana è più facile da studiare.

\begin{Def}
Un "sistema Hamiltoniano" è una terna $(M,\omega,h)$ dove $(M,\omega)$ è una varietà simplettica ed $h\in C^{\infty}(M)$ è una funzione, detta "Hamiltoniana" del sistema.
\end{Def}

\begin{Teo}
Dato un sistema Hamiltoniano $(M,\omega,h)$,$f\in C^{\infty}(M)$ una funzione, se $\theta_{t}:M\rightarrow M$ è il flusso associato al campo $X_{h}\in\mathcal{T}(M)$, allora
\begin{equation}
\frac{d}{dt}[f\circ\theta_{t}]=\{f,h\}\circ\theta_{t}
\end{equation}
\end{Teo}
\begin{proof}
Infatti si vede che
\[\frac{d}{dt}[f\circ\theta_{t}]=\theta_{t}^{*}(\mathcal{L}_{X_{h}}f)=(\mathcal{L}_{X_{h}}f)\circ\theta_{t}=X_{h}(f)\circ\theta_{t}\]
\[=df(X_{h})\circ\theta_{t}=\omega({X_{f}},X_{h})\circ\theta_{t}=\{f,h\}\circ\theta_{t}\]
\end{proof}

\begin{Def}
Sia $(M,\omega,h)$ un sistema Hamiltoniano, una funzione $f\in C^{\infty}(M)$ tale per cui $\{f,h\}=0$ è detta "integrale del moto" per il campo $X_{h}$.
\end{Def}

\begin{Def}
Sia M una varietà differenziabile, le funzioni $f_{1},...,f_{n}\in C^{\infty}(M)$ si dicono essere "indipendenti" se e solo se i loro differenziali $df_{1}\at{p},...,df_{n}\at{p}$ sono linearmente indipendenti in ogni punto p di M.
\end{Def}

\begin{flushleft}
\textbf{NOTA}:Se $(M,\omega,h)$ è un sistema Hamiltoniano ed $f_{1},...,f_{n}$ sono n integrali del moto che commutano tra loro, i campi vettoriali hamiltoniani associati agli integrali del moto $X_{f_{1}},...,X_{f_{n}}$ generano un sottospazio vettoriale isotropo di $T_{p}M$ in ogni punto p:
\end{flushleft}
\[\omega(X_{f_{i}},X_{f_{j}})=\{f_{i},f_{j}\}=0\]
Segue che n può essere al massimo pari a $\frac{1}{2}dim(M)$.


\begin{Def}
Un sistema Hamiltoniano $(M,\omega,h)$ si dice "(completamente) integrabile" se possiede $n=\frac{1}{2}dim(M)$ integrali del moto indipendenti,$f_{1}=h,f_{2},...,f_{n}$ i quali sono in involuzione mutua rispetto alle parentesi di Poisson ($\{f_{i},f_{j}\}=0$ per ogni $i,j=1,...,n$).
\end{Def}

\begin{flushleft}
$\underline{Esempio}$: Se $dim(M)=2$, ogni sistema Hamiltoniano $(M,\omega,h)$ è integrabile laddove $dh\neq0$. 

Si consideri ad esempio il pendolo semplice $M=T^{*}\mathbb{S}^{1}\simeq\mathbb{S}^{1}\times\mathbb{R}$ con Hamiltoniana 
\[h(\theta,p_{\theta})=\frac{p_{\theta}^{2}}{2ml^{2}}+mgl(1-\cos(\theta))\stackrel{m=l=1}{=}\frac{p_{\theta}^{2}}{2}+1-\cos\theta\]
I punti critici saranno tali che 
\[dh=0\iff\frac{\partial h}{\partial p_{\theta}}dp_{\theta}+\frac{\partial h}{\partial\theta}d\theta=p_{\theta} dp_{\theta}+\sin\theta d\theta=0\iff(p_{\theta}=0\hspace{0.2cm}e\hspace{0.2cm}\theta=0,\pm\pi)\]
\end{flushleft}
\begin{figure}
    \centering
    \includegraphics[scale=0.03]{RiF.jpg}
    \caption{Ritratto in fase del pendolo semplice}
    \label{fig:my_label}
\end{figure}

\begin{flushleft}
$\underline{Esempio}$:Si consideri ora il pendolo sferico $M=T^{*}\mathbb{S}^{2}$ con Hamiltoniana (grandezze fisiche unitarie)
\[h(\theta,\varphi,p_{\theta},p_{\varphi})=\frac{1}{2}[p_{\theta}^{2}+\frac{p_{\varphi}^{2}}{\sin^{2}\theta}]+\cos\theta\]
\[\Rightarrow dh=\frac{\partial h}{\partial\theta}d\theta+\frac{\partial h}{\partial\varphi}d\varphi+\frac{\partial h}{\partial p_{\theta}}dp_{\theta}+\frac{\partial h}{\partial p_{\varphi}}dp_{\varphi}=p_{\theta}dp_{\theta}+\frac{p_{\varphi}}{\sin^{2}\theta}dp_{\varphi}-[\frac{p_{\varphi}^{2}\cos\theta}{\sin^{3}\theta}+\sin\theta]d\theta\]
\[\Rightarrow dh=0\iff(p_{\theta}=p_{\varphi}=0\hspace{0.2cm}e\hspace{0.2cm}\theta=0,\pm\pi)\]
Gli integrali del moto sono h stessa ed $f_{2}=p_{\varphi}$. Si osservi infine che 
\[dh=(-\frac{p_{\varphi}^{2}}{\sin^{3}\theta}\cos\theta-\sin\theta,0,p_{\theta},\frac{p_{\varphi}}{\sin^{2}\theta})\]
\[df_{2}=(0,0,0,1)\]
da cui vedo che 
\[\lambda_{1}dh+\lambda_{2}df_{2}=(-\lambda_{1}[\frac{p_{\varphi}^{2}}{\sin^{3}\theta}\cos\theta+\sin\theta],0,\lambda_{1}p_{\theta},\lambda_{2}+\lambda_{1}\frac{p_{\varphi}}{\sin^{2}\theta})=(0,0,0,0)\]
\[\iff\begin{cases}
\lambda_{1}=\lambda_{2}=0\\
p_{\theta}=p_{\varphi}=0,\lambda_{2}=0,\theta=0,\pm\pi\\
p_{\varphi}=\sqrt{-\frac{\sin^{4}\theta}{\cos\theta}},\lambda_{2}=-\frac{\lambda_{1}p_{\varphi}}{\sin^{2}\theta},p_{\theta}=0
\end{cases}\]
Dunque in punti in cui i differenziali NON sono indipendenti tra loro sono  $(\theta,\varphi,p_{\theta},p_{\varphi})=\{(0,0,0,0),(\pm\pi,0,0,0)\}$ e punti del tipo $(\theta,\varphi,0,\sqrt{-\frac{\sin^{4}\theta}{\cos\theta}})$.
\end{flushleft}



\subsection{Varietà di Poisson}

\begin{Def}
Una "varietà di Poisson" $(M,\{,\})$ è una varietà differenziabile $M$ dotata di una algebra di Poisson $(C^{\infty}(M),\cdot,\{,\})$ sul campo $\mathbb{R}$.
\end{Def}

In particolare ciò significa che, se $(M,\{,\})$ è una varietà di Poisson, sono definite delle parentesi $\mathbb{R}$-bilineari
\[\{\cdot,\cdot\}:C^{\infty}(M)\times C^{\infty}(M)\rightarrow C^{\infty}(M)\]
che soddisfano le proprietà $(a),(c),(d)$ del Teorema 1.6 per qualsiasi scelta delle funzioni $f,g,h\in C^{\infty}(M)$:
\begin{itemize}
    \item (ANTISIMMETRIA):$\{f,g\}=-\{g,f\}$;
    \item (REGOLA DI LEIBNIZ):$\{f,gh\}=g\{f,h\}+h\{f,g\}$;
    \item (IDENTITA' DI JACOBI):$\{f,\{g,h\}\}+\{h,\{f,g\}\}+\{g,\{h,f\}\}=0$.
\end{itemize}

\begin{flushleft}
$\underline{Esempio}$:Sia $M=\mathbb{R}^{2n}$ con coordinate locali $(q^{1},...,q^{n},p^{1},...,p^{n})$; possiamo costruire una varietà di Poisson definendo le parentesi
\[\{f,g\}:=\sum_{i=1}^{n}\bigg(\frac{\partial f}{\partial p^{i}}\frac{\partial g}{\partial q^{i}}-\frac{\partial f}{\partial q^{i}}\frac{\partial g}{\partial p^{i}}\bigg)\]
\end{flushleft}

Sia ora $(M,\{,\})$ una varietà di Poisson, per ogni funzione $f\in C^{\infty}(M)$ possiamo interpretare $\{f,\cdot\}$ come una derivazione su $C^{\infty}(M)$ dal momento che viene soddisfatta la regola di Leibniz. Ora, dal momento che esiste una corrispondenza biunivoca tra campi vettoriali e derivazioni
\[\mathcal{T}(M)\leftrightarrow derivazioni\hspace{0.2cm}su\hspace{0.2cm}C^{\infty}(M)\]
\[X\mapsto\mathcal{L}_{X}(fg)=g\mathcal{L}_{X}(f)+f\mathcal{L}_{X}(g)\]
per ogni $f\in C^{\infty}(M)$ esiste ed è ben definito il campo vettoriale $X_{f}$ tale per cui
\[\{f,g\}=\mathcal{L}_{X_{f}}(g)=X_{f}(g)\]

Anche in questo contesto presentiamo la seguente

\begin{Def}
Sia $(M,\{,\})$ una varietà di Poisson ed $h\in C^{\infty}(M)$. L'unico campo vettoriale $X_{h}\in\mathcal{T}(M)$ che soddisfa
\[\{h,g\}=\mathcal{L}_{X_{h}}(g)=X_{h}(g)\hspace{0.5cm}\forall g\in C^{\infty}(M)\]
è detto "campo Hamiltoniano". La funzione h viene chiamata "Hamiltoniana".
\end{Def}

Come abbiamo già notato, se $(M,\omega)$ è una varietà simplettica, allora M possiede una struttura di Poisson data da $\{f,g\}:=\omega(X_{f},X_{g})$. Presentiamo ora alcuni risultati analoghi a quelli già riportati nelle sezioni precedenti
\begin{Prop}
Sia $(M,\omega)$ una varietà simplettica ed $f,g\in C^{\infty}(M)$. Allora la struttura di Poisson definita da $\{f,g\}:=\omega(X_{f},X_{g})$ implica che:
\begin{itemize}
    \item [($\star$)] $\mathcal {L}_{X_{f}}f=0\hspace{0.3cm}\forall f\in C^{\infty}(M)$;
    \item [($\star\star$)] $X_{\{f,g\}}=[X_{f},X_{g}]\hspace{0.3cm}\forall f,g\in C^{\infty}(M)$.
\end{itemize}
\end{Prop}
\begin{proof}
La prima è ovvia per antisimmetria ($\mathcal{L}_{X_{f}}f=\omega(X_{f},X_{f})=\{f,f\}$), mentre per la seconda basta sfruttare l'identità di Jacobi: siano infatti $f,g,h\in C^{\infty}(M)$ tre funzioni, allora
\[[X_{f},X_{g}](h)=X_{f}X_{g}(h)-X_{g}X_{f}(h)=\{f,\{g,h\}\}-\{g,\{f,h\}\}\]
\[=\{f,\{g,h\}\}+\{g,\{f,h\}\}=-\{h,\{f,g\}\}=\{\{f,g\},h\}=X_{\{f,g\}}(h)\]
\end{proof}

\begin{Def}
Sia $(M,\{,\})$ una varietà di Poisson, una funzione $f\in C^{\infty}(M)$ si dice essere "di Casimir" se e solo se $\{f,g\}=0\hspace{0.2cm}\forall g\in C^{\infty}(M)$.
\end{Def}

\begin{flushleft}
$\textbf{NOTA}$:Una varietà simplettica $(M,\omega)$ non possiede alcuna funzione di Casimir non banale dal moemnto che $\omega$ è non degenere.
\end{flushleft}

Vogliamo adesso riprendere in mano alcune definizioni che ci serviranno più avanti, in particolare dobbiamo riprendere il concetto di "campo k-vettoriale": ricordiamo anzitutto che con 
\[T_{k}(M):=\bigsqcup_{p\in M}\overbrace{T_{k}(T_{p}M)}^{T_{p}M\otimes...\otimes T_{p}M\hspace{0.2cm}k-volte}\]
indichiamo il fibrato dei tensori k-controvarianti su M. Se invece denotiamo
\[\Lambda_{k}(M):=\bigsqcup_{p\in M}\Lambda_{k}(T_{p}M)\subset T_{k}(M)\]
dove $\Lambda_{k}(T_{p}M)\subset T_{k}(T_{p}M)$ denota l'insieme dei tensori k-controvarianti alternanti, siamo pronti a dare la seguente

\begin{Def}
Sia M una varietà differenziabile , un "campo k-vettoriale" è una sezione del fibrato $\Lambda_{k}(M)$. L'insieme di tutti i campi k-vettoriali 
\[\pi:M\rightarrow\Lambda_{k}(M)\]
si denoterà come $\mathfrak{X}^{k}(M)\equiv\Gamma(\Lambda_{k}(M))$
In particolare, se $\xi_{i}\in T^{*}_{p}M\hspace{0.2cm}\forall i=1,...,k$ allora $\pi_{p}(\xi_{1},...,\xi_{k})\in\mathbb{R}$.
\end{Def}

Siano ora $(x^{1},...,x^{n})$ le coordinate locali di M centrate in $p\in M$, in tal caso  $\Lambda_{k}(T_{p}M)$ possiede una base $\{\frac{\partial}{\partial x^{i_{1}}}\at{p},...,\frac{\partial}{\partial x^{i_{k}}}\at{p}\}$ laddove $1\leq i_{1}<i_{2}<...<i_{k}\leq n=dim(M)$. Ma allora possiamo esprimere il generico campo k-vettoriale $\pi$ in coordinate locali come
\[\pi=\sum_{1\leq i_{1}<...<i_{k}\leq n}\pi^{i_{1},...,i_{k}}\frac{\partial}{\partial x^{i_{1}}}\wedge...\wedge\frac{\partial}{\partial x^{i_{k}}}\]
Dato un campo k-vettoriale $\pi\in\mathfrak{X}^{k}(M)$, possiamo poi definire una mappa multilineare ed antisimmetrica
\[C^{\infty}(M)\times...\times C^{\infty}(M)\rightarrow C^{\infty}(M)\]
\begin{equation}
    \{f_{1},...,f_{k}\}:=\pi(df_{1},...,df_{k})
\end{equation}
dove ricordiamo che $(\pi(df_{1},...,df_{k}))_{p}\equiv\pi_{p}(df_{1}\at{p},...,df_{k}\at{p})$.

\begin{Def}
Sia $(M,\{,\})$ una varietà di Poisson, l'unico campo bi-vettoriale $\pi\in\mathfrak{X}^{2}(M)$ definito da
\[\pi(df,dg):=\{f,g\}\iff\pi_{p}(df\at{p},dg\at{p})=\{f,g\}\at{p}\hspace{0.5cm}\forall p\in M\]
viene detto "tensore di Poisson".
\end{Def}

\begin{Def}
Il tensore di Poisson $\pi$ induce una mappa così definita
\[\pi^{\sharp}:\mathcal{A}(M)\rightarrow\mathfrak{X}(M)\]
\[df\mapsto\pi^{\sharp}(df):=\pi(df,\cdot)=X_{f}\]
In altre parole si ha $\pi^{\sharp}(df)\at{p}=\pi_{p}(df\at{p},\cdot): T^{*}_{p}M\rightarrow\mathbb{R}$ tale per cui $\pi^{\sharp}(df)\at{p}(dg\at{p})=\{f,g\}\at{p}$.
\end{Def}

Possiamo dunque esprimere il tensore di Poisson mediante un sistema di coordinate locali
\[\pi=\sum_{1\leq i<j\leq n}\pi^{ij}\frac{\partial}{\partial x^{i}}\wedge\frac{\partial}{\partial x^{j}}\]
\[\Rightarrow\pi(df,dg)=\sum_{1\leq i<j\leq n}\pi^{ij}\big(\frac{\partial f}{\partial x^{i}}\frac{\partial g}{\partial x^{j}}-\frac{\partial f}{\partial x^{j}}\frac{\partial g}{\partial x^{i}}\big)\]
Se ora riscrivo il tensore di Poisson con una somma non vincolata (ricordando di dividere poi per 2) ottengo
\[\pi=\sum_{i,j}\frac{\pi^{ij}}{2}\frac{\partial}{\partial x^{i}}\wedge\frac{\partial}{\partial x^{j}}\]
in modo che
\[\{f,g\}=\sum_{i,j}\frac{\pi^{ij}}{2}\big(\frac{\partial f}{\partial x^{i}}\frac{\partial g}{\partial x^{j}}-\frac{\partial f}{\partial x^{j}}\frac{\partial g}{\partial x^{i}}\big)\]
Questo significa che 
\[\{x^{i},x^{j}\}=\sum_{k,l}\frac{\pi^{kl}}{2}(\delta_{ki}\delta_{lj}-\delta_{li}\delta_{kj})=\frac{\pi^{ij}-\pi^{ji}}{2}=\pi^{ij}\]
In questo modo
\[\{x^{i},\{x^{j},x^{k}\}\}=\{x^{i},\pi^{jk}\}=\sum_{p,m}\frac{\pi^{pm}}{2}\big(\delta_{pi}\frac{\partial\pi^{jk}}{\partial x^{m}}-\frac{\partial\pi^{jk}}{\partial x^{p}}\delta_{mi}\big)=\sum_{m}\frac{\pi^{im}}{2}\frac{\partial\pi^{jk}}{\partial x^{m}}-\sum_{p}\frac{\pi^{pi}}{2}\frac{\partial\pi^{jk}}{\partial x^{p}}\]
\[=\pi^{il}\partial_{l}\pi^{jk}\]
dove naturalente vi è una sommatoria implicita nell'ultima espressione.
L'dentità di Jacobi può dunque essere espressa nella condizione
\begin{equation}
    \pi^{il}\partial_{l}\pi^{jk}+\pi^{kl}\partial_{l}\pi^{ij}+\pi^{jl}\partial_{l}\pi^{ki}=0
\end{equation}

\begin{flushleft}
$\textbf{NOTA}$:Si consideri il campo Hamiltoniano $X_{h}$, vogliamo esplicitarne la forma; sfruttando il fatto che
\[X_{h}=\sum_{i=1}^{n}(X_{h})_{i}\frac{\partial}{\partial x^{i}}\]
e che per definizione
\[X_{h}(g)=\{h,g\}=\pi(dh,dg)=\sum_{i<j}\pi^{ij}(\frac{\partial h}{\partial x^{i}}\frac{\partial g}{\partial x^{j}}-\frac{\partial h}{\partial x^{j}}\frac{\partial g}{\partial x^{i}})\]
\[=\sum_{i<j}\{x^{i},x^{j}\}(\frac{\partial h}{\partial x^{i}}\frac{\partial g}{\partial x^{j}}-\frac{\partial h}{\partial x^{j}}\frac{\partial g}{\partial x^{i}})=\sum_{i<j}\{x^{i},x^{j}\}\frac{\partial h}{\partial x^{i}}\frac{\partial g}{\partial x^{j}}-\sum_{i<j}\{x^{i},x^{j}\}\frac{\partial h}{\partial x^{j}}\frac{\partial g}{\partial x^{i}}\]
\[=\sum_{i<j}\{x^{i},x^{j}\}\frac{\partial h}{\partial x^{i}}\frac{\partial g}{\partial x^{j}}+\sum_{i<j}\{x^{j},x^{i}\}\frac{\partial h}{\partial x^{j}}\frac{\partial g}{\partial x^{i}}=\sum_{i<j}\{x^{i},x^{j}\}\frac{\partial h}{\partial x^{i}}\frac{\partial g}{\partial x^{j}}+\sum_{i>j}\{x^{i},x^{j}\}\frac{\partial h}{\partial x^{i}}\frac{\partial g}{\partial x^{j}}\]
\[=\sum_{i,j}\{x^{i},x^{j}\}\frac{\partial h}{\partial x^{j}}\frac{\partial g}{\partial x^{i}}\]
si trova infine la forma esplicita
\begin{equation}
(X_{h})_{i}=\sum_{j=1}^{n}\{x^{i},x^{j}\}\frac{\partial h}{\partial x_{j}}=\sum_{j=1}^{n}\pi^{ij}\frac{\partial h}{\partial x^{j}}
\end{equation}
\[X_{h}=\sum_{i,j=1}^{n}\pi^{ij}\frac{\partial h}{\partial x^{j}}\frac{\partial}{\partial x^{i}}\]
\end{flushleft}

L'identità di Jacobi (28) assume una forma più concisa se introduciamo le "parentesi di Schouten-Nijenhuis"
\[[\cdot,\cdot]_{SN}:\Gamma(\Lambda_{p}(M))\times\Gamma(\Lambda_{q}(M))\rightarrow\Gamma(\Lambda_{p+q-1}(M))\]
definite estendendo le parentesi di Lie (caso $q=p=1$) e la derivata di funzioni lungo campi vettoriali (caso $p=1$ ed $q=0$)
\[[X,Y]_{SN}=\mathcal{L}_{X}Y\hspace{0.3cm}e\hspace{0.3cm}[X,f]_{SN}=X(f)\hspace{0.3cm}\forall X,Y\in\mathcal{T}(M)=\mathfrak{X}(M),f\in C^{\infty}(M)\]
ed imponendone le proprietà
\begin{itemize}
    \item [(a)] \underline{Bilinearità}:
    \[[A,\lambda_{b}B+\lambda_{c}C]_{SN}=\lambda_{a}[A,B]_{SN}+\lambda_{c}[A,C]_{SN}\]
    \[[\lambda_{a}A+\lambda_{b}B,C]_{SN}=\lambda_{a}[A,C]_{SN}+\lambda_{b}[B,C]_{SN}\]
    \item [(b)]\underline{Antisimmetria rispetto al grado}:
    \[[A,B]_{SN}=-(-1)^{(p-1)(q-1)}[B,A]_{SN}\]
    \item [(c)]\underline{Regola di Leibniz rispetto al grado}:
    \[[A,B\wedge C]_{SN}=[A,B]_{SN}\wedge C+(-1)^{q(p+1)}B\wedge[A,C]_{SN}\]
    per ogni multi-vettore $X\in\Gamma(\Lambda_{p}(M)),B\in\Gamma(\Lambda_{q}(M))$ e $C\in\mathcal{T}(M)$ generico.
\end{itemize}

\begin{Def}
Siano $X_{1},...,X_{p},Y_{1},...,Y_{q}\in\mathcal{T}(M)$ campi vettoriali sulla varietà M, definiamo le "parentesi di Schouten-Nijenhuis"
\[[X_{1}\wedge...\wedge X_{p},Y_{1}\wedge...\wedge Y_{q}]_{SN}:=\sum_{1\leq i\leq p\atop 1\leq j\leq q}(-1)^{i+j}X_{1}\wedge...\wedge \widehat{X_{i}}\wedge...\wedge X_{p}\wedge[X_{i},Y_{j}]\wedge Y_{1}\wedge...\wedge\widehat{Y_{j}}\wedge...\wedge Y_{q}\]
dove l'elemento con il cappuccio va omesso.
\end{Def}

E' chiaro che tale definizione implichi le proprietà ricercate, viste le proprietà del prodotto wedge $\wedge$:
\begin{itemize}
    \item [(a)]Per linearità basta verificarlo sui generici $X_{1},...,X_{p},Y_{1},...,Y_{q},C\in\mathcal{T}(M)$
    \[[X_{1}\wedge...\wedge X_{p},Y_{1},...,Y_{q}]_{SN}=\sum_{1\leq i\leq p\atop 1\leq j\leq q}(-1)^{i+j}X_{1}\wedge...\wedge \widehat{X_{i}}\wedge...\wedge X_{p}\wedge[X_{i},Y_{j}]\wedge Y_{1}\wedge...\wedge\widehat{Y_{j}}\wedge...\wedge Y_{q}\]
    \[=-(-1)^{(p-1)(q-1)}\sum_{1\leq i\leq p\atop 1\leq j\leq q}(-1)^{j+i}Y_{1}\wedge...\wedge\widehat{Y_{j}}\wedge...\wedge Y_{q}\wedge[Y_{j},X_{i}]\wedge X_{1}\wedge...\wedge\widehat{X_{i}}\wedge...\wedge X_{p}\]
    dove ho scambiato (q-1)-volte il generico elemento $Y_{i}$ di (p-1)-posti ed utilizzato l'antisimmetria delle parentesi di Lie.
    \item [(b)]Sempre sui campi introdotti
    \[[X_{1}\wedge...\wedge X_{p},Y_{1}\wedge...\wedge Y_{q}\wedge C]_{SN}=\sum_{1\leq i\leq p\atop 1\leq j\leq q}(-1)^{i+j}X_{1}\wedge...\wedge \widehat{X_{i}}\wedge...\wedge X_{p}\wedge[X_{i},Y_{j}]\wedge Y_{1}\wedge...\wedge\widehat{Y_{j}}\wedge...\wedge Y_{q}\wedge C\]
    \[+\sum_{1\leq i\leq p}(-1)^{i+q+1}X_{1}\wedge...\wedge \widehat{X_{i}}\wedge...\wedge X_{p}\wedge[X_{i},C]\wedge Y_{1}\wedge...\wedge Y_{q}\]
    \[=[X_{1}\wedge...\wedge X_{p},Y_{1},...,Y_{q}]_{SN}\wedge C+\underbrace{(-1)^{qp}(-1)^{q}}_{=(-1)^{q(p+1)}}Y_{1}\wedge...\wedge Y_{q}\wedge[X_{1}\wedge...\wedge X_{p},C]_{SN}\]
    dove si sono scambiati i q elementi $Y_{i}$ di p-posti e trattenuto dentro le parentesi il $(-1)^{1+j}$.
\end{itemize}


Inoltre in questa maniera si vede facilmente che l'identità di Jacobi (28) si può riscrivere come 
\[[\pi,\pi]_{SN}=0\]
Per ulteriori dettagli rimandiamo al testo di Marsden[12].

\begin{Teo}[di Darboux-Weinstein]
Sia $(M,\{,\})$ una varietà di Poisson n-dimensionale, $p\in M$ un punto di M. Allora esiste una carta $(U,q^{1},...,q^{k},p^{1},...,p^{k},y^{1},...,y^{l})$ centrata in p tale per cui
\[\pi\at{U}=\sum_{i=1}^{k}\frac{\partial}{\partial q^{i}}\wedge\frac{\partial}{\partial p^{i}}+\sum_{i,j=1}^{l}\frac{c^{i,j}}{2}\frac{\partial}{\partial y^{i}}\wedge\frac{\partial}{\partial y^{j}}\]
laddove $c^{i,j}(p)=0\hspace{0.2cm}\forall i,j=1,...,l$ e $dim(M)=n=2k+l$.
\end{Teo}
Per la dimostrazione di questo Teorema rimandiamo al testo [13].



\section{Orbite Coaggiunte ed azioni}
Sia $(M,\omega,h)$ un sistema Hamiltoniano, dal Teorema di Liouville-Arnold sappiamo che quest'ultimo risulta integrabile se e solamente se possiede esattamente $\frac{1}{2}dim(M)$ integrali primi. In caso contrario non tutto è perduto, dal momento che anche pochi integrali primi permettono di restringere lo spazio delle fasi ed aiutarci nello studio della dinamica di tale sistema.

La ricerca di tali integrali primi tuttavia non è sempre facile, se non a volte impossibile. In questa sezione esponiamo un metodo che permette, sotto opportune condizioni, di trovare tali integrali del moto.

Prima di farlo però, dovremo studiare quelle che vengono chiamate "orbite coaggiunte", le quali richiedono alcune nozioni che esporremo qui sotto per facilitarne poi la lettura.



\subsection{Gruppi di Lie}
Questa parte riproduce parte di quanto presente in [8].

\begin{Def}
Un "gruppo di Lie" è una varietà differenziabile G che risulti anche un gruppo in senso algebrico, in modo che le mappe 
\[m:G\times G\rightarrow G\hspace{0.7cm}i:G\rightarrow G\]
\[m(g,h):=gh\hspace{0.7cm}i(g):=g^{-1}\]
siano entrambe differenziabili.
\end{Def}

\begin{flushleft}
$\underline{Esempio}$: Il gruppo delle matrici lineari ed invertibili $n\times n$ a valori reali $GL(n,\mathbb{R})$ è un gruppo di Lie, dal momento che costituisce un gruppo sotto moltiplicazione matriciale, ed è anche una sottovarietà di $M(n,\mathbb{R})$.
\end{flushleft}

\begin{Def}
Sia G un gruppo di Lie, ogni elemento $g\in G$ definisce due mappe 
\[L_{g}:G\rightarrow G\hspace{0.7cm}R_{g}:G\rightarrow G\]
\[h\mapsto gh=:L_{g}(h)\hspace{0.7cm}h\mapsto hg=:R_{g}(h)\]
dette rispettivamente "traslazione sinistra" e "traslazione destra".
\end{Def}

\begin{Def}
Siano G,H due gruppi di Lie, una mappa differenziabile $F:G\rightarrow H$ si dice essere un "omomorfismo di Lie" se è anche un omomorfismo tra gruppi. Se invece è anche un diffeomorfismo allora si dirà "isomorfismo di Lie", il che implica che la sua inversa sia a sua volta un omomorfismo di Lie.
\end{Def}

\begin{flushleft}
$\underline{Esempio}$:La mappa $exp:\mathbb{R}\rightarrow\mathbb{R}/\{0\}=:\mathbb{R}^{*}$ data da $e^{t+s}\mapsto e^{t}e^{s}$ è un omeomorfismo di Lie ($\mathbb{R}$ è considerato un gruppo di Lie sotto addizione mentre il secondo sotto moltiplicazione).
\end{flushleft}

\begin{Def}
Sia G un gruppo di Lie ed M una varietà liscia, un "azione sinistra" è la mappa 
\[G\times M\rightarrow M\]
\[(g,p)\mapsto g\cdot p\]
che soddisfa le relazioni
\[g_{1}\cdot(g_{2}\cdot p)=(g_{1}g_{2})\cdot p\hspace{0.5cm}e\cdot p=p\]
In modo analogo si definisce "l'azione destra"
\[M\times G\rightarrow M\]
\[(p,g)\mapsto p\cdot g\]
che soddisfa
\[(p\cdot g_{1})\cdot g_{2}=p\cdot(g_{1}g_{2})\hspace{0.5cm}p\cdot e=p\]
\end{Def}

Conviene assegnare un nome a tali azioni, ad esempio noi adotteremo $\rho:G\times M\rightarrow M$ in modo da poter scrivere l'azione dell'elemento $g\in G$ sull'elemento $p\in M$ come $\rho_{g}(p)$. Per definizione allora si ha
\[\rho_{g_{1}}\circ\rho_{g_{2}}(p)=\rho_{g_{1}g_{2}}(p)\hspace{0.3cm}e\hspace{0.3cm}\rho_{e}(p)=p\]
Diamo ora altre definizioni che ci serviranno in seguito
\begin{Def}
Sia $\rho:G\times M\rightarrow M$ un azione sinistra del gruppo di Lie G sulla varietà differenziabile M. Allora
\begin{itemize}
    \item [(i)]L'azione si dice essere "differenziabile" se risulta differenziabile come mappa da $G\times M$ in M, ossia se $\rho$ dipende in modo differenziabile da $(g,p)$. In tal caso $\rho_{g}:M\rightarrow M$ è un diffeomorfismo con inversa $\rho_{g^{-1}}$;
    \item [(ii)]Per ogni $p\in M$, definiamo "l'orbita di p" come
    \[G\cdot p:=\{g\cdot p:g\in G\}\]
    ossia l'insieme di tutte le immagini di p sotto l'azione dell'elemento g;
    \item [(iii)]Dato $p\in M$, "il gruppo d'isotropia" di p è l'insieme di quegli elementi di G che fissano p
    \[G_{p}:=\{g\in G:g\cdot p=p\}\]
    \item [(iv)]L'azione $\rho$ si dice poi "libera" se e solo se $G_{p}=\{e\}$.
\end{itemize}
\end{Def}

\begin{flushleft}
$\underline{Esempio}$:"L'azione naturale" di $GL(n,\mathbb{R})$ su $\mathbb{R}^{n}$ è quell'azione 
\[\rho:GL(n,\mathbb{R})\times\mathbb{R}^{n}\rightarrow\mathbb{R}^{n}\]
\[(A,x)\mapsto Ax\]
considerando $x\in\mathbb{R}^{n}$ come una matrice colonna. Essa è un azione sinistra poichè la moltiplicazione tra matrici è associativa $(AB)x=A(Bx)$, ed è liscia dal momento che $Ax$ dipende in modo polinomiale dalle entrate di A e dagli elementi di x.
\end{flushleft}

\subsection{Orbite coaggiunte}
I gruppi di Lie costituiscono un terreno molto "fertile" dal quale si vedrà sorgere in maniera naturale una struttura simplettica, associata a varietà differenziabili che vanno sotto al nome di "Orbite coaggiunte".

Sia ora N una varietà liscia, G un gruppo; e consideriamo l'azione sinistra
\[\rho:G\times M\rightarrow M\]
Grazie alla struttura differenziale associata, è ben definita la funzione
\[\rho_{*}:G\times\mathcal{T}(M)\rightarrow\mathcal{T}(M)\]
\[(g,X)\mapsto(\rho_{g})_{*}X\]
A noi interessa il caso in cui G sia un gruppo di Lie, in modo che il gruppo agisca su sè stesso mediante la traslazione sinistra $L_{g}$ definita precedentemente.

\begin{Def}
Dato un gruppo di Lie G, il campo $X\in\mathcal{T}(G)$ si dice essere un "invariante sinistro" se soddisfa $(L_{g})_{*}X=X$ per ogni elemento $g\in G$, ossia se X è $L_{g}$-correlato a sè stesso 
\begin{equation}
(L_{g})_{*}X_{x}=X_{L_{g}(x)}=X_{gx}\hspace{0.5cm}\forall g\in G
\end{equation}
L'insieme di tutti i campi vettoriali che soddisfano tale condizione viene indicato con $\mathcal{T}^{L}(G)$.
\end{Def}

\begin{Prop}
Sia G un gruppo di Lie con elemento neutro $e\in G$, allora l'applicazione $X\mapsto X(e)$ è un isomorfismo tra $\mathcal{T}^{L}(G)$ e lo spazio tangente $T_{e}G$.
\end{Prop}
\begin{proof}
Se $X\in\mathcal{T}^{L}(G)$ chiaramente
\[X_{h}=(L_{h})_{*}X_{e}\]
per ogni elemento $h\in G$, per cui X è completamente determinato dal suo valore calcolato in e. 

Viceversa, se scelgo $v\in T_{e}G$ e definisco $X\in\mathcal{T}(M)$ ponendo
\[X_{h}\equiv(L_{h})_{*}v\in T_{h}G\]
per ogni $h\in G$, otteniamo un campo invariante a sinistra che vale $v$ nell'elemento neutro. Infatti
\[(L_{g})_{*}X_{h}=(L_{g})_{*}(L_{h})_{*}v=(L_{gh})_{*}v=X_{gh}\]
e $X_{e}=(L_{e})_{*}v=(id)_{*}v=v$.
\end{proof}

\begin{Def}
Gli insiemi $\mathcal{T}^{L}(G)$ e $T_{e}G$ si chiamano entrambi "algebra di Lie" del gruppo G, e si indicano con $\mathfrak{g}$.
\end{Def}

\begin{Def}
Sia G un gruppo di Lie, un "sottogruppo ad un parametro" di G è un isomorfismo di Lie $F:\mathbb{R}\rightarrow G$.
\end{Def}

Mostreremo ora che ogni campo sinistro-invariante possiede come flusso un sottogruppo ad un parametro. Sono prima necessari due Lemma

\begin{Lem}
Si supponga che $F:M\rightarrow N$ sia una mappa differenziabile tra varietà liscie, $V\in\mathcal{T}(M),W\in\mathcal{T}(N)$ e si denotino con $\theta$ il flusso di V e con $\psi$ quello di W. Allora V e W sono F-correlati se e solo se $F\circ\theta_{t}=\psi_{t}\circ F\hspace{0.2cm}\forall t\in\mathbb{R}$, cioè se e solo se il diagramma
\[\begin{tikzcd}
M \arrow[r,"F"] \arrow[d,"\theta_{t}"] & N \arrow[d, "\psi_{t}"]\\
M \arrow[r,"F"] & N
\end{tikzcd}\]
commuta.
\end{Lem}
\begin{proof}
La commutatività del diagramma significa che, per ogni $(t,p)$ nel dominio di $\theta$, si abbia
\[\psi_{t}\circ F(p)=F\circ\theta_{t}(p)\]
Se poi denotiamo come al solito $\mathcal{D}_{p}$ il dominio di $\theta^{(p)}$, deve valere
\begin{equation}\psi^{(F(p))}(t)=F\circ\theta^{(p)}(t)\hspace{0.5cm}\forall t\in\mathcal{D}_{p}
\end{equation}
Si supponga ora che V sia F-correlato a W, se denotiamo $\gamma:\mathcal{D}_{p}\rightarrow N$ con $\gamma:=F\circ\theta^{(p)}$, si avrà
\[\gamma'(t)=(F\circ\theta^{(p)})'(t)=F_{*}(\theta^{(p)})'(t)=F_{*}V_{\theta^{(p)}(t)}\stackrel{F-corr.}{=}W_{F\circ\theta^{(p)}(t)}=W_{\gamma(t)}\]
dunque $\gamma$ è la curva integrale di W passante per $F\circ\theta^{(p)}(0)=F(p)$. Per unicità delle curve integrali, la massima curva integrale $\psi^{(F(p))}$ deve essere definita in $\mathcal{D}_{p}$, ed in particolare $\gamma(t)=\psi^{(F(p))}(t)$ in tale intervallo. Ciò prova (31).

Viceversa, si supponga che sia soddisfatta (31), allora per ogni punto $p\in M$ abbiamo
\[F_{*}V_{p}=F_{*}(\theta^{(p)})'(0)=(F\circ\theta^{(p)})'(0)=(\psi^{(F(p))})'(0)=W_{F(p)}\]
ossia V e W sono F-correlati.
\end{proof}

\begin{Lem}
Ogni campo sinistro-invariante è completo, ossia il suo flusso è definito per ogni valore di $\mathbb{R}$.
\end{Lem}
\begin{proof}
Sia $\mathfrak{g}$ l'algebra di Lie di G, e $X\in\mathfrak{g}$ a cui associamo il flusso $\theta$. Si supponga ora che la curva integrale $\theta^{(g)}$ di X sia definita in $(a,b)\subset\mathbb{R}$, e si assuma che $b<+\infty$ (il caso di $a$ si tratta in modo analogo).

Dal momento che $X$ è sinistro-invariante soddisferà (vedasi Lemma precedente)
\begin{equation}
    \theta^{(g)}\circ L_{g}=L_{g}\circ\theta^{(g)}
\end{equation}
Si osservi ora che la curva integrale passante per l'identità $\theta^{(e)}$ è definita per almeno un certo intervallo $(-\epsilon,\epsilon)$ con $\epsilon>0$. 

Si scelga ora $s\in(b-\epsilon,b)$ e definiamo la curva $\gamma:(a,s+\epsilon)\rightarrow G$ come
\[\gamma(t):=\begin{cases}
\theta^{(g)}(t)\hspace{1.8cm}t\in(a,b)\\
L_{\theta_{s}(g)}(\theta_{t-s}(e))\hspace{0.5cm}t\in(s-\epsilon,s+\epsilon)
\end{cases}\]
Dalla (32), se $s\in(a,b)$ e $|t-s|<\epsilon$ si ha
\[L_{\theta_{s}(g)}(\theta_{t-s}(e))\stackrel{(32)}{=}\theta_{t-s}(L_{\theta_{s}(g)}(e))=\theta_{t-s}(\theta_{s}(g))=\theta_{t}(g)=\theta^{(g)}(t)\]
dunque le due definizioni di gamma coincidono.

Ora, dal momento che $\gamma$ è chiaramente una curva integrale di X su $(a,b)$, per $t_{0}\in(s-\epsilon,s+\epsilon)$ sfruttiamo la sinistro-invarianza di X per ottenere
\[\gamma'(t_{0})=\frac{d}{dt}\at[\bigg]{t=t_{0}}L_{\theta_{s}(g)}(\theta_{t-s}(e))=(L_{\theta_{s}(g)})_{*}\frac{d}{dt}\at[\bigg]{t=t_{0}}\theta^{(e)}(t-s)\]
\[=(L_{\theta_{s}(g)})_{*}X_{\theta^{(e)}(t_{0}-s)}\stackrel{(30)}{=}X_{\gamma(t_{0})}\]
Ma allora $\gamma$ è una curva integrale di X su $(a,s+\epsilon)$, ma siccome $s+\epsilon>b$ questo contraddice la massimalità di $(a,b)$.
\end{proof}

\begin{Prop}
Sia G un gruppo di Lie e $X\in\mathfrak{g}$. La curva integrale di X passante per l'identità è un sottogruppo ad un parametro di G.
\end{Prop}
\begin{proof}
Sia $\theta$ il flusso di X, in modo che $\theta^{(e)}:\mathbb{R}\rightarrow G$ sia la curva integrale in questione. Chiaramente essa è differenziabile, dunque basta dimostrare che sia un omomorfismo tra gruppi, ossia che soddisfa $\theta^{(e)}(s+t)=\theta^{(e)}(s)\theta^{(e)}(t)\hspace{0.3cm}\forall t,s\in\mathbb{R}$. In effetti
\[\theta^{(e)}(s)\theta^{(e)}(t)=L_{\theta^{(e)}(s)}\overbrace{\theta^{(e)}(t)}^{=\theta_{t}(e)}\stackrel{(31)}{=}\theta_{t}(L_{\theta^{(e)}(s)}(e))=\theta_{t}(\theta^{(e)}(s))\]
\[=\theta_{t}(\theta_{s}(e))=\theta_{t+s}(e)=\theta^{(e)}(t+s)\]
\end{proof}

\begin{Teo}
Sia G un gruppo di Lie, allora esiste una corrispondenza biunivoca tra sottogrupi ad un parametro di G e curve integrali di campi vettoriali sinistro invarianti su G
\[\{sottogruppi\hspace{0.1cm}ad\hspace{0.1cm}un\hspace{0.1cm}parametro\hspace{0.1cm}di\hspace{0.1cm}G\}\longleftrightarrow\mathfrak{g}\longleftrightarrow T_{e}G\]
In particolare, ogni sottogruppo ad un parametro di G è univocamente determinato dal vettore tangente all'identità.
\end{Teo}
\begin{proof}
Partiamo con il seguente
\begin{Lem}
Siano G,H due gruppi di Lie e $\mathfrak{g},\mathfrak{h}$ le loro algebre di Lie associate. Se $F:G\rightarrow H$ è un omomorfismo di Lie, per ogni $X\in\mathfrak{g}$ esiste un unico $Y\in\mathfrak{h}$ che sia F-correlato ad X. 
\end{Lem}
\begin{proof}
Osseriviamo che, essendo F un omomorfismo di Lie si ha
\[F(g_{1}g_{2})=F(g_{1})F(g_{2})\Rightarrow F(L_{g_{1}}(g_{2}))=L_{F(g_{1})}F(g_{2})\]
\[\Rightarrow F\circ L_{g}=L_{F(g)}\circ F\Rightarrow F_{*}\circ (L_{g})_{*}=(L_{F(g)})_{*}\circ F_{*}\]
da cui ricavo
\[F_{*}X_{g}=F_{*}((L_{g})_{*}X_{e})=(L_{F(g)})_{*}F_{*}X_{e}=(L_{F(g)})_{*}Y_{\underbrace{F(e)}_{=e}}=Y_{F(g)}\]
cioè X ed Y sono F-correlati.
\end{proof}

Tornando al Teorema, sia $F:\mathbb{R}\rightarrow G$ un sottogruppo ad un parametro di G, e sia $X=F_{*}\frac{d}{dt}\in\mathfrak{g}$, dove $\frac{d}{dt}$ è pensato come un campo vettoriale sinistro-invariante su $\mathbb{R}$. Basterà dimostrare che $F(t)$ sia una curva integrale di X.

Ricordiamo che $F_{*}\frac{d}{dt}$ è definito come l'unico campo sinistro invariante in G che sia F-correlato ad $\frac{d}{dt}$, dunque per ogni istante $t_{0}\in\mathbb{R}$
\[F'(t_{0})=F_{*}\frac{d}{dt}\at[\bigg]{t_{0}}=X_{F(t_{0})}\]
cioè F è una curva integrale di X.
\end{proof}

\begin{Def}
Dato $X\in\mathfrak{g}$, chiameremo l'unico sottogruppo determinato in questo maniera come "sottogruppo ad un parametro generato da X".
\end{Def}

Facciamo ora una breve digressione ed analizziamo un po' più da vicino i flussi dei campi vettoriali sinistro-invarianti.

\begin{Def}
Sia G un gruppo di Lie, e sia $\mathfrak{g}$ la sua algebra di Lie. La "mappa esponenziale di G" è quella mappa 
\[exp:\mathfrak{g}\rightarrow G\]
\[X\mapsto exp(X)\]
costruita ponendo $exp(X):=F(1)$, laddove F è il sottogruppo ad un parametro di G generato da X.
\end{Def}

\begin{Prop}
Sia G un gruppo di Lie e $\mathfrak{g}$ la sua algebra di Lie, allora
\begin{itemize}
    \item [(a)]Per ogni $X\in\mathfrak{g}$, $F(t):=exp(tX)$ è il sottogruppo ad un parametro di G generato da X;
    \item [(b)]Per ogni $X\in\mathfrak{g}$, si ha $exp[(s+t)X]=exp(sX)exp(tX)$;
    \item [(c)]Il flusso $\theta$ di un campo sinistro-invariante X è $\theta_{t}=R_{exp(tX)}$(con R la traslazione destra definita prima).
\end{itemize}
\end{Prop}
\begin{proof}
Per dimostrare (a) basta provare che $exp(tX)=\theta_{X}^{(e)}(t)\equiv\theta_{(X),t}(e)$, o in altre parole
\[\theta^{(e)}_{tX}(1)=\theta^{(e)}_{X}(t)\]
Noi dimostreremo che
\[\theta^{(e)}_{tX}(s)=\theta^{(e)}_{X}(st)\]
che implica immediatamente quella precedente. Per farlo, si fissi $t\in\mathbb{R}$ e definiamo $\gamma:\mathbb{R}\rightarrow G$ mediante
\[\gamma(s):=\theta_{X}^{(e)}(st)\]
In tal modo, per la regola della catena
\[\gamma'(s)=t(\theta^{(e)}_{X})'(st)=tX_{\gamma(s)}\]
perciò $\gamma$ è la curva integrale di $tX$. Dal momento che $\gamma(0)=e$, per unicità delle curve integrali bisognerà avere $\gamma(s)=\theta^{(e)}_{tX}(s)$, che è quanto voluto.

Il punto (b) segue da (a), poichè quest'ultimo mostra che $t\mapsto exp(tX)$ è un omomorfismo tra gruppi.

Per (c) invece, usiamo (b) e la (32) per ottenere
\[R_{exp(tX)}(g)=g\hspace{0.1cm}exp(tX)=L_{g}(exp(tX))=L_{g}(\theta_{(X),t}(e))\]
\[\stackrel{(32)}{=}\theta_{(X),t}(e)(L_{g}(e))=\theta_{(X),t}(g)\]
\end{proof}

Ora, dal momento che $\mathfrak{g}$ può essere considerato lo spazio tangente di G nel punto $e$, ad ogni azione $\rho:G\times G\rightarrow G$ che fissi l'identità (ossia tale per cui $\rho_{g}(e)=e$ per ogni $g\in G$) possiamo associare un azione sull'algebra di Lie definendo
\[\mathfrak{p}:G\times \mathfrak{g}\rightarrow \mathfrak{g}\]
\[(g,X_{e})\mapsto(\rho_{g})_{*}X_{e} \]
In questo caso infatti si noti che $(\rho_{g})_{*}X_{e}=X_{\rho_{g}(e)}=X_{e}$ dunque tale mappa è ben definita.

La più semplice di queste azioni è la seuguente
\begin{Def}
"L'azione coniugata" è quella mappa
\[I:G\times G\rightarrow G\]
\[(g,x)\mapsto gxg^{-1}=:I_{g}(x)\]
In particolare è un azione che fissa l'identità $I_{g}(e)=e$.
\end{Def}

Sia ora $X\in\mathcal{T}(G)$ un campo vettoriale, denotiamo il flusso ad esso asscoiato con
\[\phi^{X}:(-\epsilon,\epsilon)\times G\rightarrow G \]
\[(t,g)\mapsto\phi^{X}_{t}(g)\]
Diamo ora la forma esplicita di $\mathfrak{p}$ sfruttando l'azione $I_{g}$
 che fissa l'identità
 
\begin{Def}
"L'azione aggiunta" del gruppo di Lie G è la mappa
\[Ad:G\times\mathfrak{g}\rightarrow\mathfrak{g}\]
\[(g,X_{e})\mapsto Ad_{g}(X_{e}):=(I_{g})_{*}X_{e}\]
\end{Def}

\begin{flushleft}
$\textbf{NOTA}$:In modo equivalente
\begin{equation}Ad_{g}(X_{e})=\frac{d}{dt}\big[g\phi^{X}_{t}(e)g^{-1}\big]\at[\big]{t=0}
\end{equation}
Questo perchè, ricordando che l'azione fissa l'identità, la mappa
\[\gamma:J\subset \mathbb{R}\rightarrow G\hspace{0.5cm}t\mapsto g\phi^{X}_{t}(e)g^{-1}\]
risulta essere la curva integrale di $Ad_{g}(X_{e})$,infatti $\gamma(0)=geg^{-1}=e$ e per definizione del flusso
\[X_{e}=(\phi^{X}(e))_{*}\frac{d}{dt}\at[\bigg]{t=0}\Rightarrow Ad_{g}(X_{e})=(I_{g})_{*}(\phi^{X}(e))_{*}\frac{d}{dt}\at[\bigg]{t=0}=\underbrace{(I_{g}\circ\phi^{X}(e))_{*}}_{=(g\phi^{X}(e)g^{-1})_{*}}\frac{d}{dt}\at[\bigg]{t=0}\]
\end{flushleft}

Si consideri ora un azione $\rho:G\times M\rightarrow M$ sulla varietà differenziabile M (G gruppo di Lie),$X\in\mathfrak{g}=\mathcal{T}^{L}(G)$ un elemento dell'algebra di Lie, allora il flusso passante per l'identità $\phi^{X}(e):(-\epsilon,\epsilon)\rightarrow G$ induce un gruppo ad un parametro di diffeomorfismi su M
\[\phi^{X_{M}}:(-\epsilon.\epsilon)\times M\rightarrow M\]
\[(t,p)\mapsto\rho_{\phi^{X}_{t}(e)}(p)\]
La sua derivata definisce così un campo vettoriale $X_{M}\in\mathcal{T}(M)$ tale per cui
\begin{equation}X_{M}\at{p}=\frac{d}{dt}\at[\bigg]{t=0}[\rho_{\phi^{X}_{t}(e)}(p)]\hspace{0.5cm}\forall p\in M
\end{equation}
In altre parole il campo $X_{M}$ possiede la curva integrale
\[t\mapsto\rho_{\phi^{X}_{t}(e)}(p)\]
dove si osservi che $\phi^{X_{M}}(0,p)=\rho_{e}(p)=p$.
Tutto ciò può essere riassunto nella seguente

\begin{Def}
"L'azione infinitesima" di $\mathfrak{g}$ su M è la mappa
\[\mathfrak{g}\rightarrow\mathcal{T}(M)\]
\[X\mapsto X_{M}\]
\end{Def}

Il significato di tale definizione è che la trasformazione infinitesima generata da $X$ su $\mathfrak{g}$ genera una trasformazione infinitesima
su M descritta da $X_{M}$.

Considerando ora $\mathfrak{g}$ anzichè M ottengo la seguente

\begin{Def}
L'azione infinitesima associata associata all'azione aggiunta $Ad:G\times\mathfrak{g}\rightarrow\mathfrak{g}$ è la mappa
\[ad:\mathfrak{g}\rightarrow\mathcal{T}(\mathfrak{g})\]
\[X\mapsto X_{\mathfrak{g}}=:ad_{X}\]
\end{Def}

\begin{Prop}
L'azione infinitesima precedente soddisfa
\begin{equation}
    ad_{X}\at{Y}=[X,Y]\hspace{0.5cm}\forall X,Y\in\mathfrak{g}
\end{equation}
\end{Prop}
\begin{proof}
Si scelga una qualsivoglia funzione $f\in C^{\infty}(G)$ e si consideri $ad_{X}\at{Y}$ come un vettore di $T_{e}G$, in questo modo
\[ad_{X}\at{Y}(f)\stackrel{(33)}{=}\frac{d}{dt}\at[\bigg]{t=0}[Ad_{\phi^{X}_{t}(e)}Y](f)\stackrel{(32)}{=}\frac{\partial^{2}}{\partial t\partial s}\at[\bigg]{t=s=0}f[\phi^{X}_{t}(e)\phi^{Y}_{s}(e)\phi^{X}_{-t}(e)]\]
\[\stackrel{Leibniz}{=}\frac{\partial^{2}}{\partial t_{1}\partial s}\at[\bigg]{t_{1}=t_{2}=s=0}f[\phi^{X}_{t_{1}}(e)\phi^{Y}_{s}(e)\phi^{X}_{-t_{2}}(e)]+\frac{\partial^{2}}{\partial t_{2}\partial s}\at[\bigg]{t_{1}=t_{2}=s=0}f[\phi^{X}_{t_{1}}(e)\phi^{Y}_{s}(e)\phi^{X}_{-t_{2}}(e)]\]
ma, ricordando che il flusso passa per l'identità, ottengo
\[=\frac{\partial^{2}}{\partial t\partial s}\at[\bigg]{t=s=0}f[\phi^{X}_{t}(e)\phi^{Y}_{s}(e)]+\frac{\partial^{2}}{\partial t\partial s}\at[\bigg]{t=s=0}f[\phi^{Y}_{s}(e)\phi^{X}_{\textcolor{red}{-t}}(e)]\]
\[=\frac{\partial^{2}}{\partial t\partial s}\at[\bigg]{t=s=0}f[\phi^{X}_{t}(e)\phi^{Y}_{s}(e)]\textcolor{red}{-}\frac{\partial^{2}}{\partial t\partial s}\at[\bigg]{t=s=0}f[\phi^{Y}_{s}(e)\phi^{X}_{\textcolor{red}{t}}(e)]=X(Y(f))-Y(X(f))\]
\[=[X,Y](f)\]
dove abbiamo sfruttato il fatto che (Prop. 2.3)
\[\phi_{t}^{X}(e)=R_{exp(tX)}(e)=exp(tX)\hspace{0.5cm}\phi^{Y}_{s}(e)=R_{exp(sY)}(e)=exp(sY)\]
dunque ad esempio
\[X(Y(f))=X(\frac{d}{ds}\at[\bigg]{s=0}f(R_{exp(sY)}(e)))=\frac{\partial^{2}}{\partial s\partial t}\at[\bigg]{t=s=0}f(R_{exp(sY)}(R_{exp(tX)}(e)))\]
\[=\frac{\partial^{2}}{\partial s\partial t}\at[\bigg]{t=s=0}f(exp(tX)exp(sY))\]
\end{proof}

\begin{flushleft}
$\textbf{NOTA}$:Si osservi che $ad_{X}\at{Y}\in T_{Y}\mathfrak{g}$, ma dal momento che $\mathfrak{g}=T_{e}G$ è uno spazio vettoriale il suo spazio tangente è isomorfo a sè stesso. Perciò possiamo considerare il commutatore $[X,Y]$ sia come un elemento di $\mathfrak{g}$ (sia X che Y sono invarianti sinistri) sia come un vettore.
\end{flushleft}

\begin{Def}
"L'azione coaggiunta" è la mappa duale di Ad, ossia 
\[Ad^{*}:G\times\mathfrak{g}^{*}\rightarrow\mathfrak{g}^{*}\]
\[(g,\alpha)\mapsto Ad^{*}_{g}\alpha\]
tale per cui
\[Ad^{*}_{g}\alpha(X):=\alpha(Ad_{g^{-1}}X)\]
Invece "l'azione infinitesima coaggiunta" è quella mappa
\[ad^{*}:\mathfrak{g}\rightarrow\mathcal{T}(\mathfrak{g}^{*})\]
\[X\mapsto ad^{*}_{X}\]
che soddisfa
\[ad^{*}_{X}\at{\alpha}(Y)\stackrel{(34)}{=}\frac{d}{dt}\at[\bigg]{t=0}(Ad^{*}_{exp(tX_{e})}\alpha)(Y)=\frac{d}{dt}\at[\bigg]{t=0}\alpha(Ad_{exp(-tX_{e})}Y)\stackrel{(\clubsuit)}{=}\alpha(ad_{-X}\at{Y})\]
\[=-\alpha([X,Y])\]
dove in $(\clubsuit)$ si è sfruttata la linearità di $\alpha$ e sempre la (34).
\end{Def}

\begin{flushleft}
$\textbf{NOTA}$:Per definizione vale ovviamente
\[(ad^{*}_{X}\at{\alpha}+\lambda ad^{*}_{Y}\at{\alpha})(Z)=ad^{*}_{X}\at{\alpha}(Z)+\lambda ad^{*}_{Y}\at{\alpha}(Z)=-\alpha([X,Z])-\lambda\alpha([Y,Z])\]
\[=-\alpha([X+\lambda Y,Z])\Rightarrow ad^{*}_{X}\at{\alpha}+\lambda ad^{*}_{Y}\at{\alpha}=ad^{*}_{X+\lambda Y}\at{\alpha} \]
\end{flushleft}

\begin{Def}
Una mappa simmetrica bilineare non degenere $(\cdot,\cdot):\mathfrak{g}\times\mathfrak{g}\rightarrow\mathbb{R}$ viene detta "Ad-invariante" se e solo se soddisfa
\begin{equation}
    (Ad_{g}X,Ad_{g}Y)=(X,Y)\hspace{0.5cm}\forall g\in G,\forall X,Y\in\mathfrak{g}
\end{equation}
\end{Def}

\begin{Prop}
Data una mappa bilineare Ad-invariante $(\cdot,\cdot)$, vi è un isomorfismo \[\mathfrak{g}\simeq\mathfrak{g}^{*}\hspace{0.6cm}X\mapsto (X,\cdot)\]
che manda azioni aggiunte in azioni coaggiunte $Ad_{g}(X)\mapsto Ad_{g}^{*}((X,\cdot))$
\[\begin{tikzcd}
\mathfrak{g} \arrow[r] \arrow[d,"Ad_{g}"] & \mathfrak{g}^{*} \arrow[d, "Ad^{*}_{g}"]\\
\mathfrak{g} \arrow[r] & \mathfrak{g}^{*}
\end{tikzcd}\]
\end{Prop}
\begin{proof}
In effeti vale
\[(Ad_{g}X,Y)=(Ad_{g}X,Ad_{g}(Ad_{g^{-1}}Y))\stackrel{(36)}{=}(X,Ad_{g^{-1}}Y)=(Ad^{*}_{g}(X,\cdot))(Y)\equiv <Ad^{*}_{g}(X,\cdot),Y>\]
dove abbiamo sfruttato la definizione di $Ad^{*}_{g}$ sostituendo $(X,\cdot)$ ad $\alpha$.
\end{proof}

Si considerino ora una varietà differenziabile M, $U\subset M$ un aperto. Sia  $\iota:U\hookrightarrow\mathfrak{g}^{*}$ l'embedding dello spazio delle fasi nell'algebra di Lie duale. In tal caso ogni funzionale $F:U\rightarrow\mathbb{R}$ può essere esteso ad una funzione lisica su $\mathfrak{g}^{*}$.

\begin{Def}
Sia $C^{\infty}(\mathfrak{g}^{*})$ lo spazio di tutte le funzioni liscie su $\mathfrak{g}^{*}$ del tipo $f:=F\circ\iota^{-1}:\mathfrak{g}^{*}\rightarrow\mathbb{R}$ con $F\in C^{\infty}(U)$. Il "differenziale" $df$ può calcolato nel punto $\eta\in\mathfrak{g}^{*}$ definendo
\begin{equation}
    <\xi,df\at{\eta}>\equiv\xi(df\at{\eta}):=\frac{d}{dt}\at[\bigg]{t=0}f(\eta+t\xi)\hspace{0.5cm}\forall\xi\in\mathfrak{g}^{*}
\end{equation}
\end{Def}



\subsection{Struttura simplettica delle orbite coaggiunte}
Con riferimento agli oggetti introdotti nella precedente sezione, partiamo con la seguente

\begin{Def}
"L'orbita coaggiunta" passante per $\alpha\in\mathfrak{g}^{*}$ è
\[\Omega_{\alpha}=\{Ad^{*}_{g}\alpha:g\in G\}\subseteq\mathfrak{g}^{*}\]
\end{Def}

\begin{Teo}
Vi è un isomorfismo $T_{\alpha}\Omega_{\alpha}\simeq\mathfrak{g}^{*}/Ker(ad^{*}_{\cdot}\at{\alpha})$, dove il nucleo si riferisce all'applicazione indotta naturalmente dalla Definizione 2.15.
\end{Teo}
\begin{proof}
Iniziamo con il cercare di capire a livello intuitivo come sia fatto il nucleo di $ad^{*}_{\cdot}\alpha$. Ricordiamo che $ad^{*}$ rappresenta l'azione di G su $\mathfrak{g}^{*}$, per cui possiamo pensare al suo nucleo (nel punto $\alpha$) come lo stabilizzatore di tale azione, ossia l'insieme di quei campi vettoriali che non "muovono" $\alpha$. Possiamo allora facilmente raffigurare l'isomorfismo cercato:lo spazio tangente all'orbita coaggiunta è costituita dai vettori lungo i quali l'azione $\alpha$ può muoversi, dunque ha senso dire che coloro sono esattamente tutti i vettori dell' azione infinitesima a meno di quelli che non la muovono.

Si consideri la sommersione $G\rightarrow\Omega_{\alpha}$,$g\mapsto Ad^{*}_{g}\alpha$, allora dato $X\in\mathfrak{g}$, possiamo definire la nuova applicazione
\[X\mapsto \frac{d}{dt}\at[\bigg]{t=0}Ad^{*}_{\phi^{X}_{t}(e)}\alpha=ad^{*}_{X}\at{\alpha}\]
o meglio
\[ad^{*}_{\cdot}\at{\alpha}:\mathfrak{g}\rightarrow T_{\alpha}\Omega_{\alpha}\]
Grazie alla suriettività, l'asserzione segue dal primo Teorema d'isomorfismo.
\end{proof}

\begin{flushleft}
$\textbf{NOTA}$:Si supponga di avere una funzione $f:\mathfrak{g}^{*}\rightarrow\mathbb{R}$ del tipo $f(\alpha)=\alpha(X)$ per qualche elemento fissato $X\in\mathfrak{g}$, allora è facile capire come agisce un campo vettoriale su di essa: se $ad^{*}_{Y}\in\mathcal{T}(\mathfrak{g}^{*})$
\begin{equation}
    ad^{*}_{Y}\at{\alpha}(f)=ad^{*}_{Y}\at{\alpha}(X)
\end{equation}
\end{flushleft}

Siamo pronti ora a fornire una struttura simplettica alle orbite coaggiunte.

\begin{Teo}[Kirllov-Kostant-Souriau]
La 2-forma $\omega\in\mathcal{A}^{2}(\Omega_{\beta})$ definita da
\begin{equation}
    \omega_{\alpha}(ad^{*}_{X}\at{\alpha},ad^{*}_{Y}\at{\alpha}):=\alpha([X,Y])\hspace{0.5cm}\forall X,Y\in\mathfrak{g},\forall\alpha\in\Omega_{\beta}
\end{equation}
è simplettica.
\end{Teo}
\begin{proof}
Cominciamo con l'osservare che tale forma è ben definita: antisimmetria e bilinearità sono ovvie per le proprietà del commutatore
\[\omega_{\alpha}(ad^{*}_{X}\at{\alpha}+\lambda ad^{*}_{Y}\at{\alpha},ad^{*}_{Z}\at{\alpha})=\omega_{\alpha}(ad^{*}_{X+\lambda Y}\at{\alpha},ad^{*}_{Z}\at{\alpha})=\alpha([X+\lambda Y,Z])\]
\[=\alpha([X,Z]+\lambda[Y,Z])=\alpha([X,Z])
+\lambda\alpha([Y.Z])=\omega_{\alpha}(ad^{*}_{X}\at{\alpha},ad^{*}_{Z}\at{\alpha})+\lambda\omega_{\alpha}(ad^{*}_{Y}\at{\alpha},ad^{*}_{Z}\at{\alpha})\]
Dimostriamo adesso la sua non degeneratezza: dal momento che $ad^{*}_{\cdot}$ è suriettiva $\forall\alpha\in\mathfrak{g}^{*}$, dobbiamo provare che
\[\omega_{\alpha}(ad^{*}_{X}\at{\alpha},ad^{*}_{Y}\at{\alpha})=0\hspace{0.4cm}\forall Y\in\mathfrak{g}\iff ad^{*}_{X}\at{\alpha}=0\]
che, grazie alla Definizione 2.15, equivale a
\[\omega_{\alpha}(ad^{*}_{X}\at{\alpha},ad^{*}_{Y}\at{\alpha})=\alpha([X,Y])=-ad^{*}_{X}\at{\alpha}(Y)=0\hspace{0.4cm}\forall Y\in\mathfrak{g}\iff ad^{*}_{X}\at{\alpha}=0\]
che è ovviamente vera.
Infine resta da provare che $\omega$ sia chiusa: usiamo ora la Proposizione A.2 dell'appendice per ottenere
\[((\mathcal{L}_{ad^{*}_{Z}}\omega)(ad^{*}_{X},ad^{*}_{Y}))_{\alpha}=ad^{*}_{Z}\at{\alpha}(\omega(ad^{*}_{X},ad^{*}_{Y}))-\omega_{\alpha}([ad^{*}_{Z}\at{\alpha},ad^{*}_{X}\at{\alpha}],ad^{*}_{Y}\at{\alpha})\]
\[-\omega_{\alpha}(ad^{*}_{X}\at{\alpha},[ad^{*}_{Z}\at{\alpha},ad^{*}_{Y}\at{\alpha}])\stackrel{(38)}{=}ad^{*}_{Z}\at{\alpha}([X,Y])-\omega_{\alpha}(ad^{*}_{[Z,X]}\at{\alpha},ad^{*}_{Y}\at{\alpha})\]
\[-\omega_{\alpha}(ad^{*}_{X}\at{\alpha},ad^{*}_{[Z,Y]}\at{\alpha})=ad^{*}_{Z}\at{\alpha}([X,Y])-\alpha([[Z,X],Y])-\alpha([X,[Z,Y]])\]
\[=ad^{*}_{Z}\at{\alpha}([X,Y])+\alpha([Y,[Z,X]]+[X,[Y,Z]])\]
D'altra parte però
\[ad^{*}_{Z}\at{\alpha}(X)=\frac{d}{dt}\at[\bigg]{t=0}(Ad^{*}_{\phi^{Z}_{t}(e)}\alpha)(X)=\frac{d}{dt}\at[\bigg]{t=0}(\alpha(Ad_{\phi^{Z}_{-t}(e)}X))=-\alpha(ad_{Z}\at{X})\]
\[=-\alpha([Z,X])\Rightarrow ad^{*}_{Z}\at{\alpha}([X,Y])=-\alpha([Z,[X,Y]])=\alpha(Z,[Y,X])\]
Unendo i due risultati e sfruttando l'identità di Jacobi vedo che $\mathcal{L}_{ad^{*}_{Z}}\omega=0$ per ogni $Z\in\mathfrak{g}$.

A questo punto però si osservi anche che 
\[ad^{*}_{Z}\lrcorner\omega=\omega(ad^{*}_{Z},\cdot)=-Z(\cdot)\hspace{0.3cm}con\hspace{0.3cm}ad^{*}_{X}\mapsto ad^{*}_{X}(Z)=:Z(ad^{*}_{X})\]
dal momento che
\[\textcolor{blue}{ad^{*}_{Z}\at{\alpha}\lrcorner\omega_{\alpha}}(ad^{*}_{X}\at{\alpha})=\omega\at{\alpha}(ad^{*}_{X}\at{\alpha},ad^{*}_{Z}\at{\alpha})=-\alpha([Z,X])\stackrel{(35)}{=}\alpha(ad_{X}\at{Z})=-ad^{*}_{X}\at{\alpha}(Z)=\textcolor{blue}{-Z}(ad^{*}_{X}\at{\alpha})\]
Ora, se consideriamo Z come una mappa
\[Z:\mathfrak{g}^{*}\rightarrow\mathbb{R}\]
\[\beta\mapsto Z(\beta)=\beta_{i}(Z^{i})\equiv dZ^{i}(\beta_{i})\]
cioè $Z\in C^{\infty}(\mathfrak{g}^{*})$, allora $dZ$ sarà del tipo (con $\beta_{i}$ coordinate rispetto a $\mathfrak{g}^{*}$)
\[dZ=\frac{\partial Z}{\partial\beta^{i}}d\beta^{i}=Z^{i}d\beta_{i}=Z(\cdot)\]
Segue che 
\[ad^{*}_{Z}\lrcorner\omega\at{\alpha}=-dZ\]
\[\Rightarrow 0= (\mathcal{L}_{ad^{*}_{Z}}\omega)_{\alpha}\stackrel{Cartan}{=}d(ad^{*}_{Z}\at{\alpha}\lrcorner\omega_{\alpha})+ad^{*}_{Z}\at{\alpha}\lrcorner d\omega_{\alpha}=ad^{*}_{Z}\at{\alpha}\lrcorner d\omega_{\alpha}\hspace{0.5cm}\forall ad^{*}_{Z}\at{\alpha}\]
da cui concludiamo che $d\omega=0$ dal momento che $ad^{*}_{Z}\at{\alpha}$ genera $T_{\alpha}\Omega_{\beta}$.
\end{proof}

\begin{flushleft}
$\underline{Esempio}$:Sia $G=SO(3)\equiv\{g\in M_{3\times 3}(\mathbb{R}):det(g)=1,gg^{T}=e\}$ identiticato con la varietà differenziabile associata alle matrici quadrate ortogonali e con determinante positivo.

Possiamo ottenere una descrizione esplicita della sua algebra di Lie $T_{\mathbb{1}_{3}}SO(3)$ differenziando la relazione $\mathbb{1}_{3}=XX^{T}$:
\[\mathfrak{so}(3)=\{\frac{d}{dt}\at[\bigg]{t=0}X:X\in\mathbb{R}^{3\times 3},X(0)=\mathbb{1}_{3},XX^{T}=\mathbb{1}_{3}\}\]
\[=\{X\in\mathbb{R}^{3\times 3}:X+X^{T}=\mathbb{0}\}\]
poichè 
\[\mathbb{0}=\frac{d}{dt}\at[\bigg]{t=0}(XX^{T})=\dot{X}_{0}X^{T}_{0}+X_{0}\dot{X}^{T}_{0}\iff \dot{X}_{0}+\dot{X}_{0}^{T}=\mathbb{0}\]
Le parentesi di Lie su tale algebra sono il solito commutatore tra matrici. C'è poi un isomorfismo
\[\phi:\mathbb{R}^{3}\rightarrow\mathfrak{so}(3)\]
\[(x,y,z)\mapsto\begin{pmatrix}
0  & -z & y\\
z  & 0  & -x\\
-y & x & 0
\end{pmatrix}\]
che manda il prodotto vettoriale nel commutatore
\[\phi(v\times w)=[\phi(v),\phi(w)]\]
perciò
\[(\mathbb{R}^{3},\times)\cong(\mathfrak{so}(3),[\cdot,\cdot])\]
Si osservi anche che per definizione abbiamo
\[\phi(v)u\equiv v\times u\]
e questo significa che, siccome gli elementi di $\mathfrak{s0}(3)$ non modificano il prodotto vettore $gv\times gw=g(v\times w)$, si ha
\[Ad_{g}(\phi(v))gw=g\phi(v)g^{-1}gw=g\phi(v)w=g(v\times w)=gv\times gw=\phi(gv)gw\]
\[\Rightarrow Ad_{g}\phi(v)=\phi(g\cdot v)\]
In altre parole l'azione aggiunta $Ad_{g}$ applicata agli elementi di $\mathfrak{so}(3)$ risulta in una rotazione dei vettori
\[Ad_{g}(\phi(v))\stackrel{(33)}{=}g\phi(v)\overbrace{g^{-1}}^{=g^{T}}=\phi(g\cdot v)\hspace{0.5cm}\forall g\in SO(3),v\in\mathbb{R}^{3}\]
Possiamo poi definire la seguente applicazione bilineare Ad-invariante
\[(\phi(v),\phi(w)):=v\cdot w\hspace{0.5cm}\forall v,w\in\mathbb{R}^{3}\]
che porta alla seguente catena di isomorfismi (vedasi Prop. 2.5)
\[
\begin{tikzcd}
\mathbb{R}^{3}\arrow[r,"\sim"] & \mathfrak{so}(3) \arrow[r,"\sim"] & \mathfrak{so}^{*}(3)
\end{tikzcd}
\]
e manda Ad in $Ad^{*}$ come nella dimostrazione della Proposizione appena citata
\[Ad^{*}_{g}((\phi(v),\cdot))=(Ad_{g}(\phi(v)),\cdot)=(\phi(g\cdot v),\cdot)\]
Si osservi in particolare che
\[(Ad_{g}\phi(v),Ad_{g}(w))=(g\cdot v,g\cdot w)=(gv)\cdot(gw)=v\cdot g^{T}gw=v\cdot w\]
Si considerino ora le orbite coaggiunte
\[\Omega_{(\phi(v),\cdot)}=\{Ad^{*}_{g}((\phi(v),\cdot)):g\in SO(3)\}=\{(\phi(g\cdot v),\cdot):g\in SO(3)\}\]
La forma simplettica su di esse definita è 
\[(\omega(ad^{*}_{\phi(v)},ad^{*}_{\phi(w)}))_{(\phi(u),\cdot)}\stackrel{(37)}{=}(\phi(u),[\phi(v),\phi(w)])=(\phi(u),\phi(v\times w))=u\cdot (v\times w)\]
Lo spazio $T_{v}\Omega_{v}$ descrive le azioni infinitesime (rotazioni infinitesime in questo caso) a partire da $v$. Possiamo dunque riassumere 
\begin{itemize}
    \item [(•)]$\mathfrak{so}(3)$:rotazioni infinitesime della sfera;
    \item [(••)]$ker(ad^{*}\at{v})$:rotazioni infinitesime attorno a $v$, ossia azioni che non modificano $(\phi(v),\cdot)$. In particolare 
    \[Ad^{*}_{g}(\phi(v),\cdot)=(\phi(v),\cdot)\iff(\phi(g\cdot v),\cdot)=(\phi(v),\cdot)\iff g\cdot v=v\]
    Dunque il nucleo è costituito effettivamnte dalle rotazioni attorno a $v$;
    \item [(•••)]$T_{v}\Omega_{v}$:movimenti infinitesimi nella sfera che partono da $v$.
\end{itemize}
\end{flushleft}



\section{$R$-Matrice}

\begin{Def}
Sia $\mathfrak{g}$ un algebra di Lie, una "R-matrice" è una mappa $R:\mathfrak{g}\rightarrow\mathfrak{g}$ con la proprietà che le "parentesi modificate"
\begin{equation}
[X,Y]_{R}:=[RX,Y]+[X,RY]\hspace{0.5cm}\forall X,Y\in\mathfrak{g}
\end{equation}
definiscono a loro volta un algebra di Lie su $\mathfrak{g}$, che chiameremo $(\mathfrak{g},[\cdot,\cdot]_{R})$.
\end{Def}

\begin{Prop}
Sia $(\mathfrak{g},[\cdot,\cdot])$ un algebra di Lie, la mappa lineare $R:\mathfrak{g}\rightarrow\mathfrak{g}$ è una R-matrice se e solo se
\begin{equation}
[B_{R}(X,Y),Z]+[B_{R}(Y,Z),X]+[B_{R}(Z,X),Y]=0
\end{equation}
laddove 
\[B_{R}(X,Y)=[RX,RY]-R[X,Y]_{R}\]
\end{Prop}
\begin{proof}
Vogliamo dimostrare che valga l'identità di Jacobi rispetto alle parentesi modificate (l'antisimmetria è ovviamente sempre rispettata da $[,]_{R}$)
\[[X,[Y,Z]_{R}]_{R}+[Z,[X,Y]_{R}]_{R}+[Y,[Z,X]_{R}]_{R}=0\]
ovvero che
\[[RX,[Y,Z]_{R}]+[X,R[Y,Z]_{R}]+[RZ,[X,Y]_{R}]+[Z,R[X,Y]_{R}]\]
\[+[RY,[Z,X]_{R}]+[Y,R[Z,X]_{R}]=0\]
\[\iff [RX,[RY,Z]+[Y,RZ]]+[X,R[Y,Z]_{R}]+[RZ,[RX,Y]+[\textcolor{red}{X},RY]]+[Z,R[X,Y]_{R}]\]
\[+[RY,[RZ,\textcolor{red}{X}]+[Z,RX]]+[Y,R[Z,X]_{R}]=0\]
\[\iff [RZ,[\textcolor{red}{X},RY]]+[RY,[RZ,\textcolor{red}{X}]]+[X,R[Y,Z]_{R}]+p.c.=0\]
\[\stackrel{Jacobi}{\iff}-[X,[RY,RZ]]+[X,R[Y,Z]_{R}]+p.c.=0\]
\[\iff [B_{R}(Y,Z),X]+p.c.=0\]
dove p.c. indica le varie permutazioni cicliche degli elementi, omessi per brevità. Abbiamo ritrovato la (40).
\end{proof}

Si consideri ora una mappa bilineare $Ad-$invariante $(\cdot,\cdot):\mathfrak{g}\times\mathfrak{g}\rightarrow\mathbb{R}$, denoteremo l'isomorfismo definito nella Proposizione 2.5 come
\[\chi:\mathfrak{g}\rightarrow\mathfrak{g}^{*}\]
\[X\mapsto(X,\cdot)\]
In tal modo 
\[(X,Y)=\chi(X)(Y)\]

\begin{Def}
Si consideri un algebra di Lie $\mathfrak{g}$, se $[\cdot,\cdot]_{R}$ sono le parentesi modificate associate ad una R-matrice, definiamo le parentesi di Poisson indotte su $\mathfrak{g}^{*}$ come
\[\{\cdot,\cdot\}_{R}:C^{\infty}(\mathfrak{g}^{*})\times C^{\infty}(\mathfrak{g}^{*})\rightarrow C^{\infty}(\mathfrak{g}^{*})\]
\[(f,g)\mapsto \{f,g\}_{R}\]
tali per cui
\[\{f,g\}_{R}\at{\xi}:=\xi([dg\at{\xi},df\at{\xi}]_{R})\hspace{0.5cm}\forall f,g\in C^{\infty}(\mathfrak{g}^{*}),\xi\in\mathfrak{g}^{*}\]
\end{Def}

Si noti infatti che $df\at{\xi}\in T^{*}_{\xi}\mathfrak{g}^{*}=(\mathfrak{g}^{*})^{*}\simeq\mathfrak{g}$, dunque tali parentesi sono ben definite. Inoltre si possono costruire ovviamente per qualsiasi parentesi di Lie $[,]$. 

\begin{Cor}
Condizione sufficiente affinchè la mappa $R:\mathfrak{g}\rightarrow\mathfrak{g}$ sia una R-matrice è che soddisfi "l'equazione modificata di Yang-Baxter"
\begin{equation}
    B_{R}(X,Y)=-c[X,Y]\hspace{0.5cm}\forall X,Y\in\mathfrak{g}
\end{equation}
per qualche costante $c\in\mathbb{R}$.
\end{Cor}
\begin{proof}
E' ovvio poichè se vale la (41) allora (40) diventa
\[-c([[X,Y],Z]+[[Y,Z],X]+[[Z,X],Y])=0\]
per l'identità di Jacobi.
\end{proof}

Chiaramente vi sono due possibilità (a costo di riscalare tutto):$c=0$ nel quale R è un automorfismo dell'algebra oppure $c=1$. Nell'ultimo caso c'è una classe speciale di soluzioni che emerge in modo naturale.

\begin{flushleft}
$\underline{Esempio}$:Si consideri $(\mathfrak{g},[\cdot,\cdot])$ che si decomponga (nel senso di spazio vettoriale) in somma diretta $\mathfrak{g}=\mathfrak{g}_{+}\oplus\mathfrak{g}_{-}$. NON richiediamo che $[\mathfrak{g}_{+},\mathfrak{g}_{+}]=0$; chiamiamo ora con $pr_{\pm}$ le proiezioni sui due sottospazi $\mathfrak{g}_{\pm}$ e denotiamo $X_{\pm}:=pr_{\pm}X$ per ogni vettore $X\in\mathfrak{g}$.

Dimostriamo che $R:=\frac{1}{2}(pr_{+}-pr_{-})$ è una R-matrice:
\[[X,Y]_{R}=[RX,Y]+[X,RY]=\frac{1}{2}[X_{+}-X_{-},Y]+\frac{1}{2}[X,Y_{+}-Y_{-}]\]
\[\frac{1}{2}[X_{+}-X_{-},Y_{+}+Y_{-}]+\frac{1}{2}[X_{+}+X_{-},Y_{+}-Y_{-}]\]
\[=\frac{1}{2}([X_{+},Y_{+}+Y_{-}]-[X_{-},Y_{+}+Y_{-}]+[X_{+},Y_{+}-Y_{-}]+[X_{-},Y_{+}-Y_{-}])\]
\[=[X_{+},Y_{+}]-[X_{-},Y_{-}]\]
Vediamo allora che è possibile decomporre l'identità di Jacobi per $[\cdot,\cdot]_{R}$ in due identità di Jacobi per $[\cdot,\cdot]$, le quali sono vere per ipotesi. Abbiamo dimostrato quanto volevamo.
\end{flushleft}

\begin{Prop}
Se la funzione $h\in C^{\infty}(\mathfrak{g}^{*})$ è di Casimir per la varietà di Poisson $(\mathfrak{g}^{*},\{,\})$ allora il campo vettoriale Hamiltoniano ad essa associato rispetto alle parentesi di Poisson modificate $\{,\}_{R}$ è
\begin{equation}
    X^{R}_{h}\at{\alpha}=ad^{*}_{R(dh\at{\alpha})}\at{\alpha}
\end{equation}
\end{Prop}
\begin{proof}
Per ogni $f\in C^{\infty}(\mathfrak{g}^{*})$ e $\alpha\in\mathfrak{g}$ si ha che (vedasi Definizione 1.28 per campi Hamiltoniani su varietà di Poisson)
\[X^{R}_{h}\at{\alpha}(f)=\{h,f\}_{R}\at{\alpha}\stackrel{Def. 3.2}{=}\alpha([df\at{\alpha},dh\at{\alpha}]_{R})\]
\[=\alpha([Rdf\at{\alpha},dh\at{\alpha}])+\alpha([df\at{\alpha},Rdh\at{\alpha}])=\alpha(ad_{Rdf\at{\alpha}}\at{dh\at{\alpha}})+\alpha(\underbrace{ad_{df\at{\alpha}}\at{Rdh\at{\alpha}}}_{=-ad_{Rdh\at{\alpha}}\at{df\at{\alpha}}})\]
\[=-ad^{*}_{Rdf\at{\alpha}}\at{\alpha}(dh\at{\alpha})+ad^{*}_{Rdh\at{\alpha}}\at{\alpha}(df\at{\alpha})\equiv-ad^{*}_{Rdf\at{\alpha}}\at{\alpha}(h)+ad^{*}_{Rdh\at{\alpha}}\at{\alpha}(f)\]
Ma dal momento che per ipotesi $h$ è di Casimir il primo addendo si annulla 
\[-ad^{*}_{Rdf\at{\alpha}}\at{\alpha}(dh\at{\alpha})=\alpha(ad_{Rdf\at{\alpha}}\at{dh\at{\alpha}})=\alpha([Rdf\at{\alpha},dh\at{\alpha}])\equiv \{h,k\}\at{\alpha}=0\]
in cui la cui forma di $k\in C^{\infty}(\mathfrak{g}^{*})$ non ci interessa. Abbiamo concluso.
\end{proof}

\begin{Cor}
Sia $(\mathfrak{g}^{*},\{,\})$ una varietà di Poisson, se $f,g\in C^{\infty}(\mathfrak{g}^{*})$ sono di Casimir allora commutano rispetto alle parentesi indotte $\{,\}_{R}$.
\end{Cor}
\begin{proof}
Infatti per ogni $\alpha\in\mathfrak{g}^{*}$
\[\{f,g\}_{R}\at{\alpha}=X^{R}_{f}\at{\alpha}(g)=0\]
e, per quanto visto nella dimostrazione precedente, il risultato segue immediatamente perchè entrambre di Casimir.
\end{proof}

\begin{flushleft}
$\textbf{NOTA}$:Il flusso di $X^{R}_{h}$ per ogni funzione $h$ di Casimir risiede nell'orbita coaggiunta sia di $(\mathfrak{g},[\cdot,\cdot])$ che di $(\mathfrak{g},[\cdot,\cdot]_{R})$. In particolare sarà un $ad^{*}$ di qualcosa ma allo stesso tempo un $ad^{*}_{R}$ di qualcos'altro. Questo perchè posso svolgere lo stesso ragionamento fatto nella Proposizione precedente rispetto alle parentesi non modificate $[,]$.
\end{flushleft}



\subsection{Forme $Ad$-invarianti ed $R$-matrici}
L'idea è quella di generalizzare la formula per campi Hamiltoniani sul duale di un algebra di Lie, nel caso sia presente una $R-$matrice ed una forma $Ad-$invariante su di essa. 

Sia $(\cdot,\cdot)$ una forma bilineare (e non-degenere) Ad-invariante sull'algebra di Lie $\mathfrak{g}$. Grazie a quanto visto, possiamo identificare una corrispondenza tra $C^{\infty}(\mathfrak{g})$ e $C^{\infty}(\mathfrak{g}^{*})$

\begin{Def}
Date due funzioni $f,g\in C^{\infty}(\mathfrak{g})$, definiamo le parentesi di Poisson su $\mathfrak{g}$ come
\begin{equation}
    \{f,g\}\at{X}:=(X,[\nabla f\at{X},\nabla g\at{X}])
\end{equation}
in cui $\nabla f\at{X}:=\chi^{-1}(df\at{X})$ è il gradiente di $f$ in $X$, ossia 
\[(\nabla f\at{X},Y)=\chi(\chi^{-1}(df\at{X}))(Y)=df\at{X}(Y)=\frac{d}{dt}\at[\bigg]{t=0}f(X+tY)\]
dove si è sfruttata la (37) identificando $\mathfrak{g}$ con il suo duale $\mathfrak{g}^{*}$.
\end{Def}

\begin{Def}
La "mappa trasposta" è quell'operatore
\[R^{*}:\mathfrak{g}^{*}\rightarrow\mathfrak{g}^{*}\]
\[\alpha\mapsto R^{*}\alpha\]
tale per cui 
\[R^{*}\alpha(X)=\alpha(RX)\hspace{0.5cm}\forall X\in\mathfrak{g}\]
\end{Def}

Grazie alla Proposizione 2.5,vi è allora una naturale identificazione indotta dalla forma $(\cdot,\cdot)$ del tipo
\[\mathfrak{g}\stackrel{\sim}{\longrightarrow} \mathfrak{g}^{*}\]
\[L\mapsto (L,\cdot)=\alpha\]
\[L \longleftrightarrow \alpha\]
Tale isomorfismo induce anche le sguenti corrispondenze
\[h\in C^{\infty}(\mathfrak{g}^{*})\longleftrightarrow f\in C^{\infty}(\mathfrak{g})\]
\[dh\at{\alpha}\longleftrightarrow\nabla f\at{L}\]
\[ad^{*}_{X}\at{\alpha}\longleftrightarrow -ad_{X}\at{L}\]
\[(R^{*}:\mathfrak{g}^{*}\rightarrow\mathfrak{g}^{*})\longleftrightarrow(R^{\dagger}:\mathfrak{g}\simeq\mathfrak{g}^{*}\stackrel{R^{*}}{\rightarrow}\mathfrak{g}^{*}\simeq\mathfrak{g})\]
dove $R^{\dagger}$ è l'operatore indotto dal trasposto $R^{*}$ su $\mathfrak{g}^{*}$.

Ricordiamo ora che l'equazione di Hamilton si scrive come
\[\dot{x}=X^{R}_{h}\at{x}=X_{h}^{R}(x)\longleftrightarrow\frac{dL}{dt}=X_{h}^{R}\at{L}\]
Possiamo dunque formulare la seguente

\begin{Prop}
Data $h\in C^{\infty}(\mathfrak{g}^{*})$, si hanno le seguenti corrispondenze per il campo vettoriale Hamiltoniano rispetto alle R-parentesi di Poisson
\[X^{R}_{h}\at{\alpha}\longleftrightarrow\frac{dL}{dt}=-ad_{R\nabla f\at{L}}\at{L}-R^{\dagger}ad_{\nabla f\at{L}}\at{L}=[L,R\nabla f\at{L}]+R^{\dagger}[L,\nabla f\at{L}]\]
\end{Prop}
\begin{proof}
Per ogni $Y\in\mathfrak{g}$, usando la dimostrazione della Proposizione 3.2
\[X^{R}_{h}\at{\alpha}(Y)=\alpha([Y,dh\at{\alpha}]_{R})=\alpha([RY,dh\at{\alpha}])+\alpha([Y,Rdh\at{\alpha}])\]
\[=\alpha(-ad_{dh\at{\alpha}}\at{RY})+\alpha(-ad_{Rdh\at{\alpha}}\at{Y})=ad^{*}_{dh\at{\alpha}}\at{\alpha}(RY)+ad^{*}_{Rdh\at{\alpha}}\at{\alpha}(Y)\]
\[=(R^{*}(ad^{*}_{dh\at{\alpha}}\at{\alpha})+ad^{*}_{Rdh\at{\alpha}}\at{\alpha})(Y)\]
il che, sfruttando le identificazioni sopra introdotte, implica
\[X^{R}_{h}\at{\alpha}=R^{*}(ad^{*}_{dh\at{\alpha}}\at{\alpha})+ad^{*}_{Rdh\at{\alpha}}\at{\alpha}\longleftrightarrow -R^{\dagger}ad_{\nabla f\at{L}}\at{L}-ad_{R\nabla f\at{L}}\at{L}\]
\end{proof}

\begin{Cor}
Usando la precedente notazione, se $h\in C^{\infty}(\mathfrak{g}^{*})$ è di Casimir per la varietà di Poisson $(\mathfrak{g}^{*},\{,\})$, allora
\begin{equation}
    \frac{dL}{dt}=-ad_{R\nabla f\at{L}}\at{L}=[L,R\nabla f\at{L}]
\end{equation}
\end{Cor}
\begin{proof}
Segue immediatamente sfruttando lo stesso procedimento svolto nella Proposizione 3.2.
\end{proof}

\begin{flushleft}
$\textbf{NOTA}$: Si consideri un elemento $\alpha\in\mathfrak{{g}^{*}}$ che sia Ad-invariante, ossia tale per cui $Ad^{*}_{X}\alpha=\alpha\hspace{0.3cm}\forall X\in\mathfrak{g}^{*}$, in tal caso
\[ad^{*}_{X}\at{\alpha}=\frac{d}{dt}\at[\bigg]{t=0}Ad^{*}_{\phi^{X}_{t}(e)}\alpha=0\]
e quindi $\alpha$ è in tal caso costante lungo le orbite coaggiunte.
 
$\underline{Esempio}$: Nel caso $\mathfrak{g}=\mathfrak{gl}(n)$, sia $Ad_{g}X=gXg^{-1}$. Definiamo $(X,Y):=Tr(XY)$, tale mappa è Ad-invariante
\[(Ad_{g}X,Ad_{g}Y)=Tr(gXg^{-1}gYg^{-1})=Tr(gXYg^{-1})\]
\[=Tr(gg^{-1}XY)=Tr(XY)=(X,Y)\]
dove si è sfruttata l'invarianza sotto permutazioni della traccia. Delle funzioni di Casimir per l'algebra di Poisson $(\mathfrak{g}\simeq\mathfrak{g}^{*},\{,\})$ sono $h_{k}(X):=Tr(X^{k})$ dal momento che
\[Ad_{g}^{*}h_{k}(X)=h_{k}(Ad_{g}X)=h_{k}(gXg^{-1})=Tr(\underbrace{gXg^{-1}gXg^{-1}...gXg^{-1}}_{k-volte})\]
\[=Tr(gX^{k}g^{-1})=Tr(X^{k})=h_{k}(X)\]
ossia tali funzioni sono Ad-invarianti, e dunque costanti lungo le orbite coaggiunte ($\{f,h_{i}\}\at{\xi}=\xi([df\at{\xi},dh_{i}\at{\xi}])=-ad^{*}_{df\at{\xi}}\at{\xi}(dh_{i}\at{\xi})=0$, vedasi NOTA precedente). In particolare si osservi come da una parte
\[(\nabla h_{k}\at{X},Y)=Tr(\nabla h_{k}\at{X}Y)\]
e dall'altra, come visto nella Definizione 3.3
\[(\nabla h_{k}\at{X},Y)=dh_{k}\at{X}(Y)=\frac{d}{dt}\at[\bigg]{t=0}h_{k}(X+tY)=\frac{d}{dt}\at[\bigg]{t=0}Tr((X+tY)^{k})\]
Se considero allora $Y=X=:L$ otteniamo
\[(\nabla h_{k}\at{L},L)=\frac{d}{dt}\at[\bigg]{t=0}h_{k}(L+tL)=\lim_{t\rightarrow0}\frac{1}{t}[h_{k}(L+tL)-h_{k}(L)]\]
\[=\lim_{t\rightarrow0}\frac{1}{t}[Tr(L^{k}\underbrace{(1+t)^{k}}_{=\sum_{m=0}^{k}\binom{k}{m}t^{m}})-Tr(L^{k})]=Tr(\binom{k}{1}L^{k})=Tr(kL^{k})\]
\[=(kL^{k-1},L)\]
In altre parole abbiamo trovato che $\nabla h_{k}\at{L}=kL^{k-1}$.

Per l'ultimo Corollario visto si ha che ogni funzione di Casimir $h$ nell'algebra di Poisson $\mathfrak{g}^{*}$ permette di trovare l'evoluzione di $X\in\mathfrak{g}$
\begin{equation}
\frac{dX}{dt}=[X,R\nabla h\at{X}]
\end{equation}
Le $h_{k},k=1,...,n$ sono in involuzione mutua e rappresentano dunque simmetrie del sistema.
\end{flushleft}



\section{Rappresentazione di Lax}
La teoria moderna sui sistemi integrabili nacque grazie all'osservazione, dovuta a Peter Lax, che a volte è possibile riscrivere delle equazioni differenziali non lineari nella "forma di Lax" 
\[\frac{dL}{dt}=[L,M]\]
dove L,M sono operatori visti come funzioni delle variabili dinamiche. Come vedremo, gli autovalori di L risultano integrali primi.

\begin{Def}
Dato un campo vettoriale $X\in\mathcal{T}(N)$ sulla varietà differenziabile N, $U\subseteq N$ un aperto, una "rappresentazione di Lax di $(X,U)$" è una coppia $(L,M)$ dove $L:U\hookrightarrow\mathfrak{gl}(n,\mathbb{R})$ è un immersione ed $M:U\rightarrow\mathfrak{gl}(n,\mathbb{R})$ è una mappa liscia in modo che sia soddisfatta la relazione
\begin{equation}
(L_{*}X)\at{L(x)}=[L(x),M(x)]
\end{equation}
o, in modo equivalente
\[\frac{dL}{dt}(x)=[L(x),M(x)]\]
visto come un campo vettoriale su $Im(L)$.
\end{Def}

\begin{Prop}
Sia $(L,M)$ una coppia di Lax, allora gli autovalori $\lambda_{1}(x),...,\lambda_{n}(x)$ di L sono integrali del moto
\[\frac{d\lambda_{i}(x(t))}{dt}=0\hspace{0.5cm}\forall i=1,...,n\]
\end{Prop}
\begin{proof}
Si definiscano i valori $t_{k}:=Tr(L^{k})$, si osservi che il polinomio caratteristico associato ad $L$ si può scrivere come
\[Det(L-\lambda\mathbb{1})=(-1)^{n}Det(\lambda\mathbb{1}-L)=(-1)^{n}(\lambda^{n}+c_{1}\lambda^{n-1}+...+c_{n})\]
dove si osservi che abbiamo considerato il coefficente corrispondente a $\lambda^{n}$ con $(-1)^{n}$, il quale è associato al prodotto dei termini lungo la diagonale (si verifica facilmente che è l'unica possibile scelta di tale coefficente).

A questo punto possiamo esprimere i coefficenti di tale polinomio in funzione dei valori $t_{k}$ definiti all'inizio, avvalendoci della "formula di Newton"
\[c_{m}=-\frac{t_{m}}{m}+\frac{1}{2!}\sum_{i+j=m}\frac{t_{i}t_{j}}{ij}-\frac{1}{3!}\sum_{i+j+k=m}\frac{t_{i}t_{j}t_{k}}{ijk}+...+(-1)^{m}\frac{t_{1}^{m}}{m!}\]
la cui dimostrazione può essere recuperata in [14].

A questo punto, siccome i coefficenti $c_{i}$ dipendono esplicitamente dalle traccie $t_{j}$, mostrando che quest'ultime sono costanti del moto avremo finito (il polinomio caratteristico rimane lo stesso, dunque pure le sue soluzioni). Ma effettivamente se si sfrutta l'invarianza ciclica della traccia
\[\frac{d}{dt}Tr(L^{k})=Tr(\frac{dL^{k}}{dt}(t))\stackrel{cycl. inv.}{=}kTr(\dot{L}L^{k-1})=kTr(\frac{dL}{dt}L^{k-1})\]
\[=kTr([L,M]L^{k-1})=kTr(LML^{k-1}-ML^{k})=0\]
dal momento che $Tr(ABC)=Tr(BCA)=Tr(CAB)$ per ogni scelta di $A,B,C\in M(n,\mathbb{R})$. 
\end{proof}

\begin{Teo}
Sia $(M,L)$ una coppia di Lax per un qualche campo vettoriale sulla varietà N, e sia $G:N\rightarrow GL(n)$ una mappa differenziabile. Allora
\[L':=GLG^{-1}\hspace{0.5cm}M':=GMG^{-1}-\frac{dG}{dt}G^{-1}\]
costituisce un ulteriore rappresentazione di Lax per il campo vettoriale considerato.
\end{Teo}
\begin{proof}
Partiamo dalla relazione 
\[\frac{d(GG^{-1})}{dt}=0\Rightarrow G\frac{dG^{-1}}{dt}+\frac{dG}{dt}G^{-1}=0\Rightarrow \frac{dG^{-1}}{dt}=-G^{-1}\frac{dG}{dt}G^{-1}\]
troviamo
\[\frac{dL'}{dt}=\frac{d}{dt}(GLG^{-1})=\frac{dG}{dt}LG^{-1}+\textcolor{red}{G[L,M]G^{-1}}-GLG^{-1}\frac{dG}{dt}G^{-1}\]
\[=\frac{dG}{dt}\overbrace{G^{-1}\textcolor{blue}{G}}^{=\mathbb{1}}\textcolor{blue}{LG^{-1}}+\textcolor{red}{[GLG^{-1},GMG^{-1}]}-\textcolor{blue}{GLG^{-1}}\frac{dG}{dt}G^{-1}\]
\[=[\frac{dG}{dt}G^{-1},\textcolor{blue}{L'}]+[L',GMG^{-1}]=[L',M']\]
\end{proof}

\begin{Cor}
Gli integrali del moto di L sono gli stessi di L'.
\end{Cor}
\begin{proof}
Ovvia dal momento che G rappresenta un semplice cambio di base.
\end{proof}



\subsection{Esempio: Toda lattice}
Si consideri un sistema meccanico di $n$ particelle $q_{1},...,q_{n}$ giacenti nella retta reale ed interagenti mediante il potenziale
\[V(q_{1},...,q_{n})=\sum_{i=1}^{n-1}e^{q_{i}-q_{i+1}}\]
ossia dove ognuna di esse interagisce con solamente la sua prima vicina mediante una forza forza di tipo esponenziale. Tale potenziale prende il nome di "Toda lattice", e per tale sistema l'Hamiltoniana assumerà la forma 
\[H(q_{1},...,q_{n},p_{1},...,p_{n})=\sum_{i=1}^{n}\frac{p_{i}^{2}}{2}+\sum_{i=1}^{n-1}e^{q_{i}-q_{i+1}}\]
Le equazioni di Hamilton per tale sistema sono
\[\dot{q_{i}}=\frac{\partial H}{\partial p_{i}}=p_{i}\hspace{0.5cm}\forall i=1,...,n\]
\[\dot{p_{i}}=-\frac{\partial H}{\partial q_{i}}=\begin{cases}
-e^{q_{1}-q_{2}}\hspace{0.5cm}se\hspace{0.2cm}i=1\\
e^{q_{i-1}-q_{i}}-e^{q_{i}-q_{i+1}}\\
e^{q_{n}-q_{n-1}}\hspace{0.5cm}se\hspace{0.2cm}i=n
\end{cases}\]
Introduciamo ora le suguenti variabili
\[\begin{cases}
a_{i}:=\frac{1}{2}e^{\frac{q_{i}-q_{i+1}}{2}}\hspace{0.5cm}i=1,...,n-1\\
b_{i}:=-\frac{p_{i}}{2}\hspace{1.4cm}i=1,...,n
\end{cases}\]
In questo modo possiamo riscrivere
\[\begin{cases}
\dot{a_{i}}=(\frac{\dot{q}_{i}-\dot{q}_{i+1}}{2})a_{i}=(\frac{p_{i}-p_{i+1}}{2})a_{i}=(b_{i+1}-b_{i})a_{i}\hspace{0.5cm}i=1,...,n-1\\
\dot{b_{i}}=-\frac{1}{2}\dot{p}_{i}=-\frac{1}{2}(e^{q_{i-1}-q_{i}}-e^{q_{i}-q_{i+1}})=2(a_{i}^{2}-a_{i-1}^{2})\hspace{0.2cm}i=1,...,n
\end{cases}\]

Vogliamo mostrare che tale sistema si può porre in forma di Lax:si consideri ora $\mathfrak{g}=\mathfrak{gl}(n)$ come fatto nella sezione precedente e la si divida nella somma diretta $\mathfrak{g}=\mathfrak{g}_{+}\oplus\mathfrak{g}_{-}$ dove $\mathfrak{g}_{+}$ è lo spazio delle matrici triangolari inferiori e $\mathfrak{g}_{-}=\mathfrak{so}(n)$ lo spazio di quelle antisimmetriche. Chiaramente infatti $\mathfrak{g}_{+}\cap\mathfrak{g}_{-}=\emptyset$ e posso sempre scrivere $A\in\mathfrak{g}$ come somma di due matrici $S\in\mathfrak{g}_{-}$ e $T\in\mathfrak{g}_{+}$, $A=S+T$ ponendo
\[\begin{pmatrix}
a_{11}&...&a_{1n}\\
.&.&.\\
.&.&.\\
a_{n1}&...&a_{nn}
\end{pmatrix}=\begin{pmatrix}
a_{11}&0&...&0\\
a_{12}+a_{21}&a_{22}&...&0\\
.&.&.&.\\
a_{1n}+a_{n1}&...&a_{n-1,n}+a_{n,n-1}&a_{nn}
\end{pmatrix}+\begin{pmatrix}
0&a_{12}&...&a_{1n}\\
-a_{12}&0&...&a_{2n}\\
.&.&.&.\\
-a_{1n}&...&-a_{n-1,n}&0
\end{pmatrix}\]

In questo modo si ha
\[\mathfrak{g}_{+}^{*}\simeq\mathfrak{g}_{-}^{\perp}=\{matrici\hspace{0.2cm}simmetriche\}\]
mentre
\[\mathfrak{g}_{-}^{*}\simeq\mathfrak{g}_{+}^{\perp}=\{matrici\hspace{0.2cm}triangolari\hspace{0.2cm}inferiori(strettamente)\}\]
e possiamo così decomporre anche il duale $\mathfrak{g}^{*}=\mathfrak{g}^{*}_{+}\oplus\mathfrak{g}^{*}_{-}$. Si considerino poi le proiezioni
\[
\begin{cases}pr_{+}(L):=\underline{l}(L)+\overline{u}(L)^{T}+d(L)\\
pr_{-}(L):=\overline{u}(L)-\overline{u}(L)^{T}\end{cases}\]
laddove $\underline{l}$ estrae la parte inferiore trangolare, $\overline{u}$ estrae la parte triangolare superiore e d la parte diagonale. Allo stesso modo definiamo le proiezioni sul duale
\[\begin{cases}
pr^{*}_{+}(L)=d(L)+\overline{u}(L)+\overline{u}(L)^{T}\\
pr^{*}_{-}(L)=\underline{l}(L)-\overline{u}(L)^{T}
\end{cases}\]
Come avevamo fatto a suo tempo definiamo la mappa R così costruita
\[\begin{cases}
R(L)=pr_{+}(L)-pr_{-}(L)=\underline{l}(L)+2\overline{u}(L)^{T}+d(L)-\overline{u}(L)\\
R^{*}(L)=pr_{+}^{*}(L)-pr_{-}^{*}(L)=d(L)+\overline{u}(L)+2\overline{u}(L)^{T}-\underline{l}(L)
\end{cases}\]
Quello che faremo ora sarà costruire una gerarchia di equazioni di Hamilton che commutino su $\mathfrak{g}$.

Associamo una struttura di Poisson a $\mathfrak{g}$
\[\pi^{\sharp}_{R}(A)\at{L}=[L,RA]+R^{*}[L,A]\]

Riprendiamo ora l'esempio della sezione precedente, e ricordiamo che su $\mathfrak{gl}(n)$ è definita la forma Ad-invariante $(X,Y):=Tr(XY)$, e che delle funzioni di Casimir per l'algebra di Poisson $(\mathfrak{g}^{*}\simeq\mathfrak{g},\{,\})$ sono le $h_{k}(L):=\frac{1}{k}Tr(L^{k})$. Denotiamo ora 
\[S:=\{b+aT_{+}+T_{-}a\}\subset\mathfrak{g}\]
laddove
\[a=\begin{pmatrix}
a_{1}&0&.&.&0\\
0&a_{2}&.&.&.\\
.&.&.&.&.&\\
.&.&.&a_{n-1}&0\\
0&.&.&0&0
\end{pmatrix},b=\begin{pmatrix}
b_{1}&0&.&.&0\\
0&b_{2}&.&.&.\\
.&.&.&.&.&\\
.&.&.&b_{n-1}&0\\
0&.&.&0&b_{n}
\end{pmatrix}\]
\[T_{+}=
\begin{pmatrix}
0&1&.&.&0\\
0&0&1&.&.\\
.&.&.&.&.&\\
.&.&.&0&1\\
0&.&.&0&0
\end{pmatrix},T_{-}=\begin{pmatrix}
0&.&.&.&0\\
1&0&.&.&.\\
.&1&.&.&.&\\
.&.&.&0&.\\
0&.&.&1&0
\end{pmatrix}\]
Si può dimostrare che per ogni $L\in S$ ed $A\in\mathfrak{g}$,$\pi_{R}^{\sharp}(A)\at{L}\in S$ e che dunque vi sia un unica struttura di Poisson su S. Restringendosi a tale spazio si ottiene poi, per quanto visto nell'ultimo esempio della sezione 3, che
\[\frac{dL}{dt}=[L,R\nabla h_{k}\at{L}]=[L,RL^{k-1}]\]
Considerando ad esempio $k=2$, si ottiene
\[\frac{dL}{dt}=[L,R\nabla h_{2}\at{L}]=[L,RL]\iff\begin{cases}
\dot{a_{i}}=(b_{i+1}-b_{i})a_{i}\\
\dot{b_{i}}=2(a_{i}^{2}-a_{i-1}^{2})
\end{cases}\]
Questo perchè
\[L=b+aT_{+}+T_{-}a\]
ma 
\[aT_{+}=\begin{pmatrix}
0&a_{1}&.&.&0\\
0&0&a_{2}&.&.\\
.&.&.&.&.&\\
.&.&.&0&a_{n-1}\\
0&.&.&.&0
\end{pmatrix},T_{-}a=\begin{pmatrix}
0&.&.&.&0\\
a_{1}&0&.&.&.\\
.&a_{2}&0&.&.&\\
.&.&.&0&.\\
0&.&.&a_{n-1}&0
\end{pmatrix}\]
\[\Rightarrow L=\begin{pmatrix}
b_{1}&a_{1}&0&.&0\\
a_{1}&b_{2}&.&.&.\\
0&.&.&.&0&\\
.&.&.&.&a_{n-1}\\
0&.&0&a_{n-1}&b_{n}
\end{pmatrix}\]
Segue allora che (vedasi forma esplicita di $R(L)$ a pagina precedente) 
\[R(L)=\begin{pmatrix}
b_{1}&-a_{1}&0&.&0\\
3a_{1}&b_{2}&.&.&.\\
0&.&.&.&0&\\
.&.&.&.&-a_{n-1}\\
0&.&0&3a_{n-1}&b_{n}
\end{pmatrix}\]
Possiamo dunque esplicitare l'equazione di Lax
\[\begin{pmatrix}
\dot{b_{1}}&\dot{a_{1}}&0&.&0\\
\dot{a_{1}}&\dot{b_{2}}&.&.&.\\
0&.&.&.&0&\\
.&.&.&.&\dot{a}_{n-1}\\
0&.&0&\dot{a}_{n-1}&\dot{b_{n}}
\end{pmatrix}=\begin{pmatrix}
b_{1}&a_{1}&0&.&0\\
a_{1}&b_{2}&.&.&.\\
0&.&.&.&0&\\
.&.&.&.&a_{n-1}\\
0&.&0&a_{n-1}&b_{n}
\end{pmatrix}\begin{pmatrix}
b_{1}&-a_{1}&0&.&0\\
3a_{1}&b_{2}&.&.&.\\
0&.&.&.&0&\\
.&.&.&.&-a_{n-1}\\
0&.&0&3a_{n-1}&b_{n}
\end{pmatrix}
\]
\[-\begin{pmatrix}b_{1}&-a_{1}&0&.&0\\
3a_{1}&b_{2}&.&.&.\\
0&.&.&.&0&\\
.&.&.&.&-a_{n-1}\\
0&.&0&3a_{n-1}&b_{n}\end{pmatrix}\begin{pmatrix}
b_{1}&a_{1}&0&.&0\\
a_{1}&b_{2}&.&.&.\\
0&.&.&.&0&\\
.&.&.&.&a_{n-1}\\
0&.&0&a_{n-1}&b_{n}
\end{pmatrix}\]
\[=\begin{pmatrix}
b_{1}^{2}+3a_{1}^{2}&-a_{1}b_{1}+a_{1}b_{2}&0&.&0\\
a_{1}b_{1}+3a_{1}b_{2}&b_{2}^{2}-a_{1}^{2}+3a_{2}^{2}&.&.&.\\
0&.&.&.&.&\\
.&.&.&.&.&.\\
0&.&.&.&.&.
\end{pmatrix}-\begin{pmatrix}
b_{1}^{2}-a_{1}^{2}&a_{1}b_{1}-a_{1}b_{2}&0&.&0\\
3a_{1}b_{1}+a_{1}b_{2}&b_{2}^{2}-a_{2}^{2}+3a_{1}^{2}&.&.&.\\
0&.&.&.&.&\\
.&.&.&.&.\\
0&.&.&&
\end{pmatrix}\]
\[=\begin{pmatrix}
4a_{1}^{2}&2a_{1}(b_{2}-b_{1})&0&.&0\\
2a_{1}(b_{2}-b_{1})&4(a_{2}^{2}-a_{1}^{2})&.&.&.\\
0&.&.&.&.&\\
.&.&.&.&\\
0&.&.&.&
\end{pmatrix}\]

Abbiamo ottenuto proprio le equazioni del Toda lattice ricavate in prima istanza (a meno di un fattore 2, questo è dovuto al fatto che in realtà andrebbe considerato $\frac{1}{2}(pr_{+}-pr_{-})=\frac{R(L)}{2}$ come visto nelle sezioni precedenti). Questo è un classico esempio di sistema che non risulta integrabile ad una prima analisi ma che, con l'utilizzo delle funzioni di Casimir, può essere integrato facilmente.





\section{Caso $\infty$-dimensionale}
A questo punto vogliamo estendere la nozione di sistema integrabile al caso in cui si abbia a che fare con sistemi ad infiniti gradi di libertà. Questa esigenza nasce qualora si voglia ad esempio descrivere come si comporti il profilo di una corda vibrante, oppure il profilo di un onda sotto opportune condizioni al contorno.

Sinora abbiamo infatti sempre considerato sistemi a finiti gradi di libertà, il che si traduce nell'avere a che fare con un numero finito di coordinate  $(q^{1},...,q^{n},p^{1},...,p^{n})$ sulla varietà. Per estendere tutto ciò al caso infinito vi sono due possibilità:
\begin{itemize}
    \item [(a)] La prima consiste nel passare al limite $n\rightarrow\infty$, mantenendo così una cardinalità discreta delle coordinate;
    \item [(b)] La seconda invece è quella di associare alla generica $q_{i}$ un indice continuo $\{q_{i}^{x}\}_{x\in\mathbb{R}}=:\{u^{i}(x)\}_{x\in\mathbb{R}}$, in modo da poter descrivere il comportamento di un infinità di punti anzichè della singola coordinata (ad esempio,per una corda avrò infiniti punti materiali che si muovono).
\end{itemize}  
Per farlo però, sono necessari alcuni accorgimenti che andremo a rimarcare quando sarà necessario, come prima cosa partiamo dalla
 
\begin{Def}
Un "sistema evolutivo di PDE" è un sistema di equazioni differenziali alle derivate parziali della forma

\begin{equation}
\frac{\partial u^{\alpha}}{\partial t}=F^{\alpha}(u,u_{1},u_{2}...)\hspace{1cm}\alpha=1,...,N
\end{equation}
laddove si sono indicate le derivate k-esime di u con $u_{k}^{\alpha}:=\partial_{x}^{k}u^{\alpha}$.
\end{Def}
\begin{flushleft}

Si osservi che le variabili $u^{\alpha}$ vanno intese come funzioni della variabile $x$!

Come prima cosa è necessario decidere che tipo di funzioni vogliamo studiare, il che consiste nel dichiarare il dominio di definizione delle varie $u^{\alpha}$:scegliamo qui di trattare le funzioni definite sul cerchio $\mathbb{S}^{1}$, il caso in cui il dominio sia tutto l'asse reale si tratta in modo analogo ma con opportune correzioni.

Tale sistema può essere interpretato come un campo vettoriale sullo spazio infinito-dimensionale di tutti i "loop" $u:\mathbb{S}^{1}\rightarrow V$, con V uno spazio vettoriale N-dimensionale con base $\{e_{1},...,e_{N}\}$ ed $x$ la coordinata su $\mathbb{S}^{1}$; in modo che $u^{\alpha}$ sia la componente lungo $e_{\alpha}$ di tale loop. E' come se avessimo dunque una varietà liscia in cui le coordinate siano funzioni, in questo caso definite sul cerchio.

Sia ora $N\geq1$ e si consideri la variabile $u^{\alpha}_{d}$ con $1\leq\alpha\leq N$ e $d\geq0$. 
\end{flushleft}

\begin{Def}
"L'anello dei polinomi differenziali" $\mathcal{A}_{u^{1},...,u^{N}}=\mathcal{A}$ è definito come l'anello dei polinomi nelle variabili $u_{i>0}^{\alpha}$ con coefficenti appartenenti all'anello delle serie formali costruite nelle variabili $u_{0}^{\alpha}$. In altre parole un elemento $f(x,u,u_{x},...)\in\mathcal{A}$ sarà esprimibile nella forma
\begin{equation}
    f(u,u_{x},...)=\sum_{m=0}^{<\infty}\sum_{I_{m},S_{m}}f_{I_{m}}^{S_{m}}(u)u_{s_{1}}^{i_{1}}...u_{s_{m}}^{i_{m}}
    =\sum_{m=0}^{<\infty}\sum_{1\leq i_{1},...,i_{m}\leq N\atop s_{1},...,s_{m}\geq1}f_{i_{1},...,i_{m}}^{s_{1},...,s_{m}}(u)u_{s_{1}}^{i_{1}}...u_{s_{m}}^{i_{m}}
\end{equation}
dove nella prima scrittura si è usata la notazione compatta con multi-indici $I_{m}:=(i_{1},...,i_{m}),S_{m}:=(s_{1},...,s_{m})$.
A loro volta i coefficenti $f_{i_{1},...,i_{m}}^{s_{1},..,s_{m}}(u)$ saranno esprimibili come serie formali in u con coefficenti in $\mathbb{C}$
\begin{equation}
    f_{i_{1},...,i_{m}}^{s_{1},..,s_{m}}(u)=\sum_{j_{1},...,j_{N}\geq0}\underbrace{C_{i_{1},...,i_{m};j_{1},...,j_{N}}^{s_{1},..,s_{m}}}_{\in\mathbb{C}}(u^{1})^{j_{1}}...(u^{N})^{j_{N}}
\end{equation}
Definiamo poi il grado del monomio $u_{i}^{\alpha}$ come $deg(u_{i}^{\alpha}):=i$, mentre indicheremo con $\mathcal{A}^{[k]}$ il sottospazio di $\mathcal{A}$ di grado k.
\end{Def}

\begin{flushleft}
$\textbf{NOTA}$:Possiamo poi differenziare i polinomi così definiti in $\mathcal{A}$ mediante l'operatore differenziale 
\[\partial_{x}:\mathcal{A}\rightarrow\mathcal{A}\]
\[\partial_{x}:=\sum_{i\geq 0}u_{i+1}^{\alpha}\frac{\partial}{\partial u^{\alpha}_{i}}\]
dove sfrutteremo la convenzione di Einstein sugli indici greci, ma non sugli indici romani.Si noti anche che, nella formula (49), gli indici $i_{j}$ si riferiscono alla variabile considerata, mentre le potenze si formano con l'avanzare dell'indice m.

Introduciamo adesso la nuova variabile formale  $\epsilon$ e definiamone il grado $deg(\epsilon):=-1$.
\end{flushleft}

\begin{Def}
Defniamo ora il nuovo anello 
\[\widehat{\mathcal{A}}:=\mathcal{A}[[\epsilon]]=\mathbb{C}[[u]][u_{>0}][[\epsilon]]\]
costituita da polinomi del tipo
\[\sum_{\mu\geq0}X_{\mu}\epsilon^{\mu}\hspace{1cm}X_{\mu}\in\mathcal{A}\]
Infine denotiamo con $\widehat{\mathcal{A}}^{[k]}$ il sottospazio di $\widehat{\mathcal{A}}$ di grado k.

Definiamo poi lo "spazio dei funzionali locali" come il quoziente $\widehat{\Lambda}:=\widehat{\mathcal{\mathcal{A}}}/(Im\partial_{x}\oplus\mathbb{C}[[\epsilon]])$, e ne denotiamo $\widehat{\Lambda}^{[k]}$ la sua componente di grado k. La classe di equivalenza $f(u;\epsilon)\in\widehat{\Lambda}$ verrà espressa in modo più suggestivo con $\overline{f}=\int_{\mathbb{S}^{1}} f(u,u_{1},u_{2},...;\epsilon)dx$.
\end{Def}

Si osservi anzitutto che $\partial_{x}$ è ben definito anche in $\widehat{\mathcal{A}}$.
Definiamo invece in $\widehat{\Lambda}$ la "derivata variazionale"
\[\frac{\delta}{\delta u^{\alpha}}:\widehat{\Lambda}\rightarrow\widehat{\mathcal{A}}\]
\[\frac{\delta\overline{f}}{\delta u^{\alpha}}:=\sum_{i\geq0}(-\partial_{x})^{i}\frac{\partial f}{\partial u^{\alpha}_{i}}\]
Inoltre ogni qualvolta $f(u,u_{1},u_{2},...;\epsilon)\in\widehat{\mathcal{A}}$ coincide con la derivata totale di una qualche $g(u,u_{1},u_{2},...;\epsilon)\in\widehat{\mathcal{A}}$ si trova
\[\overline{f}=\int_{\mathbb{S}^{1}}f(u,u_{1},u_{2},...;\epsilon)dx=\int_{0}^{2\pi}\partial_{x}g(u,u_{1},u_{2},...;\epsilon)dx=0\]

\begin{Def}
Una "trasformazione di Miura" è un cambio di coordinate del tipo
\[u^{\alpha}\mapsto\widetilde{u}^{\alpha}(u,u_{1},u_{2},...;\epsilon)=\sum_{k\geq0}f_{k}^{\alpha}(u,u_{1},u_{2},...)\epsilon^{k}\hspace{0.5cm}1\leq\alpha\leq N\]
dove
\[f_{k}^{\alpha}\in\mathcal{A}^{[k]}\hspace{0.5cm}Det(\frac{\partial\widetilde{u}}{\partial u}\at[\bigg]{\epsilon=0})=Det(\frac{\partial f_{0}^{\alpha}}{\partial u})\neq0\]
\end{Def}


I polinomi differenziali ed i loro funzionali locali possono essere descritti usando un ulteriore insieme di variabili formali, cioè mediante le componenti $p_{k},k\in\mathbb{Z}$ della trasformata di Fourier delle $u^{\alpha}(x)$. Definiamo infatti
\[u^{\alpha}_{j}=\sum_{k\in\mathbb{Z}}(ik)^{j}p_{k}^{\alpha}e^{ikx}\]
che non sono altro che le derivate j-esime delle trasformate di Fourier $u^{\alpha}=\sum_{k\in\mathbb{Z}}p_{k}^{\alpha}e^{ikx}$. Questo ci permette di riscrivere il generico termine $f(u,u_{1},u_{2},...;\epsilon)\in\widehat{\mathcal{A}}^{[d]}$ come una serire di Fourer formale
\[f(u,u_{1},u_{2},...;\epsilon)=\sum_{k_{1},...,k_{n}\in\mathbb{Z}}f_{k_{1},...,k_{n};s}^{\alpha_{1},...,\alpha_{n}}\epsilon^{s}p_{k_{1}}^{\alpha_{1}}...p_{k_{n}}^{\alpha_{n}}e^{i(\sum_{j}k_{j})x} \]
dove chiaramente $f_{k_{1},...,k_{n};s}^{\alpha_{1},...,\alpha_{n}}$ deve essere un polinomio di grado s+d(si ricordi che il grado di $\epsilon$ è stato definito negativo).



\subsection{Struttura di Poisson}
In analogia al caso finito dimensionale, ci servono delle parentesi di Poisson per poter parlare di sistemi integrabili.

Ritorniamo per un momento al caso finito, se M è la varietà del nostro sistema, quello che si faceva era associare a due funzioni $f,g\in C^{\infty}(M)$ un ulteriore funzione $\{f,g\}\in C^{\infty}(M)$. Potremmo allora pensare di associare a due funzionali $\overline{f},\overline{g}\in\widehat{\Lambda}$ un ulteriore funzionale, come espresso nella seguente

\begin{Def}
Le "parentesi di poisson" sullo spazio dei funzionali $\widehat{\Lambda}$ sono definite come
\begin{equation}
\{\overline{f},\overline{g}\}_{K}:=\int \frac{\delta\overline{f}}{\delta u^{\mu}}K^{\mu\nu}\frac{\delta\overline{g}}{\delta u^{\nu}}dx\hspace{0.5cm}K^{\mu\nu}:=\sum_{j\geq0}K^{\mu\nu}_{j}\partial_{x}^{j}\hspace{0.5cm}K^{\mu\nu}_{j}\in\widehat{\mathcal{A}}^{[-j+1]}
\end{equation}
L'operatore K si dice "operatore Hamiltoniano"(l'antisimmetria e la validità dell'dentità di Jacobi si restringono le possibili scelte di K).

Definiamo invece le "parentesi di Poisson" tra $f\in\widehat{\mathcal{A}}$ ed $\overline{g}\in\widehat{\Lambda}$ come
\begin{equation}
    \{f,\overline{g}\}:=\sum_{s\geq0}\frac{\partial f}{\partial u_{s}^{\mu}}\partial_{x}^{s}(K^{\mu\nu}\frac{\delta\overline{g}}{\delta u^{\nu}})
\end{equation}
\end{Def}

\begin{flushleft}
$\textbf{NOTA}$:Per quanto detto segue immediatamente la relazione
\begin{equation}
 \int f'dx=0\Rightarrow\int fg'dx=-\int f'gdx
 \end{equation}
Inotre si noti che (52) è compatibile con la (51), infatti 
\[\int\{f,\overline{g}\}_{K}dx=\int\sum_{s\geq0}\frac{\partial f}{\partial u_{s}^{\mu}}\partial_{x}^{s}(K^{\mu\nu}\frac{\delta\overline{g}}{\delta u^{\nu}})dx\stackrel{(53)}{=}\int\sum_{s\geq0}(-\partial_{x})^{s}\frac{\partial f}{\partial u_{s}^{\mu}}K^{\mu\nu}\frac{\delta\overline{g}}{\delta u^{\nu}}dx=\{\overline{f},\overline{g}\}_{K}\]
\end{flushleft}

Il seguente Teorema è simile a quello di Darboux (Teorema 1.8),infatti permette sotto opportune ipotesi di ricondurci ad un operatore Hamiltoniano "canonico".

\begin{Teo}
Si supponga che $K$ sia un operatore Hamiltoniano con prima componente $K\at{\epsilon=0}$ non-degenere. Allora esiste una trasformazione di Miura che induce le parentesi di Poisson alla forma
\[K^{\mu\nu}=\eta^{\mu\nu}\partial_{x}\]
dove $\eta$ è un operatore costante, simmetrico e non-degenere.
\end{Teo}

Per una dimostrazione e riformulazione più precisa rimandiamo a [16].


\subsection{Gerarchie integrabili}
Nel caso finito, quello che si faceva era associare ad ogni coordinata $q^{i}$ (o meglio alla sua derivata temporale) una parentesi di Poisson $\{q^{i},h\}$, dunque scelto il valore della coordinata ottenevo un uguaglianza tra numeri.

Nel caso infinito, la questione è un pochino diversa

\begin{Def}
Un sistema evolutivo di PDE si dice "Hamiltoniano" se si può scrivere come
\begin{equation}
    \partial_{t}u^{\alpha}=\{u^{\alpha},\overline{h}\}_{K}=K^{\alpha\nu}\frac{\delta\overline{h}}{\delta u^{\nu}}\hspace{0.5cm}\overline{h}\in\widehat{\Lambda}^{[0]}
\end{equation}
dove $\overline{h}$ è chiamata "Hamiltoniana" del sistema.

Un "sistema integrabile(o gerarchia integrabile)" è invece un sistema evolutivo di PDE Hamiltoniano infinito
\begin{equation}
    \partial_{t_{d}^{\beta}}u^{\alpha}=\{u^{\alpha},\overline{h}_{\beta,d}\}_{K}=K^{\alpha\mu}\frac{\delta\overline{h}_{\beta,d}}{\delta u^{\mu}}
\end{equation}
generato da una famiglia di Hamiltoniane $\overline{h}_{\alpha,d}\in\widehat{\Lambda}^{[0]},\alpha=1,...,N,d\geq0$ che soddisfino
\begin{equation}
\{\overline{h}_{\alpha,i},\overline{h}_{\beta,j}\}=0
\end{equation}
Una "soluzione" del sistema integrabile è quella serie formale $u^{\alpha}(x,t;\epsilon)\in\mathbb{C}[[x,t_{d}^{\beta};\epsilon]]_{0\leq\beta\leq N;d\geq0}$ che soddisfi tutte le equazioni (32) contemporaneamente. 
\end{Def}

$\textbf{NOTA}$:Si può dimostrare (vedasi [17]) che, nel caso di sistemi Hamiltoniani, il fatto che $(56)$ risulti soddisfatta equivale alla possibilità di trovare un (unica) soluzione $\{u^{\alpha}(x,t_{d}^{\beta};\epsilon)\}_{0\leq\beta\leq N;d\geq0}$ che soddisfi contemporaneamente le equazioni $(55)$.

Si osservi che in questo contesto quello che si fa è associare alla coordinata $u^{\alpha}$ (o meglio alla sua derivata temporale, che rimane una funzione definita sul cerchio) una parentesi $\{u^{\alpha},\overline{h}\}_{K}$ che rimane un elemento di $\widehat{\mathcal{A}}$. Dunque stavolta, scelta la coordinata $u^{\alpha}$ (rimarchiamo ancora che tale scelta equivale non più a fissare il valore di una coordinata, ma piuttosto a fissare la forma della funzione che costituisce ora la mia coordinata), otterrò un uguaglianza tra funzioni.

Infine si osservi come stavolta, l'Hamiltoniana del sistema può considerarsi ancora una funzione, tuttavia vista come definita sullo spazio dei loop precedentemente costruito.Come vedremo nei capitoli seguenti, conoscere la derivata temporale delle $u^{i}(x)$ permette di conoscere l'evoluzione temporale di qualsiasi funzionale, proprio come nel caso finito si riesce a conoscere l'evoluzione di ogni funzione nello spazio delle fasi una volta conosciute le derivate temporale $\dot{q}^{i}$.


\section{Esempi e applicazioni:il sistema KdV}
Il concetto di "solitone" è intimamente correlato con l'equazione KdV di cui tratteremo tra poco. Sebbene una definizione matematica ben precisa di solitone sia difficile da stabilire, lo associamo alla soluzione di qualche equazione differenziale (non lineare) ed alle seguenti proprietà:
\begin{itemize}
   \item La sua forma non cambi nel tempo;
    \item Sia localizzato;
    \item Interagisca con altri solitoni in modo che la sua forma si conservi dopo la sua collisione;
   \end{itemize}
La prima osservazione del solitone si riconduce all'ingegnere John Scott Russel, che osservò un'onda solitaria risalire la corrente nell'Union Canal per chilometri senza perdere energia.
Fu solo nel 1877 che lo scienziato Joseph Valentine de Boussinesq pubblicò un articolo che,nell'intento di descrivere tale fenomeno, includeva termini non lineari nell'equazione associata a tale onda e la cui soluzione presentava proprietà simili al solitone osseervato da Russell. Nel 1895 Johannes Korteweg e Gustav de Vries riuscirono finalmente ad includere dei termini dovuti alla tensione superficiale del fluido ed ottenere così l'equazione KdV
\[\frac{\partial u}{\partial t}=uu_{x}+\frac{\epsilon^{2}}{12}u_{xxx}\]
Tale equazione descrive in particolare un onda che si sviluppa in un fluido contenuto in un canale poco profondo ed in un certo regime, dove l'ampiezza dell'onda è piccola rispetto alla lunghezza del canale.Passiamo ora alla trattazione di tale sistema mediante la teoria del paragrafo precedente.

L'equazione di Korteweg-de-Vries costituisce l'esempio meglio conosciuto di sistema Hamiltoniano integrabile. Si tratta del sistema definito sullo spazio dei loop $\mathbb{S}^{1}\rightarrow V=\mathbb{C}$, per questo gli indici greci verranno soppressi nelle equazioni precedentemente definite. La metrica su V è semplicemente $\eta=1$, mentre la struttura di Poisson è data dall'operatore Hamiltoniano $K=\partial_{x}$. L'Hamiltoniana del sistema è invece 
\[\overline{h}_{KdV}:=\int(\frac{u^{3}}{6}+\frac{\epsilon^{2}uu_{2}}{24})dx\]
La corrispondente equazione di Hamilton (31) sarà
\[\partial_{t}u=\{u,\overline{h}_{KdV}\}_{K}\]
dove
\[\{u,\overline{h}_{KdV}\}_{K}=(\partial_{x}\circ\frac{\delta}{\delta u^{1}})\overline{h}_{KdV}=\partial_{x}(\frac{\delta\overline{h}_{KdV}}{\delta u})\]
Calcoliamo allora $(u\equiv u^{1})$
\[\frac{\delta\overline{h}_{KdV}}{\delta u}=\sum_{k\geq0}(-\partial_{x})^{k}\frac{\partial}{\partial u^{1}_{k}}(\frac{u^{3}}{6}+\frac{\epsilon^{2}uu_{2}}{24})\]
\[\frac{\partial}{\partial u}(\frac{u^{3}}{6}+\frac{\epsilon^{2}uu_{2}}{24})+(-1)^{1}(\partial_{x})\frac{\partial}{\partial u_{1}}(\frac{u^{3}}{6}+\frac{\epsilon^{2}uu_{2}}{24})+(-1)^{2}(\partial_{x})^{2}\frac{\partial}{\partial u_{2}}(\frac{u^{3}}{6}+\frac{\epsilon^{2}uu_{2}}{24})\]
\[=\frac{u^{2}}{2}+\frac{\epsilon^{2}}{24}u_{2}+\frac{\epsilon^{2}}{24}u_{2}=\frac{u^{2}}{2}+\frac{\epsilon^{2}}{12}u_{2}\]
\[\Rightarrow\partial_{t}u=uu_{1}+\frac{\epsilon^{2}}{12}u_{3}\]
L'equazione KdV è solo un esempio di sistema evolutivo che fa parte di un intera gerarchia integrabile. Ci sono diversi modi di costruire le altre Hamiltoniane che compongono tale gerarchia, in particolare qui scegliamo di utilizzare un metodo ricorsivo scoperto da A.Buryak e P.Rossi[4].

\begin{figure}
    \centering
    \includegraphics[scale=0.5]{Solitone.jpg}
    \caption{Illustrazione di un solitone ricreato nell'Union Canal(Scozia, 1995)}
    \label{fig:my_label}
\end{figure}



\subsection{Metodo Ricorsivo}
Sia $g_{-1}:=u\in\widehat{\mathcal{A}}^{[0]}$ e si costruiscano le nuove Hamiltoniane $\overline{h}_{i}\in\widehat{\Lambda}^{[0]},i\geq-1$ come $\overline{h}_{i}=\int g_{i}dx$, dove i polinomi differenziali $g_{i}\in\widehat{\mathcal{A}}^{[0]}$ sono prodotti dalla relazione ricorsiva 
\begin{equation}
    g_{i+1}:=(D-1)^{-1}\partial_{x}^{-1}\{g_{i},\overline{h}_{KdV}\}\hspace{1cm}D:=\sum_{k\geq0}(k+1)u_{k}\frac{\partial}{\partial u_{k}}
\end{equation}
Si noti che, ad ogni step, questa equazione produce una nuova densità Hamiltoniana $g_{i}$ la cui parentesi di Poisson rispetto ad $\overline{h}_{KdV}=\overline{h}_{1}$ risulta esatta rispetto alla derivazione $\partial_{x}$, dunque è sensato applicarvi $\partial_{x}^{-1}$. Anche l'operatore D è facilmente invertibile in ogni monomio che compone il polinomio risultante,dal momento che $D$ in $\widehat{\mathcal{A}}^{[0]}$ conta semplicemente il numero di variabili $\epsilon$ ed $u_{*}$.

Si ha 
\begin{itemize}
    \item $g_{-1}=u\rightarrow deg(u)=0$;
    \item $g_{0}=(D-1)^{-1}\partial_{x}^{-1}\{u,\overline{h}_{KdV}\}$, dove 
    \[\{u,\overline{h}_{KdV}\}\stackrel{(52)}{=}\sum_{s\geq0}\frac{\partial u}{\partial u_{s}}\partial_{x}^{s}(\partial_{x}\frac{\delta\overline{h}_{KdV}}{\delta u})=\partial_{x}(\frac{\delta\overline{h}_{KdV}}{\delta u})=\partial_{x}(\frac{\delta\overline{h}_{KdV}}{\delta u})\]
    Ora
    \[\frac{\delta\overline{h}_{KdV}}{\delta u}=\sum_{i\geq0}(-\partial_{x})^{i}\frac{\partial h}{\partial u_{i}}=\frac{\partial}{\partial u}(\frac{u^{3}}{6}+\frac{\epsilon^{2}uu_{2}}{24})+(-1)^{1}\partial_{x}\circ\frac{\partial}{\partial u_{1}}(\frac{u^{3}}{6}+\frac{\epsilon^{2}uu_{2}}{24})\]
    \[+(-1)^{2}\partial_{x}^{2}\circ\frac{\partial}{\partial u_{2}}(\frac{u^{3}}{6}+\frac{\epsilon^{2}uu_{2}}{24})=\frac{\epsilon^{2}}{24}u_{2}+\frac{u^{2}}{2}+\frac{\epsilon^{2}}{24}u_{2}=\frac{u^{2}}{2}+\frac{\epsilon^{2}}{12}u_{2}\]
    Ma allora
    \[g_{0}=(D-1)^{-1}(\partial_{x}^{-1}\circ\partial_{x})(\frac{u^{2}}{2}+\frac{\epsilon^{2}}{12}u_{2})=(D-1)^{-1}\frac{u^{2}}{2}+(D-1)^{-1}\frac{\epsilon^{2}}{12}u_{2}\]
    A questo punto basta osservare che
    \[(D-1)\frac{u^{2}}{2}=u\frac{\partial}{\partial u}(\frac{u^{2}}{2})-\frac{u^{2}}{2}=\frac{u^{2}}{2}\Rightarrow(D-1)^{-1}\at{\frac{u^{2}}{2}}=id\]
    \[(D-1)\frac{\epsilon^{2}u_{2}}{12}=\frac{\epsilon^{2}}{12}[3u_{2}\frac{\partial u_{2}}{\partial u_{2}}-u_{2}]=\frac{2}{12}\epsilon^{2}u_{2}\Rightarrow(D-1)^{-1}\at{\frac{\epsilon^{2}u_{2}}{12}}=\frac{1}{2}\]
    Segue subito che 
    \[g_{0}=\frac{u^{2}}{2}+\frac{\epsilon^{2}}{24}u_{2}\rightarrow deg(g_{0})=2\cdot0=0,+2-2=0;\]
    
    \item $g_{1}=(D-1)^{-1}\partial_{x}^{-1}\{g_{o},\overline{h}_{KdV}\}$, dove (sfruttando quanto trovato al punto precedente)
    \[\{g_{0},\overline{h}_{KdV}\}=\sum_{s\geq0}\frac{\partial g_{0}}{\partial u_{s}}(\partial_{x}^{s}\circ\partial_{x})\frac{\delta\overline{h}_{KdV}}{\delta u}=\sum_{s\geq0}\frac{\partial}{\partial u_{s}}(\frac{u^{2}}{2}+\frac{\epsilon^{2}}{24}u_{2})(\partial_{x}^{s}\circ\partial_{x})(\frac{\delta\overline{h}_{KdV}}{\delta u})\]
    \[=u(\partial_{x}\frac{\delta\overline{h}_{KdV}}{\delta u})+\frac{\epsilon^{2}}{24}\partial_{x}^{3}(\frac{\delta\overline{h}_{KdV}}{\delta u})=u^{2}u_{1}+\frac{\epsilon^{2}}{12}uu_{3}+\frac{\epsilon^{4}}{24}\partial_{x}^{2}(\frac{\delta\overline{h}_{KdV}}{\delta u})\]
    \[\Rightarrow \partial_{x}^{-1}\{g_{0},\overline{h}_{KdV}\}=\partial_{x}^{-1}(u_{1}u^{2}+\frac{\epsilon^{2}}{12}uu_{3})+\frac{\epsilon^{4}}{24}\partial_{x}(\frac{\delta\overline{h}_{KdV}}{\delta u})=\star\]
    A questo punto usiamo il fatto che applicare $\partial_{x}^{-1}$ equivale ad integrare il singolo monomio considerato, dunque
    \[\star=\frac{u^{3}}{3}+\frac{\epsilon^{2}}{12}\partial_{x}^{-1}(uu_{3})+\frac{\epsilon^{2}}{24}\partial_{x}(uu_{1}+\frac{\epsilon^{2}}{12}u_{3})=\frac{u^{3}}{3}+\frac{\epsilon^{2}}{12}\partial_{x}^{-1}(uu_{3})+\frac{\epsilon^{2}}{24}\partial_{x}(uu_{1}+\frac{\epsilon^{2}}{12}u_{3})\]
    \[=\frac{u^{3}}{3}+\frac{\epsilon^{2}}{12}\partial_{x}^{-1}(uu_{3})+\frac{\epsilon^{2}}{24}(uu_{2}+u_{1}^{2}+\frac{\epsilon^{2}}{12}u_{4})=\frac{u^{3}}{3}+\frac{\epsilon^{2}}{12}\partial_{x}^{-1}(uu_{3})+\frac{\epsilon^{2}}{24}(uu_{2}+u_{1}^{2})+\frac{\epsilon^{2}}{288}u_{4}\]
    \[\stackrel{(1)}{=}\frac{u^{3}}{3}+\frac{\epsilon^{2}}{12}(uu_{2}-\frac{u_{1}^{2}}{2})+\frac{\epsilon^{2}}{24}(uu_{2}+u_{1}^{2})+\frac{\epsilon^{2}}{288}u_{4}=\frac{u^{3}}{3}+\frac{\epsilon^{2}}{8}uu_{2}+\frac{\epsilon^{4}}{288}u_{4}\]
    dove in (1) abbiamo sfruttato gli integrali
    \[\int u_{1}u_{2}dx=u_{1}^{2}-\int u_{2}u_{1}dx\Rightarrow\int u_{1}u_{2}dx=\frac{u_{1}^{2}}{2}\] e
    \[\int uu_{3}dx=uu_{2}-\int u_{1}u_{2}dx=uu_{2}-\frac{u_{1}^{2}}{2}\]
    A questo punto basta osservare che
    \[(D-1)^{-1}u^{3}=3u^{3}-u^{3}=2u^{3}\Rightarrow(D-1)^{-1}\at{u^{3}}=\frac{1}{2}\]
    \[(D-1)^{-1}uu_{2}=uu_{2}+3uu_{2}-uu_{2}=3uu_{2}\Rightarrow(D-1)^{-1}\at{uu_{2}}=\frac{1}{3}\]
    \[(D-1)^{-1}u_{4}=5u_{4}-u_{4}=4u_{4}\Rightarrow(D-1)^{-1}\at{u_{4}}=\frac{1}{4}\]
    In definitiva
    \[g_{1}=(D-1)^{-1}\frac{u^{3}}{3}+(D-1)^{-1}\frac{\epsilon^{2}}{8}uu_{2}+(D-1)^{-1}\frac{\epsilon^{4}}{288}u_{4}\]
    \[=\frac{u^{3}}{6}+\frac{\epsilon^{2}}{24}uu_{2}+\frac{\epsilon^{4}}{1152}u_{4}\rightarrow deg(g_{1})=0,2\cdot(-1)+2\cdot1=0,4\cdot(-1)+4\cdot1=0;\]
    
    \item $g_{2}=(D-1)^{-1}\partial_{x}^{-1}\{g_{1},\overline{h}_{KdV}\}$ laddove
    \[\{g_{1},\overline{h}_{KdV}\}=\sum_{s\geq0}\frac{\partial g_{1}}{\partial u_{s}}\partial_{x}^{s}(\partial_{x}\frac{\delta\overline{h}_{KdV}}{\delta u})=\sum_{s\geq0}\frac{\partial}{\partial u_{s}}(\frac{u^{3}}{6}+\frac{\epsilon^{2}}{24}uu_{2}+\frac{\epsilon^{4}}{1152}u_{4})(\partial_{x}^{s}\circ\partial_{x})(\frac{\delta\overline{h}_{KdV}}{\delta u})\]
    \[=(\frac{u^{2}}{2}+\frac{\epsilon^{2}}{24}u_{2})\partial_{x}(\frac{\delta\overline{h}_{KdV}}{\delta u})+\frac{\epsilon^{2}}{24}u\partial_{x}^{3}(\frac{\delta\overline{h}_{KdV}}{\delta u})+\frac{\epsilon^{4}}{1152}\partial_{x}^{5}(\frac{\delta\overline{h}_{KdV}}{\delta u})\]
    \[=(\frac{u^{2}}{2}+\frac{\epsilon^{2}}{24}u_{2})(uu_{1}+\frac{\epsilon^{2}}{12}u_{3})+\frac{\epsilon^{2}}{24}u\partial_{x}^{2}(uu_{1}+\frac{\epsilon^{2}}{12}u_{3})+\frac{\epsilon^{4}}{1152}\partial_{x}^{5}(\frac{\delta\overline{h}_{KdV}}{\delta u})\]
    \[=(\frac{u^{2}}{2}+\frac{\epsilon^{2}}{24}u_{2})(uu_{1}+\frac{\epsilon^{2}}{12}u_{3})+\frac{\epsilon^{2}}{24}u\partial_{x}(u_{1}^{2}+uu_{2}+\frac{\epsilon^{2}}{12}u_{4})+\frac{\epsilon^{4}}{1152}\partial_{x}^{5}(\frac{\delta\overline{h}_{KdV}}{\delta u})\]
    \[=(\frac{u^{2}}{2}+\frac{\epsilon^{2}}{24}u_{2})(uu_{1}+\frac{\epsilon^{2}}{12}u_{3})+\frac{\epsilon^{2}}{24}u(2u_{2}u_{1}+u_{1}u_{2}+uu_{3}+\frac{\epsilon^{2}}{12}u_{5})+\frac{\epsilon^{4}}{1152}\partial_{x}^{5}(\frac{u^{2}}{2}+\frac{\epsilon^{2}}{12}u_{2})\]
    \[\Rightarrow\partial_{x}^{-1}\{g_{1},\overline{h}_{KdV}\}=\partial_{x}^{-1}\bigr[(\frac{u^{3}u_{1}}{2}+\frac{\epsilon^{2}}{24}u^{2}u_{3}+\frac{\epsilon^{2}}{24}uu_{1}u_{2}+\frac{\epsilon^{4}}{288}u_{2}u_{3})+\frac{\epsilon^{2}}{24}(3uu_{1}u_{2}+u^{2}u_{3}+\frac{\epsilon^{2}}{12}uu_{5}) \bigl]\]
    \[+\frac{\epsilon^{4}}{1152}\partial_{x}^{4}(\frac{u^{2}}{2}+\frac{\epsilon^{2}}{12}u_{2}) \]
    \[=\frac{u^{4}}{8}+\frac{\epsilon^{2}}{24}\partial_{x}^{-1}(u^{2}u_{3})+\frac{\epsilon^{2}}{24}\partial_{x}^{-1}(uu_{1}u_{2})+\frac{\epsilon^{4}}{288}\frac{u_{2}^{2}}{2}+\frac{3\epsilon^{2}}{24}\partial_{x}^{-1}(uu_{1}u_{2})+\frac{\epsilon^{2}}{24}\partial_{x}^{-1}(u^{2}u_{3})+\frac{\epsilon^{4}}{288}\partial_{x}^{-1}(uu_{5})\]
    \[+\frac{\epsilon^{4}}{1152}\partial_{x}^{4}(\frac{u^{2}}{2}+\frac{\epsilon^{2}}{12}u_{2})\]
    \[=\frac{u^{4}}{8}+\frac{\epsilon^{2}}{24}(2\partial_{x}(u^{2}u_{3})+4\partial_{x}(uu_{1}u_{2}))+\frac{\epsilon^{4}}{288}\frac{u_{2}^{2}}{2}+\frac{\epsilon^{4}}{288}\partial_{x}^{-1}(uu_{5})+\frac{\epsilon^{4}}{1152}\partial_{x}^{4}(\frac{u^{2}}{2}+\frac{\epsilon^{2}}{12}u_{2})=\clubsuit\]
    Ora sfruttiamo i seguenti integrali
    \[\int u^{2}u_{3}dx=u^{2}u_{2}-2\int uu_{1}u_{2}dx\Rightarrow\int u^{2}u_{3}dx+2\int uu_{1}u_{2}dx=u^{2}u_{2} \]
    \[\int uu_{5}dx=uu_{4}-\int u_{4}u_{1}dx=uu_{4}-(u_{1}u_{3}-\int u_{2}u_{3}dx)=uu_{4}-u_{1}u_{3}+\frac{u_{2}^{2}}{2}\]
    In questa maniera posso proseguire scrivendo
    \[\clubsuit=\frac{u^{4}}{8}+\frac{\epsilon^{2}}{12}u^{2}u_{2}+\frac{\epsilon^{4}}{288}(\frac{u_{2}^{2}}{2}+uu_{4}-u_{1}u_{3}+\frac{u_{2}^{2}}{2})+\frac{\epsilon^{4}}{1152}\partial_{x}^{3}(uu_{1}+\frac{\epsilon^{2}}{12}u_{3})\]
    \[=\frac{u^{4}}{8}+\frac{\epsilon^{2}}{12}u^{2}u_{2}+\frac{\epsilon^{2}}{288}(u_{2}^{2}+uu_{4}-u_{1}u_{3})+\frac{\epsilon^{4}}{1152}\partial_{x}^{2}(u_{1}^{2}+uu_{2}+\frac{\epsilon^{2}}{12}u_{4})\]
    \[=\frac{u^{4}}{8}+\frac{\epsilon^{2}}{12}u^{2}u_{2}+\frac{\epsilon^{2}}{288}(u_{2}^{2}+uu_{4}-u_{1}u_{3})+\frac{\epsilon^{4}}{1152}\partial_{x}(2u_{1}u_{2}+u_{1}u_{2}+uu_{3}+\frac{\epsilon^{2}}{12}u_{5})\]
    \[=\frac{u^{4}}{8}+\frac{\epsilon^{2}}{12}u^{2}u_{2}+\frac{\epsilon^{2}}{288}(u_{2}^{2}+uu_{4}-u_{1}u_{3})+\frac{\epsilon^{4}}{1152}(2u_{1}u_{3}+2u_{2}^{2}+u_{2}^{2}+u_{1}u_{3}+u_{1}u_{3}+uu_{4}+\frac{\epsilon^{2}}{12}u_{6})\]
    \[=\frac{u^{4}}{8}+\frac{\epsilon^{2}}{12}u^{2}u_{2}+\frac{\epsilon^{2}}{288}(u_{2}^{2}+uu_{4}-u_{1}u_{3})+\underbrace{\frac{\epsilon^{4}}{1152}}_{=\frac{\epsilon^{4}}{288\cdot4}}(4u_{1}u_{3}+3u_{2}^{2}+uu_{4}+\frac{\epsilon^{2}}{12}u_{6})\]
    \[=\frac{u^{4}}{8}+\frac{\epsilon^{2}}{12}u^{2}u_{2}+\frac{\epsilon^{4}}{288}(\frac{7}{4}u_{2}^{2}+\frac{5}{4}uu_{4}+\frac{\epsilon^{2}}{48}u_{6})=\frac{u^{4}}{8}+\frac{\epsilon^{2}}{12}u^{2}u_{2}+\frac{\epsilon^{4}}{288}(\frac{7}{4}u_{2}^{2}+\frac{5}{4}uu_{4})+\frac{\epsilon^{6}}{13824}u_{6}\]
    Infine basterà osservare che 
    \[(D-1)u^{4}=4u^{4}-u^{4}=3u^{4}\Rightarrow(D-1)^{-1}\at{u^{4}}=\frac{1}{3}\]
    \[(D-1)u^{2}u_{2}=2u^{2}u_{2}+3u^{2}u_{2}-u^{2}u_{2}=4u^{2}u_{2}\Rightarrow(D-1)^{-1}\at{u^{2}u_{2}}=\frac{1}{4}\]
    \[(D-1)u_{2}^{2}=6u_{2}^{2}-u_{2}^{2}=5u_{2}^{2}\Rightarrow(D-1)^{-1}\at{u_{2}^{2}}=\frac{1}{5}\]
    \[(D-1)uu_{4}=uu_{4}+5uu_{4}-uu_{4}=5uu_{4}\Rightarrow(D-1)^{-1}\at{uu_{4}}=\frac{1}{5}\]
    \[(D-1)u_{6}=7u_{6}-u_{6}=6u_{6}\Rightarrow(D-1)^{-1}\at{u_{6}}=\frac{1}{6}\]
    \[\Rightarrow g_{2}=\frac{u^{4}}{24}+\frac{\epsilon^{2}}{48}u^{2}u_{2}+\epsilon^{4}(\frac{7}{5760}u_{2}^{2}+\frac{uu_{4}}{1152})+\frac{\epsilon^{6}}{82944}u_{6}\]
\end{itemize}



\section{Gerarchie di Gelfand-Dickey}
Questa sezione segue principalmente il percroso di [1].

Quello che abbiamo fatto nella sezione 5 è stato dare un formalismo generale per sistemi integrabili in un contesto a dimensione infinita, vorremmo tuttavia cercare (in analogia a quanto fatto per il caso finito) un metodo e/o criterio che permetta di ricavare (se presenti) delle quantità conservate. Per farlo sarà anzitutto necessario introdurre la rappresentazione di Lax in questo contesto.
Introdurremo delle nuove gerarchie, quelle di Gelfand-Dickey. Tali gerarchie, essendo composte da sistemi integrabili, possiedono degli integrali primi ben definiti, e ciò che faremo in questa sezione sarà costruirli in maniera opportuna.

Si consideri nuovamente l'algebra (commutativa) $\mathcal{A}$ dei polinomi di simboli formali $\{u_{j}^{i}\}_{j\geq0}^{i=1,..,n-2}$ per qualche intero n,e ricordiamo che in tale algebra è definita la derivazione
\[\partial_{x}(fg)=f\partial_{x} g+g\partial_{x} f\hspace{1cm}\forall f,g\in\mathcal{A}\]
dove
\[\partial_{x} u_{j}^{i}:=u_{j+1}^{i}\]
$\forall i,j$; In tal modo abbiamo costruito un algebra differenziale $\mathcal{A}$.

In tale algebra si possono costruire chiaramente altre derivazioni:sia $\xi$ una di queste con $\xi u_{j}^{i}:=a_{i,j}$. Allora per ogni $f\in\mathcal{A}$ psso scrivere
\[\xi f=\sum_{i,k}a_{i,k}\frac{\partial f}{\partial u_{k}^{i}}\]
Le derivazioni che commutano con $\partial_{x}$ giocano un ruolo particolarmente importante, infatti se $[\partial_{x},\xi]=0$ allora
\[a_{i,k+1}=\xi u_{k+1}^{i}=\xi\partial_{x} u_{k}^{i}=\partial_{x}\xi u_{k}^{i}=\partial_{x} a_{i,k}\Rightarrow a_{i,k}=\partial_{x}^{k}a_{i}=a_{k}^{i}\]
\[\Rightarrow\xi f=\sum_{i,k}a_{k}^{i}\frac{\partial f}{\partial u_{k}^{i}}\]

\begin{Def}
Sia $a:=(a^{1},...,a^{n-2})$ con $a^{i}\in\mathcal{A}\hspace{0.2cm}\forall i=1,...,n-2$. Definiamo la derivazione
\begin{equation}
\partial_{a}:=\sum_{i=1}^{n-2}\sum_{k=0}^{\infty}a_{k}^{i}\frac{\partial}{\partial u_{k}^{i}}
\end{equation}
\end{Def}

\begin{Prop}
Le derivazioni $\partial_{x}$ e $\partial_{a}$ commutano.
\end{Prop}
\begin{proof}
Dobbiamo mostrare che, per ogni $f\in\mathcal{A}$, si abbbia 
\[(\partial_{x}\partial_{a}-\partial_{a}\partial_{x})f=0\]
In effetti, per le proprietà dell'algebra, sarà sufficiente dimostrarlo per i generatori di $\mathcal{A}$, cioè per $f=u_{j}^{i}$.

Ma da una parte si ha
\[\partial_{x}\partial_{a}(u_{j}^{i})=\partial_{x}(\sum_{l=1}^{n-2}\sum_{k=0}^{\infty}a^{l}_{k}\frac{\partial u^{i}_{j}}{\partial u_{k}^{l}})=\partial_{x} a_{j}^{i}=a_{j+1}^{i}\]
D'altra parte però
\[\partial_{a}\partial_{x}(u^{i}_{j})=\partial_{a}u^{i}_{j+1}=\sum_{l=1}^{n-2}\sum_{k=0}^{\infty} a_{k}^{l}\frac{\partial u_{j+1}^{i}}{\partial u_{k}^{l}}=a_{j+1}^{i}\]
\end{proof}

\begin{Prop}
La derivazione (58) può essere applicata ad un funzionale $\overline{f}\in\widehat{\Lambda}$ mediante l'uguaglianza 
\begin{equation}
    \partial_{a}\overline{f}=\int\partial_{a}fdx
\end{equation}
\end{Prop}
\begin{proof}
Per la Prop.7.1 sappiamo che $\partial_{a}$ e $\partial_{x}$ commutano, dunque se $f,g\in\mathcal{A}$ differiscono per un certo $\partial_{x} h(h\in\mathcal{A})$ si ha 
\[f-g=\partial_{x} h\Rightarrow \partial_{a}f-\partial_{a}g=\partial_{a}(f-g)=\partial_{a}\partial_{x} h=\partial_{x}(\partial_{a}h)\]
Ma allora $\partial_{a}f,\partial_{a}g$ risiedono nella stessa classe di equivalenza. Chiaramente il discorso si estende in maniera ovvia al caso in cui $f,g$ differiscano anche per una costante $c\in\mathbb{C}[[\epsilon]]$.
\end{proof}

\begin{Prop}
Vale la seguente relazione
\begin{equation}
    \partial_{a}\overline{f}=\int (\sum_{i=1}^{n-2}\frac{\delta f}{\delta u^{i}}a^{i})dx
\end{equation}
dove si ricordi la "derivata variazionale" $\frac{\delta f}{\delta u^{i}}:=\sum_{k=0}^{\infty}(-\partial_{x})^{k}\frac{\partial f}{\partial u_{k}^{i}}$
\end{Prop}
\begin{proof}
Si tratta di un calcolo diretto
\[\partial_{a}\overline{f}=\int\partial_{a}fdx=\int\sum_{i}\sum_{k=0}^{\infty}a_{k}^{i}\frac{\partial f}{\partial u_{k}^{i}} dx\stackrel{(53)}{=}\int\sum_{i}a^{i}\sum_{k=0}^{\infty}(-1)^{k}\partial_{x}^{k}(\frac{\partial f}{\partial u_{k}^{i}})dx=\int\sum_{i}\frac{\delta f}{\delta u^{i}}a^{i}dx\]
\end{proof}
$\textbf{NOTA}$:Dalla Proposizione precedente segue che, se $f=\partial_{x} g$, allora deve essere $\frac{\delta f}{\delta u^{i}}=0\forall i=1,...,n-2$. Questo perchè il primo membro  della (60) diventa nullo, ed il secondo dovrà annullarsi per ogni scelta di $a$.



\subsection{Anello di Pseudooperatori differenziali}

\begin{Def}
Sia $m\in\mathbb{Z}$ arbitrario, definiamo la serie formale
\[X=\sum_{-\infty}^{m}X_{i}\partial_{x}^{i}\hspace{0.5cm}X_{i}\in\mathcal{A}\]
Tale serie può essere aggiunta e moltiplicata per elementi di $\mathcal{A}$ (nonchè per ulteriori serie formali) una volta definita la regola di commutazione
\[\partial_{x}^{k}f:=f\partial_{x}^{k}+\binom{k}{1}f^{'}\partial_{x}^{k-1}+...=\sum_{i=0}^{\infty}\binom{k}{i}f^{(i)}\partial_{x}^{k-i}\]
dove, per evitare abusi di notazione $f^{(i)}\equiv\partial_{x}^{i}(f)$ è la derivazione $\partial_{x}$ applicata $i$ volte ad $f$, e ricordiamo che
\[\binom{k}{i}:=\frac{(k)(k-1)...(k-i+1)}{i!}=\frac{k!}{(k-i)!i!}\bigg(=0\hspace{0.2cm}se\hspace{0.2cm} i>k\geq0\bigg)\]
L'insieme di tutte le serie formali così definite viene chiamato "anello degli operatori pseudo-differenziali"($\Psi DO$). Adotteremo poi la seguente notazione
\[R_{+}:=\{\sum_{i\geq0}X_{i}\partial_{x}^{i}\}\hspace{0.5cm}R_{-}:=\{\sum_{i<0}X_{i}\partial_{x}^{i}\}\]
il primo detto "anello di operatori differenziali", il secondo detto "anello di operatori integrali".
Invece indicheremo con $X_{+}:=\sum_{i\geq0}X_{i}\partial_{x}^{i},X_{-}:=\sum_{i<0}X_{i}\partial_{x}^{i}$ ed infine $res(X):=X_{-1}$.
Infine scriveremo
\[R_{n}:=\{\sum_{i=0}^{n-1}X_{i}\partial_{x}^{i}\}\hspace{0.5cm}\partial_{x}^{n}R_{-}:=\{\sum_{-\infty}^{n-1}X_{i}\partial_{x}^{i}\}\]
\end{Def}

\begin{Def}
Definiamo la mappa
\[\langle\cdot,\cdot\rangle:\mathcal{A}\times\mathcal{A}\mapsto\widehat{\Lambda}\]
\[\langle X,Y\rangle:=\int res(XY)dx\in\widehat{\Lambda} \]
\end{Def}
In particolare $\langle R_{+},R_{+}\rangle=\langle R_{-}, R_{-}\rangle=0$. 

\begin{Prop}
La moltiplicazione definita nella Definizione 7.2 è associativa.
\end{Prop}
\begin{proof}
Per linearità è sufficiente mostrare che, dati
\[(A\partial_{x}^{a}B\partial_{x}^{b})C\partial_{x}^{c}=A\partial_{x}^{a}(B\partial_{x}^{b}C\partial_{x}^{c})\hspace{0.5cm}a,b,c\in\mathbb{Z};A,B,C\in\Psi DO\]
A primo membro si ha
\[(\sum_{i\geq0}\binom{a}{i}AB^{(i)}\partial_{x}^{a+b-i})C\partial_{x}^{c}=\sum_{i,j\geq0}\binom{a}{i}\binom{a+b-i}{j}AB^{(i)}C^{(j)}\partial_{x}^{a+b+c-i-j}\]
A secondo membro invece ho
\[A\partial_{x}^{a}(\sum_{l\geq0}\binom{b}{l}BC^{(l)}\partial_{x}^{c+b-l})=\sum_{k,l\geq0}\binom{b}{l}\binom{a}{k}A(BC^{(l)})^{(k)}\partial_{x}^{a+b+c-k-l}\]
\[=\sum_{l,k\geq0}\sum_{\nu=0}^{k}\binom{b}{l}\binom{a}{k}\binom{k}{\nu}AB^{(k-\nu)}C^{(l+\nu)}\partial_{x}^{a+b+c-k-l}\hspace{0.5cm}\begin{cases} k-\nu=:i\\ \nu=0\iff i=k,\nu=k\iff i=0\end{cases}\]
\[=\sum_{l,k\geq0}\sum_{i=0}^{k}\binom{a}{k}\binom{b}{l}\binom{k}{k-i}AB^{(i)}C^{(l+k-i)}\partial_{x}^{a+b+c-l-k}\hspace{0.5cm}\begin{cases} l+k-i=:j\\ l=0\iff j=k-i   \end{cases}\]
\[=\sum_{k\geq0}\sum_{i=0}^{k}\sum_{j\geq k-i}\binom{a}{k}\binom{b}{i+j-k}\binom{k}{i}AB^{(i)}C^{(j)}\partial_{x}^{a+b+c-i-j}\]
Ora, si ha $j\geq k-i\Rightarrow i\leq k\leq i+j$ mentre $i,j\geq0$. Rimango allora con 
\[\sum_{i,j\geq0}\sum_{k=i}^{i+j}\binom{a}{k}\binom{b}{i+j-k}\binom{k}{i}AB^{(i)}C^{(j)}\partial_{x}^{a+b+c-i-j}\]
Resta allora da verificare se
\[\sum_{k=i}^{i+j}\binom{a}{k}\binom{b}{i+j-k}\binom{k}{i}=\binom{a}{i}\binom{a+b-i}{j}\]
Tuttavia (vedi Appendice)
\[\sum_{k=i}^{i+j}\binom{a}{k}\binom{b}{i+j-k}\binom{k}{i}=\sum_{k=i}^{i+j}\binom{a}{i}\binom{a-i}{k-i}\binom{b}{i+j-k}\]
Rimane allora da verificare la relazione
\[\sum_{k=i}^{i+j}\binom{a-i}{k-i}\binom{b}{i+j-k}=\binom{a+b-i}{j}\]
Ma quest'ultima è sempre verificata, dal momento che vale la "convoluzione di Vandermonde" (vedasi appendice per una dimostrazione)
\[\sum_{r=0}^{k}\binom{n}{r}\binom{m}{k-r}=\binom{n+m}{k}\hspace{0.5cm}\forall n,m\in\mathbb{N}\]
Basta allora sostituire $r$ con $k-i$,$n$ con $a-i$,$m$ con $b$ e $k$ con $j$.
\end{proof}

\begin{Prop}
Per ogni serie formale X e Y esiste un $h\in\mathcal{A}$ tale per cui
\begin{equation}
    res[X,Y]=\partial_{x} h
\end{equation}
\end{Prop}
\begin{proof}
Per linearità basterà verificarlo per i monomi $X=X_{i}\partial_{x}^{i}$ ed $Y=Y_{j}\partial_{x}^{j}$, con $i\geq0,j<0,i+j\geq-1$ dal momento che negli altri casi basterà scambiare i con j(si ricordi che $\langle R_{-},R_{-}\rangle=\langle R_{+},R_{+}\rangle=0$). Verifichiamolo usando
\[h=\binom{i}{i+j+1}\sum_{\alpha=0}^{i+j}(-1)^{\alpha}X_{i}^{(\alpha)}Y_{j}^{(i+j-\alpha)}\]
In effetti
\[[X,Y]=X_{i}\partial_{x}^{i}(Y_{j})\partial_{x}^{j}-Y_{j}\partial_{x}^{j}(X_{i})\partial_{x}^{i}=X_{i}\sum_{k=0}^{\infty}\binom{i}{k}Y_{j}^{(k)}\partial_{x}^{i+j-k}-Y_{j}\sum_{k=0}^{\infty}\binom{j}{k}X_{i}^{(k)}\partial_{x}^{j+i-k}\]
\[\Rightarrow res[X,Y]=X_{i}\binom{i}{i+j+1}Y_{j}^{(i+j+1)}-Y_{j}\binom{j}{i+j+1}X_{i}^{(i+j+1)}\]
ma 
\[\binom{j}{i+j+1}=\frac{j(j-1)...(j-i-j-1+1)}{(i+j+1)!}=\frac{j(j-1)...(-i)}{(i+j+1)!}=(-1)^{i+j+1}\frac{i(i-1)...(-j)}{(i+j+1)!}\]
\[=(-1)^{i+j+1}\binom{i}{i+j+1}\]
perciò
\[res[X,Y]=\binom{i}{i+j+1}(X_{i}Y_{j}^{(i+j+1)}+(-1)^{i+j}Y_{j}X_{i}^{(i+j+1)})=\partial_{x} h\]
In particolare, nell'ultimo passaggio, abbiamo utilizzato il fatto che
\[\partial_{x} h=\binom{i}{i+j+1}\sum_{\alpha=0}^{i+j}(-1)^{\alpha}(X_{i}^{(\alpha)}Y_{j}^{(i+j-\alpha+1)}+X_{i}^{(\alpha+1)}Y_{j}^{(i+j-\alpha)})=\]
\[=\binom{i}{i+j+1}(X_{i}Y_{j}^{(i+j+1)}+X_{i}^{(1)}Y_{j}^{(i+j)}-X_{i}^{(1)}Y_{j}^{(i+j)}-X_{i}^{(2)}Y_{j}^{(i+j-1)}+...\]
\[...+(-1)^{i+j}X_{i}^{(i+j)}Y_{j}^{(1)}+(-1)^{i+j}X_{i}^{(i+j+1)}Y_{j})=res[X,Y]\]
poichè i termini non estremali si annullano vicendevolente.
\end{proof}


\begin{Cor}
\begin{equation}
\int res(XY)dx=\int res(YX)dx
\end{equation}
\end{Cor}

\begin{Prop}
Se $X=\sum_{-\infty}^{m}X_{i}\partial_{x}^{i}$ ed $X_{m}=1$, allora esiste un unico $\Psi DO$ inverso $X^{-1}$ ed un unico $X^{1/m}$. Essi commutano con X.
\end{Prop}
\begin{proof}
Sia 
\[X^{-1}=\partial_{x}^{-m}+Y_{-m-1}\partial_{x}^{-m-1}+...\]
a coefficenti indefiniti. Imponendo $XX^{-1}=1$ ottengo
\[\sum_{i=-\infty}^{m}X_{i}\partial_{x}^{i}(\partial_{x}^{-m}+Y_{-m-1}\partial_{x}^{-m-1}+...)=\overbrace{X_{m}}^{=1}(1+Y_{-m-1}\partial_{x}^{-1}+\binom{m}{1}Y_{-m-1}^{'}\partial_{x}^{-2}+...\]
\[+Y_{-m-2}\partial_{x}^{-2}+...)+X_{m-1}(\partial_{x}^{-1}+Y_{-m-1}\partial_{x}^{-2}+...+Y_{-m-2}\partial_{x}^{-3}+...)+X_{m-2}(\partial_{x}^{-2}+...)...\]
\[...=1+(X_{m-1}+Y_{-m-1})\partial_{x}^{-1}+(X_{m-2}+Y_{-m-2}+Y_{-m-1}X_{m-1}+mY^{'}_{-m-1})\partial_{x}^{-2}+...\equiv 1\]
Otteniamo così una sequenza di equazioni della forma $Y_{-m-k}=-X_{m-k}+Q_{k}$ dove $Q_{k}$ sono polinomi differenziali in $\{X_{i}\},\{Y_{j}\}(j>-m-k)$. Allo stesso modo si può costruire $X^{1/m}$ definendolo come $(X^{1/m})^{m}=X$. Segue che $[X,X^{-1}]=1-1=0$, mentre
\[[X,X^{1/m}]=\overbrace{X^{1/m}...X^{1/m}}^{m-volte}X^{1/m}-X^{1/m}\overbrace{X^{1/m}...X^{1/m}}^{m-volte}=0\]
\end{proof}

\begin{Cor}
Per ogni intero p l'operatore $X^{\frac{p}{m}}$ può essere costruito in modo che il suo termine maggiore sia $\partial_{x}^{p}$ e che commuti con X. 
\end{Cor}

\begin{Def}
Se $X=\sum_{i=0}^{n-2}X_{i}\partial_{x}^{i}\in R_{n-1}$, allora denoteremo con $\partial_{X}$ la derivazione (57) corrispondente ad $a\equiv(X_{0},...,X_{n-2})$. 

Definiamo poi
\[\frac{\delta f}{\delta L}:=\sum_{i=0}^{n-2}\partial_{x}^{-i-1}\frac{\delta f}{\delta u^{i}}\in\Psi DO\]
Tale operatore possiede in particolare un numero infinito indici negativi, dal momento che le potenze negative di $\partial_{x}$ vanno a creare infiniti termini (si veda la Definizione 7.2).
\end{Def}

\begin{Prop}
Se $X=\sum_{i=0}^{n-2}X_{i}\partial_{x}^{i}\in R_{n-1}$, allora
\begin{equation}
    \partial_{X}\overline{f}=\int res(X\cdot\frac{\delta f}{\delta L})dx=\langle X,\frac{\delta f}{\delta L}\rangle
\end{equation}
\end{Prop}
\begin{proof}
Abbiamo
\[\partial_{X}\overline{f}\stackrel{(59)}{=}\int\partial_{X}f dx=\int\sum_{i=0}^{n-2}\sum_{k=0}^{\infty}X_{i}^{(k)}\frac{\partial f}{\partial u^{i}_{k}} dx\stackrel{(58)}{=}\int\sum_{i=0}^{n-2}X_{i}\frac{\delta f}{\delta u^{i}}dx=\int res(X\frac{\delta f}{\delta L})dx\]
In particolare si è sfruttato il fatto che
\[res(X\frac{\delta f}{\delta L})=res(\sum_{i}^{n-2}X_{i}\partial_{x}^{i}\sum_{j=0}^{n-2}\partial_{x}^{-j-1}\frac{\delta f}{\delta u^{j}})=res(\sum_{i}^{n-2}X_{i}\partial_{x}^{i}\sum_{j=0}^{n-2}\sum_{\alpha=0}^{\infty}\binom{-j-1}{\alpha}\frac{\delta f}{\delta u^{j}}^{(\alpha)}\partial_{x}^{-j-1-\alpha})\]
\[=res(\sum_{i,j=0}^{n-2}\sum_{\alpha,\beta=0}^{\infty}\binom{-j-1}{\alpha}\binom{i}{\beta}X_{i}\frac{\delta f}{\delta u^{j}}^{(\beta+\alpha)}\partial_{x}^{-j-1+i-\alpha-\beta})\stackrel{i=j;\alpha=-\beta=0}{=}\sum_{i=0}^{n-2}X_{i}\frac{\delta f}{\delta u^{i}}\]
\end{proof}

\begin{Lem}
Si ha $\partial_{X}L=X$, dove L è definito dalla (64).
\end{Lem}
\begin{proof}
Da un calcolo diretto ottengo
\[\partial_{X}L=\sum_{i=0}^{n-2}\partial_{X}u^{i}\partial_{x}^{i}=\sum_{i=0}^{n-2}\sum_{j=0}^{n-2}\sum_{k=0}^{\infty}X_{j}^{(k)}\frac{\partial u^{i}}{\partial u_{k}^{j}}\partial_{x}^{i}=\sum_{i=0}^{n-2}X_{i}\partial_{x}^{i}=X\]
\end{proof}



\subsection{Coppie di Lax, gerarchie di Gelfand-Dickey}
Vogliamo generalizzare i concetti visti nella sezione 4

\begin{Def}
Una volta definito l'operatore (pseudo-)differenziale
\begin{equation}
    L:=\partial_{x}^{n}+u^{n-2}\partial_{x}^{n-2}+...+u^{0}
\end{equation}
dove $\{u^{i}\}^{i=0,...,n-2}\in\mathcal{A}$, si definisce il nuovo operatore
\[P_{m}:=(L^{m/n})_{+}\equiv L_{+}^{m/n}\]
\end{Def}

\begin{Prop}
Il commutatore $[P_{m},L]$ appartiene ad $R_{n-1}$.
\end{Prop}
\begin{proof}
Infatti
\[[P_{m},L]=[L^{m/n}-L_{-}^{m/n},L]\stackrel{Prop.  7.6}{=}-[L_{-}^{m/n},L]\]
L'operatore a membro di sinistra è ovviamente in $R_{+}$(questo perchè sia $P_{m}$ che $L$ possiedono potenze positive di $\partial_{x}$, dunque moltiplicandoli tra loro si ottengono delle serie con finite con potenze $\geq0$ vista la definizione 7.2); mentre quello di destra è un commutatore di due operatori di ordine -1($L_{-}^{\frac{m}{n}}$) ed n($L$), dunque il suo ordine sarà inferiore o uguale ad $-1+n-1=n-2$.
\end{proof}

\begin{Def}
Diremo che l'operatore differenziale P, assieme ad L, definiscono una "coppia di Lax" se $[P,L]\in R_{n-1}$.
\end{Def}

\begin{flushleft}
$\textbf{NOTA}$:Dalla Proposizione 7.8 segue che, per ogni intero $m>0$, l'operatore $P_{m}$ costituisce una coppia di Lax assieme ad $L$. Dal momento che $[P_{m},L]\in R_{n-1}$, l'operatore $\partial_{[P_{m},L]}$ è ben definito; in particolare il Lemma 7.1 afferma che $\partial_{[P_{m},L]}L=[P_{m},L]$.
\end{flushleft}

\begin{Def}
Sia ora $t_{m}$ un ulteriore parametro, e si supponga che le $u^{i}$ dipendano anche da quest'ultimo; la "gerarchia di Gelfand Dickey(GD)" è quel sistema di equazioni 
\begin{equation}
    \partial_{t_{m}}L\equiv\partial_{m}L=\partial_{[P_{m},L]}L
\end{equation}
o, equivalentemente (vedi Lemma 7.1)
\[\partial_{m}L=[P_{m},L]\]
associate alle $\{u^{i}\}^{i=0,...,n-2}$ e determinato da due interi:n ed m. Se n=2 si parla invece di "gerarchia di Korteweg-de Vries(KdV)".
\end{Def}

\begin{flushleft}
Si osservi che allora, per il Lemma 7.1, possiamo scrivere
\begin{equation}
    \partial_{m}=\partial_{[P_{m},L]}\Rightarrow\partial_{m}=\partial_{\dot{L}}
\end{equation}
\end{flushleft}

\begin{flushleft}
$\underline{Esempio}$:Si consideri n=2, $L=\partial_{x}^{2}+u$, in tal caso poniamo
\[L^{1/2}=\partial_{x}+v_{0}+v_{-1}\partial_{x}^{-1}+v_{-2}\partial_{x}^{-2}+...\]
Ponendo $L=L^{1/2}L^{1/2}$ vedo che
\[\partial_{x}^{2}+u=L^{1/2}L^{1/2}=(\partial_{x}+v_{0}+v_{-1}\partial_{x}^{-1}+v_{-2}\partial_{x}^{-2}+...)(\partial_{x}+v_{0}+v_{-1}\partial_{x}^{-1}+v_{-2}\partial_{x}^{-2}+...)=\]
\[=\partial_{x}^{2}+v_{0}\partial_{x}+v_{0}'+v_{-1}+v_{-1}'\partial_{x}^{-1}+v_{-2}\partial_{x}^{-1}+v_{-2}'\partial_{x}^{-2}+...+v_{0}\partial_{x}+v_{0}^{2}+v_{0}v_{-1}\partial_{x}^{-1}+v_{0}v_{-2}\partial_{x}^{-2}+...+\]
\[+v_{-1}+v_{-1}v_{0}\partial_{x}^{-1}+v_{-1}v_{0}'\partial_{x}^{-2}+...+v_{-2}\partial_{x}^{-1}+v_{-2}v_{0}\partial_{x}^{-2}+...\]
\[\iff\begin{cases}
v_{0}=0\\
2v_{-1}+\cancel{v_{0}^{2}}+\cancel{v_{0}'}=u\Rightarrow v_{-1}=\frac{u}{2}\\
v_{-1}'+2v_{-2}+2\cancel{v_{0}}v_{-1}=0\Rightarrow v_{-2}=-\frac{v_{-1}}{2}=-\frac{u'}{4}\\
v_{n\leq-3}=0
\end{cases}\]
Segue che $L^{1/2}=\partial_{x}+\frac{u}{2}\partial_{x}^{-1}-\frac{u'}{4}\partial_{x}^{-2}$ e dunque
\[P_{3}=L^{3/2}_{+}=\bigg((\partial_{x}+\frac{u}{2}\partial_{x}^{-1}-\frac{u'}{4}\partial_{x}^{-2})(\partial_{x}^{2}+u)\bigg)_{+}=\bigg(\partial_{x}^{3}+u\partial_{x}+u'+\frac{u}{2}\partial_{x}+\frac{u^{2}}{2}\partial_{x}^{-1}+o(\partial_{x}^{-1})-\frac{u'}{4}-\frac{u'u}{4}\partial_{x}^{-2}+\]
\[+o(\partial_{x}^{-2})\bigg)_{+}=\partial_{x}^{3}+\frac{3}{2}u\partial_{x}+\frac{3}{4}u'\]
\[\Rightarrow [P_{3},L]=(\partial_{x}^{3}+\frac{3}{2}u\partial_{x}+\frac{3}{4}u')(\partial_{x}^{2}+u)-(\partial_{x}^{2}+u)(\partial_{x}^{3}+\frac{3}{2}u\partial_{x}+\frac{3}{4}u')\]
\[=(\partial_{x}^{5}+u\partial_{x}^{3}+\binom{3}{1}u'\partial_{x}^{2}+\binom{3}{2}u''\partial_{x}+\binom{3}{3}u^{'''})+(\frac{3}{2}u\partial_{x}^{3}+\frac{3}{2}u(u\partial_{x}+u'))+\]
\[+(\frac{3}{4}u'\partial_{x}^{2}+\frac{3}{4}uu')-\{(\partial_{x}^{5}+\frac{3}{2}u\partial_{x}^{3}+\binom{2}{1}\frac{3}{2}u'\partial_{x}^{2}+\binom{2}{2}\frac{3}{2}u''\partial_{x}+\frac{3}{4}u'\partial_{x}^{2}+\]
\[+\binom{2}{1}\frac{3}{4}u''\partial_{x}+\binom{2}{2}\frac{3}{4}u''')+(u\partial_{x}^{3}+\frac{3}{2}u^{2}\partial_{x}+\frac{3}{4}u'u)\}=\frac{u'''}{4}+\frac{3}{2}uu'\]
Si conclude che
\[\partial_{m}L\equiv\partial_{t}(\partial_{x}^{2}+u)=\frac{u'''}{4}+\frac{3}{2}uu'\iff 4\partial_{t}u=u'''+6uu'\]
che è proprio l'equazione di Korteweg-de Vries.
\end{flushleft}



\subsection{Integrali primi}
Vogliamo adesso costruire degli integrali primi per ogni sistema associato alle gerarchie GD definite nella sezione precedente. Anzitutto bisogna capire cosa si intende con "integrale primo"

\begin{Def}
Un "integrale primo" è un funzionale $\overline{f}=\int fdx$ che viene conservato, ossia che soddisfi
\[0=\partial_{m}\overline{f}=\int\partial_{m}f dx\stackrel{(\star)}{=}\int res([P_{m},L]\frac{\delta f}{\delta L}) dx\]
\end{Def}

\begin{flushleft}
$\textbf{NOTA}$:in $(\star)$ abbiamo sfruttato il fatto che, anzitutto $(\partial_{m}u^{i}_{k}=\dot{u}^{i}_{k}+(66))$
\[\partial_{m}f=\sum_{i=0}^{n-2}\sum_{k=0}^{\infty}\frac{\partial f}{\partial u_{k}^{i}}\dot{u}_{k}^{i}=\sum_{i=0}^{n-2}\sum_{k=0}^{\infty}\frac{\partial f}{\partial u_{k}^{i}}(\dot{u}^{i})^{(k)}\]
\[\Rightarrow\int\partial_{m}fdx=\int\sum_{i=0}^{n-2}\sum_{k=0}^{\infty}\frac{\partial f}{\partial u_{k}^{i}}(\dot{u}^{i})^{(k)}dx=\int\sum_{i=0}^{n-2}\sum_{k=0}^{\infty}(-\partial)^{k}\frac{\partial f}{\partial u_{k}^{i}}\dot{u}^{i}dx=\int\sum_{i=0}^{n-2}\frac{\delta f}{\delta u^{i}}\dot{u}^{i}dx\]
Infine si noti che (in modo analogo alla Prop. 7.7)
\[res\bigg(\dot{L}\frac{\delta f}{\delta L} \bigg)=res(\sum_{j=0}^{n-2}\dot{u}^{j}\partial_{x}^{j}\sum_{i=0}^{n-2}\partial_{x}^{-i-1}\frac{\delta f}{\delta u^{i}})=\sum_{i=0}^{n-2}\dot{u}^{i}\frac{\delta f}{\delta u^{i}}\]
\[\Rightarrow\int\sum_{i=0}^{n-2}\frac{\delta f}{\delta u^{i}}\dot{u}^{i}dx=\int res(\dot{L}\frac{\delta f}{\delta L})dx=\int res([P_{m},L]\frac{\delta f}{\delta L})dx\]
\end{flushleft}

\begin{Lem}
Per ogni $k\in\mathbb{Z}$, l'equazione (65) implica che
\begin{equation}
    \partial_{m}L^{k/n}=[P_{m},L^{k/n}]
\end{equation}
\end{Lem}
\begin{proof}
Sia k=1. Supponiamo che $L^{1/n}=\partial_{x}+v_{-1}\partial_{x}^{-1}+...$ con $v_{i}$ polinomi differenziali in $\{u^{i}\}$. Se definiamo $\partial_{m}v_{i}$ in modo arbitrario e poi mostriamo che $\partial_{m}u^{i}$ risulta in accordo con $\partial_{m}L=[P_{m},L]$, significa che $\partial_{m}v_{i}$ era stato scelto correttamente. 

Supponiamo allora che $\partial_{m}L^{1/n}=[P_{m},L^{1/n}]$ e verfichiamo:
\[\partial_{m}L=\partial_{m}(L^{1/n})^{n}\stackrel{Leibniz}{=}\sum_{i=0}^{n-1}(L^{1/n})^{i}\partial_{m}L^{1/n}(L^{1/n})^{n-1-i}=\sum_{i=0}^{n-1}(L^{1/n})^{i}[P_{m},L^{1/n}]L^{n-1-i}\]
\[=[P_{m},(L^{1/n})^{n}]=[P_{m},L]\]
Per $k>1$ si usa lo stesso ragionamento, mentre per k negativi si consideri $k>0$, la tesi segue ora dal fatto $\partial_{m}$ è una derivazione:
\[0=\partial_{m}(L^{-k/n}L^{k/n})=(\partial_{m}L^{-k/n})L^{k/n}+L^{-k/n}\partial_{m}L^{k/n}\]
\[\Rightarrow\partial_{m}L^{-k/n}=-L^{-k/n}(\partial_{m}L^{k/n})L^{-k/n}\]
Dunque
\[\partial_{m}L^{-k/n}=-L^{-k/n}[P_{m},L^{k/n}]L^{-k/n}=-(L^{-k/n}P_{m}-P_{m}L^{-k/n})=-[L^{-k/n},P_{m}]\]
\[=[P_{m},L^{-k/n}]\]
\end{proof}

Terminiamo la sezione con il risultato che ci ervavamo proposti di trovare

\begin{Prop}
I funzionali
\begin{equation}
    J_{k}:=\int res(L^{k/n})dx\hspace{0.5cm}k=1,2,...
\end{equation}
sono integrali primi di tutte le equzioni della gerarchia n-esima. 
\end{Prop}
\begin{proof}
Infatti
\[\partial_{m}J_{k}=\int\partial_{m}res(L^{k/n})dx=\int res(\partial_{m}L^{k/n})dx=\int res([P_{m},L^{k/n}])dx=0\]
grazie alla Prop. 7.5.
\end{proof}

Risulta così evidente la proprietà più importante delle gerarchie GD: esse ammettono infiniti integrali primi.



\subsection{Compatibilità delle equazioni di Gelfand-Dickey}
Un ulteriore peculitarità delle equazioni $(65)$ è la loro "compatibilità" per diversi valori di $m$:ciò significa che è possibile trovare delle soluzioni $\{u^{\alpha}(x,t_{m})\}_{m\in\mathbb{N}}$ che soddisfino le equazioni della gerarchia GD per ogni valore di $m$ contemporaneamente. Si può poi dimostrare che, affinchè ciò avvenga, è necessario che $\partial_{m}$ e $\partial_{l}$ commutino per ogni scelta di $m,l$. Si noti anche che, per quanto riguarda i sistemi Hamiltoniani, tale condizione è equivalente proprio alla definizione di integrabilità!
Rimandiamo al testo di Buryak per ulteriori approfondimenti [17].

\begin{Prop}
Vale la seguente uguaglianza
\[\partial_{l}P_{m}-\partial_{m}P_{l}=[P_{l},P_{m}]\]
\end{Prop}
\begin{proof}
Utilizzando il Lemma 7.2 avremo che (indicando con $\oplus$,$\ominus$ rispettivamente gli indici $+,-$ indicando che possono essere tralasciati)
\[\overbrace{\partial_{l}L^{m/n}_{+}}^{=(\partial_{l}L^{m/n})_{+}}-\overbrace{\partial_{m}L_{+}^{l/n}}^{=(\partial_{m}L^{l/n})_{+}}\stackrel{Lem. 7.2}{=}[L_{+}^{l/n},\underbrace{L^{m/n}}_{=L^{m/n}_{+}+L^{m/n}_{-}}]_{+}-[L_{+}^{m/n},L^{l/n}]_{+}=\]
\[=[L_{+}^{l/n},L_{+}^{m/n}]_{\oplus}+[L_{\oplus}^{l/n},L_{-}^{m/n}]_{+}-\underbrace{[L_{+}^{m/n},L^{l/n}]_{+}}_{=-[L^{l/n},L_{+}^{m/n}]_{+}}\]
\[=[L_{+}^{l/n},L_{+}^{m/n}]+\underbrace{[L^{l/n},L^{m/n}]}_{=0}_{+}=[L_{+}^{l/n},L_{+}^{m/n}].\]
\end{proof}

\begin{Prop}
Si ha che $[\partial_{l},\partial_{m}]=0$ per ogni valore di m ed l.
\end{Prop}
\begin{proof}
Dal momento che $\partial_{m,l}$ sono derivazioni in $\mathcal{A}$, basta dimostrarlo per i generatori di tale algebra (ossia basta verificarlo sull'azione in L). Si ha allora
\[[\partial_{l},\partial_{m}]L=\partial_{l}\partial_{m}L-\partial_{m}\partial_{l}L\stackrel{Lem. 7.2}{=}\partial_{l}[L^{m/n}_{+},L]-\partial_{m}[L_{+}^{l/n},L]\stackrel{Leibniz}{=}\]
\[=[\partial_{l}L_{+}^{m/n}-\partial_{m}L_{+}^{l/n},L]+[L_{+}^{m/n},\partial_{l}L]-[L_{+}^{l/n},\partial_{m}L]\]
che, sfruttando il Lemma precedente e l'equazione di Lax (65), diventa
\[=[[L_{+}^{l/n},L_{+}^{m/n}],L]+[L_{+}^{m/n},[L_{+}^{l/n},L]]-[L_{+}^{l/n},[L_{+}^{m/n},L]]=\]
\[-([L,[L_{+}^{l/n},L_{+}^{m/n}]]+[L_{+}^{m/n},[L,L_{+}^{l/n}]]+[L_{+}^{l/n},[L_{+}^{m/n},L]])=0\]
grazie all'identità di Jacobi.
\end{proof}

$\textbf{NOTA}$:Per quanto visto possiamo affermare che KdV è a tutti gli effetti un sistema integrabile!



\subsection{Solitoni come soluzioni delle gerarchie GD}
Una proprietà dell'equazione (65) è quella di possedere soluzioni analitiche,  i solitoni.

Siano $n,m\in\mathbb{N}$ fissati, $N$ il numero di solitoni (arbitrario) e si considerino le funzioni
\begin{equation}
    y_{k}(x,t):=\sum_{\epsilon}a_{k,\epsilon}\exp\bigg(\epsilon\alpha_{k}x+\epsilon^{m}\alpha_{k}^{m}t\bigg)\hspace{0.5cm}k=1,...,N
\end{equation}
dove $\epsilon$ è per definizione la radice dell'unità di $n$ ($\epsilon^{n}=1$). Invece le costanti $\alpha_{k},a_{k}$ sono numeri complessi. 

\begin{Def}
Definiamo l'operatore di ordine N 
\[\Phi:=\frac{1}{\Delta}Det\begin{pmatrix}
                            y_{1}&...&y_{N}&1\\
                            y_{1}'&...&y_{N}^{'}&\partial_{x}\\
                            .&&.\\
                            .&&.\\
                            .&&.\\
                            y_{1}^{(N)}&...&y_{N}^{(N)}&\partial_{x}^{N}
                                                         \end{pmatrix}\]
dove, nello sviluppo del determinante, gli operatori nell'ultima colonna non agiscono sulle altre colonne e $\Delta$ è il Wronskiano di ${y_{1},...,y_{N}}$.
\end{Def}

\begin{Lem}
Le funzioni $y_{k}(x,t)$ così costruite soddisfano
\[(\partial_{m}=\partial_{t})y_{k}=\partial_{x}^{m}y_{k}\hspace{0.5cm}\partial_{x}^{n}y_{k}=\alpha^{n}y_{k}\hspace{0.5cm}\Phi y_{k}=0\]
\end{Lem}
\begin{proof}
La prima è ovvia, la seconda anche dal momento che $\epsilon^{n}=1$, mentre la terza è vera poichè (non sto più considerando $\Phi$ come un $\Psi DO$ bensì piuttosto come una combinazione di derivazioni)
\[\Phi y_{k}=\frac{1}{\Delta}Det\begin{pmatrix}
                            y_{1}&...&y_{N}&y_{k}\\
                            y_{1}'&...&y_{N}^{'}&y_{k}'\\
                            .&&.\\
                            .&&.\\
                            .&&.\\
                            y_{1}^{(N)}&...&y_{N}^{(N)}&y_{k}^{(N)}
                                                         \end{pmatrix}=0\]
per le proprietà del determinante.
\end{proof}

Definiamo a questo punto l'operatore
\begin{equation}
    L:=\Phi\partial_{x}^{n}\Phi^{-1}
\end{equation}
Tale operatore soddisfa l'equaione (65), come dimostra la seguente

\begin{Prop}
Valgono le seguenti affermazioni
\begin{itemize}
    \item [(i)]L è un $\Psi DO$ il cui termine a grado maggiore è $\partial_{x}^{n}$;
    \item [(ii)]L soddisfa l'equazione $(65)$.
\end{itemize}
\end{Prop}
\begin{proof}
(i) Riscriviamo (70) come ($\partial_{x}^{N}$ commuta con $\partial_{x}^{n}$)
\[L=(\Phi\partial_{x}^{-N})\partial_{x}^{n}(\Phi\partial_{x}^{-N})^{-1}\]
L'operatore $\Phi\partial_{x}^{-N}$ è un operatore di ordine $0$(si ricordi che $\Phi$ ha come termine a grado maggiore $\partial_{x}^{N}$), dunque lo sarà anche il suo inverso $(\Phi\partial_{x}^{-N})^{-1}$
\[\Phi\partial_{x}^{-N}:= 1+a_{-1}\partial_{x}^{-1}+...+a_{-N}\partial_{x}^{-N}\Rightarrow(\Phi\partial_{x}^{-N})^{-1}=\sum_{i\geq0}b_{-i}\partial_{x}^{-i}\]
con $b_{0}=1$. Seguirà che $L=(1+a_{-1}\partial_{x}^{-1}+..+a_{-N}\partial_{x}^{-N})\partial_{x}^{n}(\sum_{i\geq0}b_{-i}\partial_{x}^{i})$ ovvero
\[L=\sum_{-\infty}^{n}c_{i}\partial_{x}^{i},\hspace{0.5cm}c_{n}=1\]
Ora, sfruttando l'esressione $L=L_{+}+L_{-}$, la (70) si può riscrivere come
\[L_{+}+L_{-}=\Phi\partial_{x}^{n}\Phi^{-1}\Rightarrow L_{+}\Phi=\Phi\partial_{x}^{n}-L_{-}\Phi\]
\[\Rightarrow L_{+}\Phi-\Phi\partial_{x}^{n}=-L_{-}\Phi\]
Il membro di destra è un operatore di ordine $<N$($L_{-}$ inizia con un termine di ordine $-1$), dunque lo sarà anche il membro di sinistra. Per il Lemma 7.3 ho poi che $(L_{+}\Phi-\Phi\partial_{x}^{n})y_{k}=0$. Se però un operatore di ordine $<N$ manda a zero $N$ funzioni linearmente indipendenti, significa che è identicamente nullo. Dunque $L_{-}\Phi=0$ da cui $L_{-}=0$ poichè $\Phi$ è invertibile.

(ii) L'equazione (70) implica che
\[L^{1/n}=\Phi\partial_{x}\Phi^{-1}\hspace{0.5cm}\overbrace{L^{m/n}=\Phi\partial_{x}^{m}\Phi^{-1}}^{(\star)}\hspace{0.5cm}\overbrace{L_{+}^{m/n}\Phi-\Phi\partial_{x}^{m}}^{:=Q}=-L_{-}^{m/n}\Phi\]
Infatti (la prima delle relazioni precedenti implica ovviamente la seconda)
\[(L^{1/n})^{n}=\underbrace{\Phi\partial_{x}\Phi^{-1}\Phi\partial_{x}\Phi^{-1}...\Phi\partial_{x}\Phi^{-1}}_{n-volte}=\Phi\partial_{x}^{n}\Phi^{-1}\]
mentre la terza si ottiene dalla seconda in modo analogo a prima. In particolare per quest'ultima si osservi che il secondo membro è un operatore di ordine $<N$, mentre l'operatore a primo membro è un operatore che dentiamo $Q$ (perciò $ord(Q)<N$).Applicando ora $\partial_{m}=\partial_{t}$ ad $\Phi y_{k}=0$ si ottiene, tenendo conto del Lemma 7.2, che
\[0=(\partial_{t}\Phi)y_{k}+\Phi\overbrace{\partial_{t}y_{k}}^{=\partial^{m}y_{k}}=(\partial_{t}\Phi)y_{k}\textcolor{red}{+L_{+}^{m/n}\Phi y_{k}-L_{+}^{m/n}\Phi y_{k}}+\Phi\partial^{m}y_{k}=(\partial_{t}\Phi)y_{k}+L_{+}^{m/n}\Phi y_{k}-Qy_{k}\]
\[=((\partial_{t}\Phi)-Q)y_{k}\]
L'operatore $(\partial_{t}\Phi)-Q$ è di ordine $<N$ ed annulla $N$ funzioni linearmente indipendenti $y_{k}$, perciò $Q=\partial_{t}\Phi$. Inoltre
\[\partial_{t}\Phi=L_{+}^{m/n}\Phi-\Phi\partial^{m}\stackrel{(\star)}{=}L_{+}^{m/n}\Phi-L^{m/n}\Phi=-L_{-}^{m/n}\Phi\]
Finalmente allora
\[\partial_{t}L=(\partial_{t}\Phi)\partial^{n}\Phi^{-1}+\Phi\partial^{n}(\partial_{t}\Phi^{-1})\stackrel{(1)}{=}(\partial_{t}\Phi)\partial^{n}\Phi^{-1}-\Phi\partial^{n}\Phi^{-1}(\partial_{t}\Phi)\Phi^{-1}\]
\[\stackrel{(2)}{=}(L_{+}^{m/n}\Phi-\Phi\partial^{m})\partial^{n}\Phi^{-1}-\Phi\partial^{n}\Phi^{-1}(L_{+}^{m/n}\Phi-\Phi\partial^{m})\Phi^{-1}=L_{+}^{m/n}\Phi\partial^{n}\Phi^{-1}-\Phi\partial^{n}\Phi^{-1}L_{+}^{m/n}\]
\[=L_{+}^{m/n}L-LL_{+}^{m/n}=[P_{m},L]\]
dove in $(2)$ si è usata la relazione trovata sopra  $\partial_{t}\Phi=Q$, ed in $(1)$ si è usato che $0=\partial_{t}(\Phi\Phi^{-1})\Rightarrow\Phi(\partial_{t}\Phi^{-1})=-(\partial_{t}\Phi)\Phi^{-1}\Rightarrow\partial_{t}\Phi^{-1}=-\Phi^{-1}(\partial_{t}\Phi)\Phi^{-1}$.
\end{proof}

\begin{flushleft}
$\underline{Esempio}$:Si consideri $n=2,m=3,N=1,L=\partial_{x}^{2}+u$. Siccome $\epsilon^{2}=1$ ho
\[\epsilon=\pm1,\hspace{0.5cm}y(x,t)=e^{\alpha x+1^{3}\alpha^{3}t}+ae^{-\alpha x+(-1)^{3}\alpha^{3}t}=e^{\alpha x+\alpha^{3}t}+ae^{-\alpha x-\alpha^{3}t}\]
In questo caso poi si ha $\Delta=y$ per cui
\[\Phi=\frac{1}{y}Det\begin{pmatrix}
y&1\\
y'&\partial_{x}
\end{pmatrix}=\partial_{x}-\frac{y'}{y}\]
Definita $\varphi:=\frac{y'}{y}$, e posto
\[\Phi=\partial_{x}^{-1}+Y_{-2}\partial_{x}^{-2}+Y_{-3}\partial_{x}^{-3}+...\]
si vede che 
\[1\equiv\Phi\Phi^{-1}\iff(\partial_{x}-\varphi)(\partial_{x}^{-1}+Y_{-2}\partial_{x}^{-2}+Y_{-3}\partial_{x}^{-3}+...)=1+Y_{-2}\partial_{x}^{-1}+Y_{-2}'\partial_{x}^{-2}+\]
\[Y_{-3}\partial_{x}^{-2}+Y_{-3}'\partial_{x}^{-3}-\varphi\partial_{x}^{-1}-\varphi Y_{-2}\partial_{x}^{-2}-\varphi Y_{-3}\partial_{x}^{-3}+...\equiv 1\]
\[\iff\begin{cases}
Y_{-2}-\varphi=0\rightarrow Y_{-2}=\varphi\\
Y_{-2}'+Y_{-3}-\varphi Y_{-2}=0\rightarrow Y_{-3}=\varphi^{2}-\varphi'
\end{cases}\]
\[\Rightarrow \Phi^{-1}=\partial_{x}^{-1}+\varphi\partial_{x}^{-2}+(\varphi^{2}-\varphi')\partial_{x}^{-3}+o(\partial_{x}^{-4})\]
In questo modo trovo
\[\Phi\partial_{x}^{2}\Phi^{-1}=\Phi(\partial_{x}+\varphi+\binom{2}{1}\varphi'\partial_{x}^{-1}+\binom{2}{2}\varphi''\partial_{x}^{-2}+(\varphi^{2}-\varphi')\partial_{x}^{-1}+\binom{2}{1}(2\varphi\varphi'-\varphi'')\partial_{x}^{-2}\]
\[+\binom{2}{2}(2\varphi'^{2}+2\varphi\varphi''-\varphi''')\partial_{x}^{-3}+...)\]
\[=(\partial_{x}-\varphi)(\partial_{x}+\varphi+(\varphi'+\varphi^{2})\partial_{x}^{-1}+(4\varphi\varphi'-\varphi'')\partial_{x}^{-2}+(2\varphi'^{2}+2\varphi\varphi''-\varphi''')\partial_{x}^{-3})\]
\[=\partial_{x}^{2}+\varphi\partial_{x}+\varphi'+(\varphi'+\varphi^{2})+(2\varphi\varphi'+\varphi'')\partial_{x}^{-1}+(4\varphi\varphi'-\varphi'')\partial_{x}^{-1}+(4\varphi'^{2}+4\varphi\varphi''-\varphi''')\partial_{x}^{-2}\]
\[+(2\varphi'^{2}+2\varphi\varphi''-\varphi''')\partial_{x}^{-2}+(4\varphi'\varphi''+2\varphi\varphi'''+2\varphi'\varphi''-\varphi^{iv})\partial_{x}^{-3}+...-(\varphi\partial+\varphi^{2}+(\varphi\varphi'+\varphi^{3})\partial_{x}^{-1}+...)=\partial_{x}^{2}+2\varphi'\]
dal momento che la Proposizione precedente mi garantisce che $L$ non possiede termini a potenze negative. Ci siamo ricondotti all'equazione
\[L=\partial_{x}^{2}+u=\partial_{x}^{2}+2\varphi'\]
\[\iff u(x,t)=2\varphi'=2(\frac{y''}{y}-\frac{y'^{2}}{y^{2}})\]
Risolvendo esplicitamente si ricava 
\[u(x,t)=2\bigg[\alpha^{2}-\frac{(\alpha e^{\alpha x+a\alpha^{3}t}-a\alpha e^{-\alpha x-\alpha^{3}t})^{2}}{(e^{\alpha x+\alpha^{3}t}+ae^{-\alpha x-\alpha^{3}t})^{2}}\bigg]\]
\[=\frac{2\alpha^{2}}{(e^{\alpha x+\alpha^{3}t}+ae^{-\alpha x-\alpha^{3}t})^{2}}\bigg[e^{2\alpha x+2\alpha^{3}t}+a^{2}e^{-2\alpha x-2\alpha^{3}t}+2a-e^{2\alpha x+2\alpha^{3}t}-a^{2}e^{-2\alpha x-2\alpha^{3}t}+2a\bigg]\]
\[=\frac{2\alpha^{2}4a}{(e^{\alpha x+\alpha^{3}t}+ae^{-\alpha x-\alpha^{3}t})^{2}}=\frac{2\alpha^{2}}{\frac{1}{4}(\frac{e^{\alpha x+\alpha^{3}t}+ae^{-\alpha x-\alpha^{3}t}}{\sqrt{a}})^{2}}=\frac{2\alpha^{2}}{\frac{1}{4}(e^{\alpha x+\alpha^{3}t-\log(a)/2}+e^{-(\alpha x+\alpha^{3}t-\log(a)/2)})^{2}}\]
ovverosia, in definitiva
\begin{equation}
u(x,t)=\frac{2\alpha^{2}}{\cosh^{2}(x\alpha+t\alpha^{3}-\frac{1}{2}\log(a))}
\end{equation}
Questo è il tipico profilo di un solitone che si propaga con velocità costante (dipendente da $\alpha$).

Si osservi come, al variare della velocità (quindi di $\alpha$), si ha che maggiore è quest'ultimo e maggiore è l'ampiezza (nonostante la larghezza diminuisca).
\end{flushleft}


    \begin{figure}
        \centerline{\includegraphics[scale=0.2]{SolitonPlot_1.jpg}}
        \caption{Plot della soluzione (71):si tratta di un solitone ad istante fissato, che per semplicità abbiamo considerato come unitario.}
        \label{fig:my_label}
    \end{figure}

    \begin{figure}
        \centerline{\includegraphics[scale=0.2]{Solitone_Din_Opt.png}}
        \caption{Plot della soluzione (71)in due dimensioni}
        \label{fig:my_label}
    \end{figure}

\section{Un altro esempio:il Toda lattice}
Riprendiamo ora l'esempio esposto nella sezione 4.1: avevamo visto come, per il sistema di Toda, le equazioni di Hamilton si possano esprimere come 
\[\begin{cases}\dot{q}_{n}=p_{n}\\
\dot{p}_{n}=e^{q_{n-1}-q_{n}}-e^{q_{n}-q_{n+1}}\end{cases}\]
Se introduciamo la nuova variabile
\[v_{n}:=e^{q_{n-1}-q_{n}}\]
allora il sistema precedente implica (ma non è equivalente) il sistema di equazioni
\begin{equation}
\begin{cases}
\dot{p}_{n}=v_{n}-v_{n+1}\\
\dot{v}_{n}=v_{n}(p_{n-1}-p_{n})
\end{cases}
\end{equation}
Abbiamo già visto la forma di Lax associata a tale sistema, nonchè il fatto che le traccie $Tr(L^{k})$ risultino degli integrali primi del sistema (una volta definito opportunamente $L$). Tuttavia a quel punto ci si può chiedere quale sia la traccia di $L^{k}$ nel caso in cui si abbia un numero infinito di particelle, caso in cui $L$ risulta di dimensione infinita. Come esposto nelle sezioni precedenti, la risposta sta nel considerare le "densità" anzichè le tracce globali.

Si considerino due funzioni $v,p:\mathbb{Z}\rightarrow\mathbb{R}$, l'insieme di tali funzioni costituisce un $\mathbb{R}-$algebra $\mathcal{A}$. In tale algebra definiamo l'operatore
\[\Delta:\mathcal{A}\rightarrow\mathcal{A}\]
\[f(n)\mapsto(\Delta f)(n):=f(n+1)\]
Se supponiamo che $p,v$ dipendano da un ulteriore parametro $t$, allora il sistema (72) può essere riscritto come 
\[\begin{cases}
p_{t}\equiv\dot{p}=v-\Delta v\\
v_{t}\equiv\dot{v}=v(\Delta^{-1}p-p)
\end{cases}\]
e ci siamo dunque ricondotti ad uguaglianze tra funzioni. Si consideri ora l'anello associativo $\mathcal{A}[\xi,\xi^{-1}]$ con coefficenti in $\mathcal{A}$, e si definisca la regola di commutazione
\[\xi^{s}f=(\Delta^{s}f)\xi^{s}\hspace{0.5cm}\forall s\in\mathbb{Z}\]
Definiamo anche gli operatori
\[L:=\xi+p+v\xi^{-1},\hspace{0.2cm}P:=v\xi^{-1}\in\mathcal{A}[\xi,\xi^{-1}]\]
In tal maniera
\[L_{t}=p_{t}+v_{t}\xi^{-1}\]
mentre
\[[P,L]=[v\xi^{-1},\xi+p+v\xi^{-1}]=v\xi^{-1}(\xi+p)-(\xi+p)v\xi^{-1}=v+v(\Delta^{-1} p)\xi^{-1}-(\Delta v)\xi\xi^{-1}-pv\xi^{-1}\]
\[=v-\Delta v+(\Delta^{-1} p-p)v\xi^{-1}\]
Uguagliandolo ad $L_{t}$, abbiamo dunque dimostrato che il sistema (72) è soddisfatto se e solo se
\begin{equation}
L_{t}=[P,L]
\end{equation}
Collegata alle equazioni di Toda, vi è un intera famiglia di equazioni
\[\frac{\partial L}{\partial t_{p}}=\frac{1}{(p+1)!}[(L^{p+1})_{+},L]\hspace{0.5cm}\frac{\partial }{\partial t_{q}}\frac{\partial L}{\partial t_{p}}=\frac{\partial }{\partial t_{p}}\frac{\partial L}{\partial t_{q}}\hspace{0.5cm}p,q\geq0\]
che assieme costituiscono la "gerarchia del Toda lattice". Chiaramente per $p=0$ si ottiene la (72). Per ulteriori informazioni rimandiamo al testo di Kupershmidt [15].




\appendix
\section{Appendice}
\subsection{Dimostrazione della formula (16)}
\begin{Prop}
Sia M una varietà differenziabile, $X\in\mathcal{T}(M)$ un campo vettoriale su M e $\sigma\in\mathcal{A}^{k}(M)$ un campo tensoriale. Allora per ogni k-upla di campi vettoriali $Y_{1},...,Y_{k}\in\mathcal{T}(M)$ si ha
\begin{equation}
\mathcal{L}_{X}(\sigma(Y_{1},...,Y_{k}))=(\mathcal{L}_{X}\sigma) (Y_{1},...,Y_{k})+\sigma(\mathcal{L}_{X}Y_{1},...,Y_{k})+...+\sigma(Y_{1},...,\mathcal{L}_{X}Y_{k})
\end{equation}
\end{Prop}
\begin{proof}
Si denoti con $\theta$ il flusso del campo X,dalla definizione si ha che per ogni punto $p\in M$
\[\mathcal{L}_{X}(\sigma(Y_{1},...,Y_{k}))_{p}=\frac{\partial}{\partial t}\at[\bigg]{t=0}(\theta_{t}^{*}\sigma(Y_{1},...,Y_{k}))_{p}=\]
\[=\lim_{t\rightarrow 0}\frac{1}{t}\{(\sigma)_{\theta_{t}(p)}(Y_{1}\at{\theta_{t}(p)},...,Y_{k}\at{\theta_{t}(p)})-\sigma_{p}(Y_{1}\at{p},...,Y_{k}\at{p})\}=\]
\[=\lim_{t\rightarrow 0}\{\frac{1}{t}[(\sigma)_{\theta_{t}(p)}(Y_{1}\at{\theta_{t}(p)},...,Y_{k}\at{\theta_{t}(p)})\textcolor{blue}{-\sigma_{p}((\theta_{-t*})Y_{1}\at{\theta_{t}(p)},...,(\theta_{-t*})Y_{k}\at{\theta_{t}(p)})}  ]+\]
\[\frac{1}{t}[\textcolor{blue}{+\sigma_{p}((\theta_{-t*})Y_{1}\at{\theta_{t}(p)},...,(\theta_{-t*})Y_{k}\at{\theta_{t}(p)})}\textcolor{orange}{-\sigma_{p}(Y_{1}\at{p},(\theta_{-t*})Y_{2}\at{\theta_{t}(p)},...,(\theta_{-t*})Y_{k}\at{\theta_{t}(p)})}]+\]
\[\frac{1}{t}[\textcolor{orange}{+\sigma_{p}(Y_{1}\at{p},(\theta_{-t*})Y_{2}\at{\theta_{t}(p)},...,(\theta_{-t*})Y_{k}\at{\theta_{t}(p)})}-\sigma_{p}(Y_{1}\at{p},Y_{2}\at{p},(\theta_{-t*})Y_{3}\at{\theta_{t}(p)},...,(\theta_{-t*})Y_{k}\at{\theta_{t}(p)})]+\]
\[+......+\]
\[+\frac{1}{t}[\sigma_{p}(Y_{1}\at{p},...,Y_{k-1}\at{p},(\theta_{-t*})Y_{k}\at{\theta_{t}(p)})-\sigma(Y_{1}\at{p},...,Y_{k}\at{p})]\}\]
dove il primo termine è pari a
\[\lim_{t\rightarrow 0}\frac{1}{t}[(\theta_{t}^{*}\sigma)_{p}((\theta_{-t*})Y_{1}\at{\theta_{t}(p)},...,(\theta_{-t*})Y_{k}\at{\theta_{t}(p)})-\sigma_{p}((\theta_{-t*})Y_{1}\at{\theta_{t}(p)},...,(\theta_{-t*})Y_{k}\at{\theta_{t}(p)})]=\]
\[=\lim_{t\rightarrow 0}[\frac{(\theta^{*}_{t}\sigma)_{p}-\sigma_{p}}{t}](((\theta_{-t*})Y_{1}\at{\theta_{t}(p)},...,(\theta_{-t*})Y_{k}\at{\theta_{t}(p)}))=(\mathcal{L}_{X}\sigma)_{p}(Y_{1}\at{p},...,Y_{k}\at{p})\]
mentre i restanti sono del tipo
\[\lim_{t\rightarrow 0}\frac{1}{t}[\sigma_{p}(Y_{1}\at{p},...,Y_{i-1}\at{p},(\theta_{-t*})Y_{i}\at{\theta_{t}(p)},...,(\theta_{-t*})Y_{k}\at{\theta_{t}(p)})-\sigma_{p}(Y_{1}\at{p},...,Y_{i}\at{p},(\theta_{-t*})Y_{i+1}\at{\theta_{t}(p)},...,(\theta_{-t*})Y_{k}\at{\theta_{t}(p)})]\]
\[=\lim_{t\rightarrow 0}\sigma_{p}(Y_{1}\at{p},...,\frac{(\theta_{-t*})Y_{i}\at{\theta{t}(p)}-Y_{i}\at{p}}{t},(\theta_{-t*})Y_{i+1}\at{\theta_{t}(p)},...,(\theta_{-t*})Y_{k}\at{\theta_{t}(p)})\]
\[=\sigma_{p}(Y_{1}\at{p},...,Y_{i-1}\at{p},(\mathcal{L}_{X}Y_{i})_{p},Y_{i+1}\at{p},...,Y_{k}\at{p})\]
Per arbitrarietà di p ottengo allora la tesi.
\end{proof}

\begin{Prop}
Se $X\in\mathcal{T}(M)$ è un campo vettoriale sulla varietà liscia M e $\sigma\in\mathcal{A}^{k}(M)$ è un campo tensoriale, la derivata di Lie $\mathcal{L}_{X}\sigma$ può essere espressa come in (16)
\[    (\mathcal{L}_{X}\sigma)(Y_{1},...,Y_{k})=X(\sigma(Y_{1},...,Y_{k}))-\sigma([X,Y_{1}],Y_{2},...,Y_{k})-...-\sigma(Y_{1},...,Y_{k-1},[X,Y_{k}])\]
\end{Prop}
\begin{proof}
Basta usare la (36) e risoverla per $\mathcal{L}_{X}\sigma$ rimpiazzando le varie $\mathcal{L}_{X}Y_{i}$ con $[X,Y_{i}]$.
\end{proof}

\subsection{Dimostrazione Prop.1.10}
Dimostriamo ora la relazione sfruttata nella Prop.1.10
\begin{Prop}
Sia M una varietà differenziabile, $X,Y\in\mathcal{T}(M)$ due campi vettoriali e $\omega\in\mathcal{A}^{k+1}(M)$ una 2-forma su M. Allora
\[\mathcal{L}_{X}(Y\lrcorner\omega)=(\mathcal{L}_{X}Y)\lrcorner\omega+Y\lrcorner(\mathcal{L}_{X}\omega)\]
\end{Prop}
\begin{proof}
Anzitutto osserviamo che per ogni punto $p\in M$, se $\theta$ è il flusso di X, si ha che (dalla definizione) $\forall Y_{1},...,Y_{k}\in T_{\theta_{t}(p)}M$
\[(Y\lrcorner\omega)_{\theta_{t}(p)}(Y_{1},...,Y_{k})=\omega_{\theta_{t}(p)}(Y_{\theta_{t}(p)},Y_{1},...,Y_{k})\]
\[\Rightarrow (Y\lrcorner\omega)_{\theta_{t}(p)}=Y_{\theta_{t}(p)}\lrcorner\omega_{\theta_{t}(p)}\]
per ogni istante $t\in\mathcal{D}_{p}$ nel dominio di $\theta^{p}$.Inoltre per ogni scelta di $X_{1},...,X_{k}\in T_{\theta_{t}(p)}M$
\[\theta_{t}^{*}(Y\lrcorner\omega)_{\theta_{t}(p)}(X_{1},...,X_{k})=(Y\lrcorner\omega)_{\theta_{t}(p)}(\theta_{t*}X_{1},...,\theta_{t*}X_{k})=\omega_{\theta_{t}(p)}(Y_{\theta_{p}(p)},\theta_{t*}X_{1},...,\theta_{t*}X_{k})\]
\[=\theta_{t}^{*}\omega_{\theta_{t}(p)}(\theta_{-t*}Y_{\theta_{t}(p)},X_{1},...,X_{k})=(\theta_{-t*}Y_{\theta_{t}(p)})\lrcorner(\theta_{t}^{*}\omega_{\theta_{t}(p)})(X_{1},...,X_{k})\]
\[\Rightarrow \theta_{t}^{*}(Y\lrcorner\omega)=(\theta_{-t*}Y)\lrcorner(\theta_{t}^{*}\omega)\]
Segue allora che dalla definizione di derivata di Lie
\[(\mathcal{L}_{X}(Y\lrcorner\omega))_{p}=\lim_{t\rightarrow 0}\frac{1}{t}[\theta_{t}^{*}(Y\lrcorner\omega)_{\theta_{t}(p)}-(Y\lrcorner\omega)_{p}]=\lim_{t\rightarrow0}\frac{1}{t}\{[(\theta_{-t*}Y_{\theta_{t}(p)})\lrcorner(\theta_{t}^{*}\omega_{\theta_{t}(p)})\textcolor{red}{-Y_{p}\lrcorner\theta_{t}^{*}\omega_{\theta_{t}(p)}}]\]
\[+\frac{1}{t}[\textcolor{red}{Y_{p}\lrcorner\theta_{t}^{*}\omega_{\theta_{t}(p)}}-(Y\lrcorner\omega)_{p}]\}=\lim_{t\rightarrow0}[(\frac{\theta_{-t*}Y_{\theta_{t}(p)}-Y_{p}}{t})\lrcorner\omega_{\theta_{t}(p)}+Y_{p}\lrcorner(\frac{\theta_{t}^{*}\omega_{\theta_{t}(p)}-\omega_{p}}{t})]\]
\[=(\mathcal{L}_{X}Y)_{p}\lrcorner\omega_{p}+Y_{p}\lrcorner(\mathcal{L}_{X}\omega)_{p}=(\mathcal{L}_{X}Y\lrcorner\omega)_{p}+(Y\lrcorner\mathcal{L}_{X}\omega)_{p}\]
\end{proof}



\subsection{Dimostrazione della regola di convoluzione di Vandermonde}
Vogliamo completare la dimostrazione della Proposizione 7.4.

\begin{Prop}
Per ogni intero $k\leq m\leq n$ vale
\[\binom{n}{m}\binom{m}{k}=\binom{n}{k}\binom{n-k}{m-k}\]
\end{Prop}
\begin{proof}
E' un semplice calcolo
\[\binom{n}{k}\binom{n-k}{m-k}=\frac{n!}{(n-k)!k!}\frac{(n-k)!}{(n-m)!(m-k)!}=\frac{n!}{(n-m)!\textcolor{brown}{m!}}\frac{\textcolor{brown}{m!}}{(m-k)!k!}\]
\[=\binom{n}{m}\binom{m}{k}\]
\end{proof}

\begin{Prop}
Siano $n,k,m\in\mathbb{N}$ degli interi non negativi. Allora vale la "regola di convoluzione di Vandermonde"
\[\sum_{r=0}^{k}\binom{n}{r}\binom{m}{k-r}=\binom{n+m}{k}\]
\end{Prop}
\begin{proof}
Si consideri un generico polinomio $1+x$, per il Teorema binomiale 
\[(1+x)^{m+n}=\sum_{r=0}^{m+n}\binom{m+n}{r}x^{r}\hspace{0.5cm}\forall m,n\in\mathbb{N}\]
D'altra parte però posso scrivere anche
\[(1+x)^{m+n}=(1+x)^{m}(1+x)^{n}=(\sum_{i=0}^{m}\binom{m}{i}x^{i})(\sum_{j=0}^{n}\binom{n}{j}x^{j})\]
\[=\sum_{r=0}^{m+n}(\sum_{k=0}^{r}\binom{m}{k}\binom{n}{r-k})x^{r}\]
\end{proof}



\section{Bibliografia}
 
\begin{thebibliography}{30}

\bibitem[1]{} Dickey L.A.,Soliton equations and Hamiltonian systems, Advanced Series in Mathematical Physics, Vol.26,2nd ed.,World Scientific Publishing Co.,Inc.,River Edge,NJ,2003.

\bibitem[2]{} Buryak A.,Double ramification cycle and integrable hierarchies,Comm. Math. Phys. 336(2015),1085-1107, arXiv:1403.1719.

\bibitem[3]{} Buryak A.,Rossi P., Double ramification cycle and quantum integrable systems,Lett.Math.Phys.106(2016),289-317,arXiv:1503.03687.

\bibitem[4]{} Buryak A.,Rossi P.,Recursion relations for double ramification hierarchies, Comm.Math.Phys.342(2016),533-568,arXiv:1411.6797.

\bibitem[5]{} Rossi P.,Integrable systems and moduli spaces of curves, Mémoire d'Habilitation à Diriger des Recherches,Universìte de Bourgogne(Dijon),2016, available at:https://hal.archives-ouvertes,fr/tel-01444750.

\bibitem[6]{} Buryak A.,Dubrovin B., Guéré J., Rossi P., Tau-structure for the double ramification hierarchies, arXiv:1602.05423.

\bibitem[7]{} Buryak A.,Dubrovin B.,Guéré J., Rossi P.,Integrable systems of double ramification type, arXiv:1609.04059.

\bibitem[8]{} Lee J.M., Introduction to smooth manifolds,2000,University of Washington,WA 98195-4350.

\bibitem[9]{} Ana Cannas da Silva, Lectures on Symplectic Geometry,2006,Springer-Verlag,available at:www.springerlink.com.

\bibitem[10]{} J.P.Fortney, A visual introduction to differential forms and calculus on manifolds,Springer nature Switzerland AG 2018.

\bibitem[11]{} Schalch N.,Korteweg-de Vries equation,ETH ZÜRICH,March 26,2018.

\bibitem[12]{} J.E.Marsden, Introduction to Mechanics and Symmetry, 2nd edition, 15 July 1998.

\bibitem[13]{} Y.Jiang,H.Najafi and J.pozny, Poisson Geometry lectures, May 17 2004.

\bibitem[14]{}The Caratheristic Polynomial, Physics 116A, Winter 2011.

\bibitem[15]{}Discrete Lax equations and differential-difference calculus, B.A. Kupershmidt, Société mathématique de France, 1985, tous droits reserves. 

\bibitem[16]{}GetzlerE.,A Darboux theorem for Hamiltonian operators in the formal calculus of variations,DukeMath.J.
111 (2002), 535-560.

\bibitem[17]{}A. Buryak, Integrable systems as systems of PDEs with an infinite dimensional algebra of simmetries.
 
\end{thebibliography}
\end{document}

